%=====================================================================
% INSTALACOES ELETRICAS PREDIAIS - PROJETO E DIMENSIONAMENTO
% Parte 1 de 2 do material predial
% Atualizacao: agosto de 2026
%=====================================================================
\documentclass[11pt,aspectratio=169]{beamer}

\usepackage[utf8]{inputenc}
\usepackage[T1]{fontenc}
\usepackage[brazil]{babel}
\usepackage{lmodern}
\usepackage{amsmath,amssymb,bm}
\usepackage{siunitx}
\usepackage{booktabs,array,tabularx,multirow}
\usepackage[most]{tcolorbox}
\usepackage{tikz}
\usepackage{pgfplots}
\usepackage[european,american currents]{circuitikz}
\usepackage{ragged2e}
\usepackage{microtype}
\usepackage{xcolor}
\usetikzlibrary{arrows.meta,positioning,calc,fit,shapes.geometric,patterns,decorations.pathmorphing,decorations.markings,matrix}
\pgfplotsset{compat=1.18}
\sisetup{locale=DE,per-mode=symbol,detect-all}

\newcommand{\sen}{\operatorname{sen}}
\newcommand{\tg}{\operatorname{tg}}
\newcommand{\fase}[1]{\textcolor{corPrim}{\textbf{#1}}}
\newcommand{\alerta}[1]{\textcolor{corAt}{\textbf{#1}}}
\newcommand{\Disciplina}{}
\newcommand{\TituloCurto}{Projeto}
%=====================================================================
% ESTILO DA CASA -- TEMPLATE BEAMER (padrao da colecao)
% Arquivo central de identidade visual para todos os materiais.
% Azul-petroleo + ambar | tcolorbox | TikZ/pgfplots/circuitikz
%=====================================================================

% Dependencias minimas do proprio estilo.
\RequirePackage{xcolor}
\RequirePackage{tikz}
\usetikzlibrary{calc}
\RequirePackage{tcolorbox}
\tcbuselibrary{skins,breakable}
\RequirePackage{listings}
\RequirePackage{xparse}
\RequirePackage{etoolbox}

\makeatletter

%====================== PALETA OFICIAL ================================
\definecolor{corPrim}{RGB}{13,48,82}      % azul-petroleo
\definecolor{corAccent}{RGB}{222,138,28}  % ambar
\definecolor{corDef}{RGB}{31,111,168}     % azul medio
\definecolor{corNorma}{RGB}{0,121,107}    % verde-petroleo
\definecolor{corEx}{RGB}{46,125,50}       % verde
\definecolor{corFix}{RGB}{204,102,0}      % laranja
\definecolor{corAt}{RGB}{176,58,46}       % vermelho
\definecolor{codBg}{RGB}{252,250,245}
\definecolor{codKw}{RGB}{13,48,82}
\definecolor{codStr}{RGB}{46,125,50}
\definecolor{codCom}{RGB}{120,130,138}

% Compatibilidade com nomes usados historicamente nos materiais.
\colorlet{Petroleo}{corPrim}
\colorlet{PetroleoEscuro}{corPrim!82!black}
\colorlet{PetroleoMedio}{corDef!82!corPrim}
\colorlet{PetroleoClaro}{corDef!7}
\colorlet{Ambar}{corAccent}
\colorlet{AmbarEscuro}{corAccent!72!black}
\colorlet{AmbarClaro}{corAccent!16}
\colorlet{AzulPetroleo}{corPrim}
\colorlet{AzulPetroleoEscuro}{corPrim!82!black}
\colorlet{AzulPetroleoClaro}{corDef!7}
\colorlet{AzulMedio}{corDef}
\colorlet{AzulClaro}{corDef!7}
\colorlet{CinzaEscuro}{black!82}
\colorlet{CinzaClaro}{black!5}
\colorlet{CinzaSuave}{black!3}
\colorlet{CinzaMedio}{black!55}
\colorlet{CinzaTexto}{black!82}
\colorlet{CinzaBox}{black!3}
\colorlet{Verde}{corEx}
\colorlet{Vermelho}{corAt}
\colorlet{Turquesa}{corNorma!70!corDef}
\colorlet{Roxo}{corDef!58!corAt}
\colorlet{ifmagreen}{corPrim}
\colorlet{ifmagreen2}{corDef}
\colorlet{ifmalight}{corDef!7}
\colorlet{accent}{corAccent}
\colorlet{warn}{corAt}
\colorlet{ok}{corEx}
\colorlet{bluex}{corDef}
\colorlet{grayx}{black!58}

% Variantes em minusculas presentes em materiais antigos.
\colorlet{azulPetroleo}{corPrim}
\colorlet{azulPetroleoEscuro}{corPrim!82!black}
\colorlet{azulPetroleoClaro}{corDef!7}
\colorlet{azulMedio}{corDef}
\colorlet{azulClaro}{corDef!7}
\colorlet{azulEscuro}{corPrim!82!black}
\colorlet{azul}{corPrim}
\colorlet{azul2}{corDef}
\colorlet{ambar}{corAccent}
\colorlet{ambarClaro}{corAccent!16}
\colorlet{cinza}{black!5}
\colorlet{cinzaClaro}{black!5}
\colorlet{cinzaMedio}{black!55}
\colorlet{cinzaEscuro}{black!82}
\colorlet{cinzaTexto}{black!82}
\colorlet{verde}{corEx}
\colorlet{verdeclaro}{corEx!10}
\colorlet{vermelho}{corAt}
\colorlet{vermelhoclaro}{corAt!9}
\colorlet{laranja}{corFix}
\colorlet{roxo}{corDef!58!corAt}

%====================== TEMA BEAMER ==================================
\usetheme{default}
\usefonttheme{professionalfonts}
\setbeamercolor{normal text}{fg=black!85,bg=white}
\setbeamercolor{structure}{fg=corPrim}
\setbeamercolor{frametitle}{bg=corPrim,fg=white}
\setbeamercolor{footsec}{bg=corPrim,fg=white!92!corPrim}
\setbeamercolor{alerted text}{fg=corAccent!85!black}
\setbeamerfont{frametitle}{series=\bfseries,size=\large}
\setbeamerfont{footline}{size=\tiny}
\setbeamertemplate{navigation symbols}{}
\setbeamertemplate{caption}{\raggedright\insertcaption\par}
\setbeamersize{text margin left=0.55cm,text margin right=0.55cm}

\setbeamertemplate{itemize item}{\raise0.35ex\hbox{\color{corAccent}\rule{0.62ex}{0.62ex}}}
\setbeamertemplate{itemize subitem}{\color{corDef}\raise0.5pt\hbox{--}}
\setbeamertemplate{itemize subsubitem}{\color{corPrim}\tiny$\circ$}
\setbeamertemplate{enumerate item}{\color{corPrim}\textbf{\insertenumlabel.}}
\setbeamertemplate{section in toc}{%
  \leavevmode\textcolor{corAccent}{\textbf{\inserttocsectionnumber.}}~\inserttocsection\par\vskip1.5pt}
\setbeamertemplate{subsection in toc}{%
  \leavevmode\leftskip=0.95cm\textcolor{corDef}{\textbullet}~\inserttocsubsection\par}

%====================== TITULO DO SLIDE ==============================
\setbeamertemplate{frametitle}{%
  \nointerlineskip%
  \begin{beamercolorbox}[wd=\paperwidth,ht=2.7ex,dp=1.3ex,leftskip=0.55cm,rightskip=0.55cm]{frametitle}%
    \usebeamerfont{frametitle}\insertframetitle%
  \end{beamercolorbox}%
  {\color{corAccent}\rule{\paperwidth}{1.3pt}}\par\vskip2pt%
}

%====================== RODAPE DUPLO ================================
\newcommand{\CasaRodapeEsq}{%
  \ifcsname Disciplina\endcsname\Disciplina\else\insertshorttitle\fi}
\newcommand{\CasaRodapeDir}{%
  \ifcsname TituloCurto\endcsname\TituloCurto\else\insertshorttitle\fi}
\newcommand{\CasaMetaCapa}{%
  \ifcsname Disciplina\endcsname
    \if\relax\detokenize\expandafter{\Disciplina}\relax
    \else
      \Disciplina\ \textbullet\ \insertdate
    \fi
  \else
    \insertdate
  \fi}
\setbeamertemplate{footline}{%
  \leavevmode%
  \hbox{%
    \begin{beamercolorbox}[wd=.5\paperwidth,ht=2.4ex,dp=1ex,leftskip=0.5cm]{footsec}%
      \usebeamerfont{footline}\CasaRodapeEsq%
    \end{beamercolorbox}%
    \begin{beamercolorbox}[wd=.5\paperwidth,ht=2.4ex,dp=1ex,rightskip=0.5cm]{footsec}%
      \hfill\usebeamerfont{footline}\CasaRodapeDir\quad\textbar\quad\insertframenumber%
    \end{beamercolorbox}%
  }\vskip0pt%
}

%====================== DIVISORIAS DE SECAO =========================
\newcommand{\CasaSecNum}{\ifnum\value{section}<10 0\fi\arabic{section}}
\newcommand{\CasaDivisoriaManual}[2]{%
  \begin{frame}[plain,noframenumbering]
    \makebox[0pt][l]{%
      \raisebox{0pt}[0pt][0pt]{%
        \begin{tikzpicture}[remember picture,overlay]
          \fill[corPrim] (current page.south west) rectangle (current page.north east);
          \if\relax\detokenize{#1}\relax\else
            \node[anchor=west,text=corAccent,font=\bfseries\fontsize{40}{40}\selectfont]
              at ([xshift=1.6cm,yshift=1.1cm]current page.west) {#1};
          \fi
          \fill[corAccent] ([xshift=1.7cm,yshift=0.35cm]current page.west)
            rectangle ++(2.4,0.09);
          \node[anchor=west,text=white,font=\bfseries\LARGE,text width=0.75\paperwidth]
            at ([xshift=1.6cm,yshift=-0.5cm]current page.west) {#2};
        \end{tikzpicture}%
      }%
    }%
  \end{frame}%
}
\AtBeginSection[]{\CasaDivisoriaManual{\CasaSecNum}{\insertsection}}

% Alguns materiais antigos usam \secslide{titulo}; outros, \secslide{numero}{titulo}.
\@ifundefined{secslide}
  {\NewDocumentCommand{\secslide}{m g}{\IfNoValueTF{#2}{\CasaDivisoriaManual{}{#1}}{\CasaDivisoriaManual{#1}{#2}}}}
  {\RenewDocumentCommand{\secslide}{m g}{\IfNoValueTF{#2}{\CasaDivisoriaManual{}{#1}}{\CasaDivisoriaManual{#1}{#2}}}}
\@ifundefined{Divisoria}
  {\NewDocumentCommand{\Divisoria}{m g}{\IfNoValueTF{#2}{\CasaDivisoriaManual{}{#1}}{\CasaDivisoriaManual{#1}{#2}}}}
  {\RenewDocumentCommand{\Divisoria}{m g}{\IfNoValueTF{#2}{\CasaDivisoriaManual{}{#1}}{\CasaDivisoriaManual{#1}{#2}}}}

%====================== CAPA ========================================
\defbeamertemplate*{title page}{estilo da casa}{%
  \vspace{0.8cm}%
  {\color{white}\bfseries\huge\inserttitle\par}%
  \vspace{2mm}%
  {\color{corAccent}\rule{3.2cm}{2.2pt}\par}%
  \vspace{4mm}%
  {\color{white!85}\large\insertsubtitle\par}%
  \vspace{8mm}%
  {\color{white}\normalsize\insertauthor\par}%
  \vspace{1mm}%
  {\color{white!70}\small\CasaMetaCapa\par}%
}
\setbeamertemplate{title page}[estilo da casa]

%====================== CAIXAS =======================================
% Materiais antigos usam convencoes diferentes para o argumento opcional
% das caixas. Para preservar a semantica, caixas existentes nao sao
% redefinidas; apenas sua paleta historica e recolorida acima.
\tcbset{
  caixabase/.style={
    enhanced,sharp corners=downhill,boxrule=0.6pt,left=2mm,right=2mm,
    top=1.2mm,bottom=1.2mm,fonttitle=\bfseries\small,
    attach boxed title to top left={xshift=3mm,yshift=-2.4mm},
    boxed title style={sharp corners,boxrule=0pt},
    before skip=3pt,after skip=3pt
  }
}
\@ifundefined{caixaDef}{\newtcolorbox{caixaDef}[1][]{caixabase,colback=corDef!5,colframe=corDef,coltitle=white,colbacktitle=corDef,title=Definição,#1}}{}
\@ifundefined{caixaNorma}{\newtcolorbox{caixaNorma}[1][]{caixabase,colback=corNorma!5,colframe=corNorma,coltitle=white,colbacktitle=corNorma,title=Norma / Método,#1}}{}
\@ifundefined{caixaEx}{\newtcolorbox{caixaEx}[1][]{caixabase,colback=corEx!5,colframe=corEx,coltitle=white,colbacktitle=corEx,title=Exemplo resolvido,#1}}{}
\@ifundefined{caixaFix}{\newtcolorbox{caixaFix}[1][]{caixabase,colback=corFix!5,colframe=corFix,coltitle=white,colbacktitle=corFix,title=Fixação,#1}}{}
\@ifundefined{caixaAt}{\newtcolorbox{caixaAt}[1][]{caixabase,colback=corAt!5,colframe=corAt,coltitle=white,colbacktitle=corAt,title=Atenção,#1}}{}
\@ifundefined{caixaForm}{\newtcolorbox{caixaForm}[1][]{caixabase,colback=corPrim!4,colframe=corPrim,coltitle=white,colbacktitle=corPrim,title=Formulário,#1}}{}
\@ifundefined{caixaObs}{\newtcolorbox{caixaObs}[1][]{caixabase,colback=black!3,colframe=corPrim,coltitle=white,colbacktitle=corPrim,title=Observação,#1}}{}
\@ifundefined{caixaAlerta}{\newtcolorbox{caixaAlerta}[1][]{caixabase,colback=corAt!5,colframe=corAt,coltitle=white,colbacktitle=corAt,title=Atenção,#1}}{}
\@ifundefined{caixaProjeto}{\newtcolorbox{caixaProjeto}[1][]{caixabase,colback=corNorma!5,colframe=corNorma,coltitle=white,colbacktitle=corNorma,title=Projeto prático,#1}}{}
\@ifundefined{caixaProj}{\newtcolorbox{caixaProj}[1][]{caixabase,colback=corNorma!5,colframe=corNorma,coltitle=white,colbacktitle=corNorma,title=Projeto prático,#1}}{}

%====================== CODIGO =======================================
\lstdefinestyle{codCasa}{
  basicstyle=\ttfamily\scriptsize,
  keywordstyle=\color{codKw}\bfseries,
  stringstyle=\color{codStr},commentstyle=\color{codCom}\itshape,
  numbers=left,numberstyle=\tiny\color{codCom},numbersep=5pt,
  frame=single,rulecolor=\color{corPrim!25},backgroundcolor=\color{codBg},
  breaklines=true,showstringspaces=false,tabsize=2,columns=fullflexible
}
\lstdefinestyle{codC}{style=codCasa,language=C}
\lstdefinestyle{codPy}{style=codCasa,language=Python}
\lstdefinestyle{py}{style=codCasa,language=Python}

%====================== EXTENSOES MODULARES ==========================
% Se existir <jobname>_antes_laboratorios.tex, o arquivo e carregado uma
% unica vez imediatamente antes da secao de laboratorios. Isso permite
% ampliar um material sem duplicar o arquivo principal.
\newif\ifCasaExtensaoAntesLabsInserida
\CasaExtensaoAntesLabsInseridafalse
\AtBeginDocument{%
  \IfFileExists{\jobname_antes_laboratorios.tex}{%
    \let\CasaSectionOriginal\section
    \renewcommand{\section}[1]{%
      \ifCasaExtensaoAntesLabsInserida\else
        \ifstrequal{##1}{Laboratórios, projetos e comissionamento}{%
          \global\CasaExtensaoAntesLabsInseridatrue
          \input{\jobname_antes_laboratorios.tex}%
        }{}%
      \fi
      \CasaSectionOriginal{##1}%
    }%
  }{}%
}

\makeatother

%=====================================================================
% SIMBOLOS ELETRICOS TIKZ — DEFINICOES SEM SLIDES
% Carregado no preambulo dos materiais que precisam desenhar plantas.
%=====================================================================

\providecommand{\SimLuz}[2]{\begin{scope}[shift={(#1,#2)}]
\draw[corAccent,thick] (0,0) circle (0.14);
\draw[corAccent,thick] (-0.10,-0.10)--(0.10,0.10);
\draw[corAccent,thick] (-0.10,0.10)--(0.10,-0.10);
\end{scope}}
\providecommand{\SimTUG}[2]{\begin{scope}[shift={(#1,#2)}]
\draw[corPrim,thick] (0,0) circle (0.13);
\draw[corPrim,thick] (-0.16,0)--(0.16,0);
\end{scope}}
\providecommand{\SimTUGM}[2]{\begin{scope}[shift={(#1,#2)}]
\draw[corPrim,thick] (0,0) circle (0.13);
\draw[corPrim,thick] (-0.16,0)--(0.16,0);
\draw[corPrim,thick] (-0.16,0.065)--(0.16,0.065);
\end{scope}}
\providecommand{\SimTUE}[3]{\begin{scope}[shift={(#1,#2)}]
\draw[corFix,thick] (0,0) circle (0.16);
\node[font=\tiny\bfseries,corFix] at (0,0) {#3};
\end{scope}}
\providecommand{\SimInterruptor}[3]{\begin{scope}[shift={(#1,#2)}]
\draw[corNorma,thick] (0,0) circle (0.11);
\node[font=\tiny\bfseries,corNorma,anchor=west] at (0.14,0) {#3};
\end{scope}}
\providecommand{\SimQD}[3]{\begin{scope}[shift={(#1,#2)}]
\draw[corPrim,fill=corDef!10,thick,rounded corners=1pt] (-0.32,-0.20) rectangle (0.32,0.20);
\node[font=\tiny\bfseries] at (0,0) {#3};
\end{scope}}
\providecommand{\SimDados}[2]{\begin{scope}[shift={(#1,#2)}]
\draw[corDef,thick] (0,0.15)--(-0.14,-0.11)--(0.14,-0.11)--cycle;
\node[font=\tiny\bfseries,corDef] at (0,-0.01) {D};
\end{scope}}
\providecommand{\SimMotor}[3]{\begin{scope}[shift={(#1,#2)}]
\draw[corAt,thick] (0,0) circle (0.20);
\node[font=\tiny\bfseries,corAt] at (0,0) {#3};
\end{scope}}
\providecommand{\SimVFD}[3]{\begin{scope}[shift={(#1,#2)}]
\draw[corDef,fill=corDef!8,thick,rounded corners=1pt] (-0.28,-0.18) rectangle (0.28,0.18);
\node[font=\tiny\bfseries,corDef] at (0,0) {#3};
\end{scope}}
\providecommand{\SimTomadaTri}[2]{\begin{scope}[shift={(#1,#2)}]
\draw[corFix,thick] (0,0) circle (0.16);
\node[font=\tiny\bfseries,corFix] at (0,0) {$3\sim$};
\end{scope}}
\providecommand{\SimEmerg}[2]{\begin{scope}[shift={(#1,#2)}]
\draw[corAt,fill=corAt!12,thick] (0,0) circle (0.16);
\node[font=\tiny\bfseries,corAt] at (0,0) {!};
\end{scope}}

\providecommand{\RotaL}[4]{\draw[corAccent!80!black,dashed,thick] (#1,#2)--(#3,#4);}
\providecommand{\RotaT}[4]{\draw[corPrim!75,dashed,thick] (#1,#2)--(#3,#4);}
\providecommand{\RotaF}[4]{\draw[corFix!80,dashed,thick] (#1,#2)--(#3,#4);}

\providecommand{\CotaH}[4]{%
\draw[thin] (#1,#3)--(#1,#3+0.28);
\draw[thin] (#2,#3)--(#2,#3+0.28);
\draw[{Latex}-{Latex},thin] (#1,#3+0.20)--(#2,#3+0.20)
node[midway,fill=white,inner sep=1pt,font=\tiny]{#4};}
\providecommand{\CotaV}[4]{%
\draw[thin] (#3,#1)--(#3+0.28,#1);
\draw[thin] (#3,#2)--(#3+0.28,#2);
\draw[{Latex}-{Latex},thin] (#3+0.20,#1)--(#3+0.20,#2)
node[midway,fill=white,inner sep=1pt,font=\tiny,rotate=90]{#4};}


\title{Instalações Elétricas Prediais}
\subtitle{Parte 1 — Leitura, projeto, cálculo e dimensionamento}
\author{Francisco}
\institute{Do levantamento arquitetônico ao projeto executivo}
\date{}

\begin{document}

{\setbeamercolor{background canvas}{bg=corPrim}\setbeamertemplate{footline}{}
\begin{frame}[plain,noframenumbering]\titlepage\end{frame}}

\begin{frame}{O que este Beamer entrega}
\small
\begin{columns}[T,onlytextwidth]
\column{0.49\textwidth}
\begin{caixaProjeto}[title=Projeto e cálculo]
\begin{itemize}
\item leitura da arquitetura e levantamento de cargas;
\item iluminação, TUG, TUE e cargas modernas;
\item demanda, circuitos e balanceamento de fases;
\item condutores, eletrodutos e queda de tensão;
\item curto-circuito, proteção, DR, DPS e aterramento.
\end{itemize}
\end{caixaProjeto}
\column{0.49\textwidth}
\begin{caixaNorma}[title=Da entrada à entrega]
\begin{itemize}
\item padrão de entrada Equatorial Maranhão;
\item documentação, unifilar, quadro de cargas e memorial;
\item projetos residenciais e loja de shopping completos;
\item exemplos resolvidos e exercícios graduados;
\item critérios declarados, rastreáveis e verificáveis.
\end{itemize}
\end{caixaNorma}
\end{columns}
\end{frame}

\section{Normas, escopo e ciclo de projeto}

\begin{frame}{Instalação elétrica: visão de sistema}
\small
\begin{columns}[T,onlytextwidth]
\column{0.52\textwidth}
\begin{caixaDef}
Uma instalação elétrica é um \textbf{sistema coordenado}: fonte, entrada, medição, distribuição, proteção, condutores, aterramento, cargas, comando, documentação e procedimentos de operação/manutenção.
\end{caixaDef}
\begin{itemize}
\item Predial: residência, comércio, escola, hospital, laboratório e serviços.
\item Industrial: máquinas, motores, fornos, automação, instrumentação, CCM e redes internas.
\item O objetivo não é apenas ``funcionar'': é garantir segurança, desempenho, continuidade e manutenção.
\end{itemize}
\column{0.46\textwidth}
\centering
\begin{tikzpicture}[scale=0.74,every node/.style={font=\scriptsize}]
\tikzset{b/.style={draw=corPrim,fill=corDef!7,rounded corners,minimum width=3.0cm,minimum height=0.65cm,align=center}}
\node[b] (r) at (0,3.2) {Rede / fonte local};
\node[b] (e) at (0,2.2) {Entrada e medição};
\node[b] (q) at (0,1.2) {QGBT / QD};
\node[b] (c1) at (-1.8,0.1) {Circuitos\\prediais};
\node[b] (c2) at (1.8,0.1) {Circuitos\\industriais};
\node[b] (s) at (0,-1.1) {PE + equipotencial\\+ DPS + SPDA};
\draw[-{Stealth},thick,corPrim] (r)--(e)--(q);
\draw[-{Stealth},thick,corPrim] (q)--(c1);
\draw[-{Stealth},thick,corPrim] (q)--(c2);
\draw[-{Stealth},thick,corNorma] (s)--(q);
\draw[-{Stealth},thick,corNorma] (s)--(c1);
\draw[-{Stealth},thick,corNorma] (s)--(c2);
\end{tikzpicture}
\end{columns}
\end{frame}

\begin{frame}[shrink=7]{Mapa normativo — situação em agosto de 2026}
\scriptsize
\begin{tabularx}{\textwidth}{>{\bfseries}p{2.8cm}p{3.1cm}X}
\toprule
Documento & Situação usada no material & Aplicação \\
\midrule
ABNT NBR 5410:2004 & Referência vigente adotada para BT; a própria NT.00001.EQTL Rev.09 a cita expressamente. & Instalações elétricas de baixa tensão: choque, condutores, circuitos, DR, aterramento, verificação etc. \\
ABNT NBR 14039:2021 & Referência para instalações de média tensão. & Instalações de 1,0 kV a 36,2 kV. \\
ABNT NBR 5419-1:2026 & Nova edição publicada em 2026. & Princípios e determinação da necessidade de proteção contra descargas atmosféricas; conferir as partes aplicáveis da série no projeto. \\
ABNT NBR 17019:2022 & Referência para alimentação de veículos elétricos. & Instalações fixas para recarga de VE; também referenciada pela NT.00042.EQTL Rev.04. \\
NR-10 + Portaria MTE 737/2026 & A nova redação foi publicada, mas a Portaria entra em vigor um ano após sua publicação. & Segurança em instalações e serviços com eletricidade; distinguir requisito vigente de requisito já publicado com vigência futura. \\
Série ABNT NBR IEC 61439 & Usar edição aplicável ao conjunto. & Quadros e conjuntos de manobra e comando de baixa tensão. \\
\bottomrule
\end{tabularx}
\end{frame}

\begin{frame}[shrink=8]{Equatorial Energia — referências atuais para o Maranhão}
\scriptsize
\begin{tabularx}{\textwidth}{>{\bfseries}p{2.65cm}p{2.15cm}X}
\toprule
Documento & Revisão/data & Quando usar \\
\midrule
NT.00001.EQTL & Rev.09 — 29/05/2025 & Fornecimento de energia em baixa tensão para unidade consumidora individual. É a referência principal do projeto residencial deste material. \\
NT.00004.EQTL & Rev.08 — 31/12/2025 & Empreendimentos/edificações de múltiplas unidades consumidoras — condomínios, edifícios e centros de medição. \\
NT.00020.EQTL & Rev.06 — 31/12/2025 & Conexão de micro e minigeração distribuída; formulários do Grupo B atualizados em 2026. \\
NT.00030.EQTL & Rev.03 — 31/12/2025 & Padrões construtivos das caixas de medição e proteção. \\
NT.00042.EQTL & Rev.04 — 31/12/2025 & Critérios de conexão de estações de recarga de veículos elétricos. \\
NT.00045.EQTL & Rev.01 — 29/03/2023 & Postes para padrão de entrada. \\
\bottomrule
\end{tabularx}
\begin{caixaAt}
As normas da concessionária tratam da interface com a rede e dos padrões de conexão. Elas \textbf{não substituem} o dimensionamento da instalação interna conforme ABNT e demais requisitos legais.
\end{caixaAt}
\end{frame}

\begin{frame}{Equatorial Maranhão: limites que afetam o projeto residencial}
\small
\begin{columns}[T,onlytextwidth]
\column{0.50\textwidth}
\begin{caixaNorma}[title=NT.00001.EQTL Rev.09]
A norma atende unidades individuais em BT com carga instalada de até 75 kW. Para o Maranhão:
\begin{itemize}
\item ligação monofásica urbana: 220 V, dentro dos limites previstos;
\item ligação trifásica urbana: \textbf{380/220 V, 3F+N};
\item acima dos limites e nos casos especiais, é necessária análise específica.
\end{itemize}
\end{caixaNorma}
\column{0.48\textwidth}
\begin{caixaProjeto}[title=Tabela 1 — 220/380 V]
Exemplos de enquadramento por carga instalada:
\begin{itemize}
\item até 5 kW: mono 25 A;
\item 5,1--8 kW: mono 40 A;
\item 8,1--12 kW: mono 63 A;
\item 12,1--24 kW: tri 40 A;
\item 24,1--38 kW: tri 63 A.
\end{itemize}
\end{caixaProjeto}
\end{columns}
\begin{caixaObs}
O projeto residencial completo deste material fecha em aproximadamente 19,88 kW instalados e, por isso, usa o enquadramento trifásico de 40 A como \textbf{consequência} da carga calculada.
\end{caixaObs}
\end{frame}

\begin{frame}{Projeto interno e projeto para a concessionária não são a mesma coisa}
\small
\begin{columns}[T,onlytextwidth]
\column{0.49\textwidth}
\begin{caixaNorma}[title=Unidade individual em BT]
A NT.00001.EQTL informa que, em regra, não é necessário apresentar à concessionária um projeto elétrico separado para novas instalações/reformas de unidade individual atendida em tensão secundária.
\end{caixaNorma}
\column{0.49\textwidth}
\begin{caixaAt}[title=Mas continuam obrigatórios]
\begin{itemize}
\item projeto interno tecnicamente correto;
\item padrão de entrada conforme concessionária;
\item ABNT, segurança e Corpo de Bombeiros quando aplicável;
\item normas específicas para GD, VE e EMUC;
\item responsabilidade técnica quando exigida.
\end{itemize}
\end{caixaAt}
\end{columns}
\end{frame}

\begin{frame}{NR-10: nova redação publicada, vigência futura}
\small
\begin{columns}[T,onlytextwidth]
\column{0.52\textwidth}
\begin{caixaNorma}[title=Portaria MTE nº 737/2026]
A Portaria é de 29/05/2026, foi publicada no DOU de 01/06/2026 e estabelece que sua nova redação da NR-10 \textbf{entra em vigor um ano após a publicação}.
\end{caixaNorma}
\begin{itemize}
\item integra explicitamente a gestão com o GRO/NR-1;
\item reforça choque e arco elétrico;
\item detalha segurança em projetos e atualização documental;
\item prevê estudo de energia incidente quando aplicável.
\end{itemize}
\column{0.46\textwidth}
\begin{caixaAt}[title=Como aplicar em agosto de 2026]
Não tratar a redação nova como se já estivesse vigente. No material:
\begin{itemize}
\item requisito atual é apresentado como atual;
\item texto publicado para 2027 é identificado como futuro;
\item boas práticas podem ser antecipadas sem falsa alegação normativa.
\end{itemize}
\end{caixaAt}
\end{columns}
\end{frame}

\begin{frame}{Ciclo completo de um projeto elétrico}
\small
\centering
\begin{tikzpicture}[node distance=0.32cm,every node/.style={font=\scriptsize}]
\tikzset{et/.style={draw=corPrim,fill=corDef!7,rounded corners,minimum width=2.55cm,minimum height=0.72cm,align=center}}
\node[et] (a) {1. Levantamento\\arquitetura e cargas};
\node[et,right=of a] (b) {2. Normas, risco\\e premissas};
\node[et,right=of b] (c) {3. Circuitos e\\unifilar};
\node[et,below=of c] (d) {4. Cálculos\\e proteção};
\node[et,left=of d] (e) {5. Quadros e\\infraestrutura};
\node[et,left=of e] (f) {6. Memorial, BIM\\e especificação};
\node[et,below=of f] (g) {7. Execução e\\fiscalização};
\node[et,right=of g] (h) {8. Ensaios e\\comissionamento};
\node[et,right=of h] (i) {9. As built e\\manutenção};
\foreach \u/\v in {a/b,b/c,c/d,d/e,e/f,f/g,g/h,h/i}{\draw[-{Stealth},thick,corPrim] (\u)--(\v);}
\draw[-{Stealth},thick,corAccent] (i.east) to[out=0,in=-70] node[right,align=center]{retroalimentar\\o projeto} (a.south east);
\end{tikzpicture}
\end{frame}

\begin{frame}{Entregáveis mínimos de um bom projeto}
\small
\begin{columns}[T,onlytextwidth]
\column{0.49\textwidth}
\begin{itemize}
\item Planta com pontos, rotas e identificação.
\item Quadro de cargas, demanda e balanceamento.
\item Diagramas unifilares e, quando necessário, multifilares.
\item Memorial de cálculo: corrente, cabos, queda, curto, proteção e aterramento.
\item Detalhes de quadros, BEP, DPS, SPDA e interfaces.
\end{itemize}
\column{0.49\textwidth}
\begin{itemize}
\item Especificação técnica e lista de materiais.
\item Estudos adicionais quando aplicáveis.
\item Plano de inspeção/comissionamento.
\item As built e registro de alterações.
\item ART/RRT e demais responsabilidades conforme competência profissional.
\end{itemize}
\end{columns}
\end{frame}

%=====================================================================

%=====================================================================
% PROJETO PREDIAL - LEITURA E INTERPRETACAO DA DOCUMENTACAO
%=====================================================================

\section{Leitura e interpretação de projetos elétricos}

\begin{frame}{Projeto elétrico é um conjunto coordenado de documentos}
\centering
\begin{tikzpicture}[
  b/.style={draw=corPrim,fill=corDef!7,rounded corners,minimum width=2.25cm,minimum height=0.75cm,align=center,font=\scriptsize},
  a/.style={-{Stealth},thick,corPrim}]
\node[b](pl)at(0,2.4){Plantas\\localização e rotas};
\node[b](dg)at(4.0,1.2){Diagramas\\lógica elétrica};
\node[b](qc)at(2.5,-1.8){Quadros e listas\\dados dos circuitos};
\node[b](mem)at(-2.5,-1.8){Memorial e cálculos\\critérios e resultados};
\node[b](det)at(-4.0,1.2){Detalhes\\como executar};
\node[draw=corAccent,fill=corAccent!12,rounded corners,minimum width=2.4cm,minimum height=0.85cm,align=center,font=\small\bfseries](obra)at(0,0.2){Instalação\\executável};
\foreach \x in {pl,dg,qc,mem,det}{\draw[a](\x)--(obra);}
\end{tikzpicture}
\vspace{2mm}
\begin{caixaAt}
Nenhum desenho deve ser lido isoladamente. A planta diz onde; o diagrama diz como; o quadro diz quanto; o memorial explica por quê; o detalhe mostra como executar.
\end{caixaAt}
\end{frame}

\begin{frame}{Antes de ler o desenho, leia a identificação}
\small
\begin{columns}[T,onlytextwidth]
\column{0.50\textwidth}
\begin{tikzpicture}[scale=0.78]
\draw[very thick] (0,0) rectangle (7.2,4.4);
\draw[thick] (4.45,0)--(4.45,1.45);
\draw[thick] (4.45,1.45)--(7.2,1.45);
\draw[thick] (5.85,0)--(5.85,1.45);
\node[font=\scriptsize\bfseries] at (5.83,1.15){PROJETO ELÉTRICO};
\node[font=\tiny,align=left,anchor=west] at (4.62,0.72){Prancha: PE-03\\Pavimento: térreo};
\node[font=\tiny,align=left,anchor=west] at (6.0,0.72){Rev.: 04\\Esc.: 1:50};
\node[font=\tiny,anchor=west] at (0.25,4.05){PLANTA DE ILUMINAÇÃO};
\draw[corAccent,thick,dashed] (4.55,0.15) rectangle (7.05,1.3);
\node[font=\tiny,corAccent] at (5.8,-0.35){carimbo e revisão};
\end{tikzpicture}
\column{0.48\textwidth}
\begin{enumerate}
\item título e finalidade da prancha;
\item edifício, pavimento e área;
\item código, escala e unidade;
\item revisão e data de emissão;
\item responsável e status do documento;
\item notas gerais e referências cruzadas;
\item legenda aplicável àquela prancha.
\end{enumerate}
\end{columns}
\begin{caixaAt}
Não executar prancha superada, sem revisão ou emitida apenas para estudo/coordenação.
\end{caixaAt}
\end{frame}

\begin{frame}{Linhas, símbolos e textos formam uma linguagem}
\small
\begin{tabularx}{\textwidth}{>{\bfseries}p{2.5cm}p{3.0cm}X}
\toprule
Elemento & Exemplo gráfico & O que interpretar \\
\midrule
Linha contínua & \tikz\draw[corPrim,thick](0,0)--(2.2,0); & ligação ou limite físico conforme a legenda \\
Linha tracejada & \tikz\draw[corAccent,thick,dashed](0,0)--(2.2,0); & comando, circuito, rota oculta ou referência \\
Seta & \tikz\draw[-{Stealth},corPrim,thick](0,0)--(2.2,0); & sentido funcional, fluxo ou continuação \\
Nó preenchido & \tikz\fill[corPrim](0,0)circle(2.4pt); & conexão elétrica efetiva \\
Cruzamento sem nó & \tikz{\draw[corPrim](0,-0.25)--(0,0.25);\draw[corPrim](-0.4,0)--(0.4,0);} & linhas cruzam sem conexão, salvo convenção diferente \\
Etiqueta C7 & \fbox{C7} & circuito que remete ao quadro de cargas \\
Referência 4/PE-08 & \fbox{4/PE-08} & detalhe 4 na prancha PE-08 \\
\bottomrule
\end{tabularx}
\begin{caixaObs}
Convenções variam. A legenda e as notas do projeto prevalecem sobre hábitos de leitura.
\end{caixaObs}
\end{frame}

\begin{frame}{Símbolo só tem significado completo com atributos}
\small
\begin{columns}[T,onlytextwidth]
\column{0.50\textwidth}
\begin{tikzpicture}[every node/.style={font=\scriptsize}]
\SimLuz{0.4}{3.9}\node[anchor=west]at(1.0,3.9){L1 / C1 / teto};
\SimInterruptor{0.4}{3.05}{S1}\node[anchor=west]at(1.5,3.05){comanda L1};
\SimTUG{0.4}{2.2}\node[anchor=west]at(1.0,2.2){T1 / C3 / 2P+T / 220 V};
\SimTUGM{0.4}{1.35}\node[anchor=west]at(1.0,1.35){T5 / C5 / bancada};
\SimTUE{0.4}{0.5}{E}\node[anchor=west]at(1.0,0.5){AC-1 / C12 / dedicado};
\end{tikzpicture}
\column{0.48\textwidth}
Para cada ponto, procurar:
\begin{itemize}
\item identificador único;
\item circuito e fase;
\item tensão e número de polos;
\item altura/cota de montagem;
\item potência ou equipamento;
\item comando associado;
\item nota ou detalhe construtivo.
\end{itemize}
\end{columns}
\begin{caixaAt}
Um círculo na planta não diz sozinho qual cabo, proteção, altura ou ligação usar. Esses dados vêm da etiqueta e dos documentos associados.
\end{caixaAt}
\end{frame}

\begin{frame}{Da arquitetura à planta elétrica}
\centering
\begin{tikzpicture}[scale=0.78]
\begin{scope}[xshift=-7.0cm]
\draw[very thick](0,0)rectangle(5.8,4.2);
\draw[thick](0,1.5)--(5.8,1.5);\draw[thick](3.7,1.5)--(3.7,4.2);
\node[font=\scriptsize\bfseries]at(1.85,3.7){SALA};
\node[font=\scriptsize\bfseries]at(4.75,3.7){COZINHA};
\node[font=\scriptsize\bfseries]at(2.9,0.75){QUARTO};
\draw[white,line width=5pt](1.0,1.5)--(1.9,1.5);\draw(1.0,1.5)--(1.9,2.4);
\node[font=\scriptsize]at(2.9,-0.55){planta arquitetônica};
\end{scope}
\draw[-{Stealth},very thick,corAccent](-0.8,2.1)--(0.8,2.1);
\begin{scope}[xshift=1.6cm]
\draw[very thick](0,0)rectangle(5.8,4.2);
\draw[thick](0,1.5)--(5.8,1.5);\draw[thick](3.7,1.5)--(3.7,4.2);
\node[font=\scriptsize\bfseries]at(1.85,3.7){SALA};
\node[font=\scriptsize\bfseries]at(4.75,3.7){COZINHA};
\node[font=\scriptsize\bfseries]at(2.9,0.75){QUARTO};
\draw[white,line width=5pt](1.0,1.5)--(1.9,1.5);\draw(1.0,1.5)--(1.9,2.4);
\SimQD{0.4}{2.1}{QD}
\SimLuz{1.9}{2.8}\SimLuz{4.75}{2.8}\SimLuz{2.9}{0.65}
\SimInterruptor{1.05}{1.75}{S1}
\foreach \x/\y in {0.25/2.3,3.45/2.3,4.0/1.8,5.55/2.2,0.25/0.3,5.55/0.3}{\SimTUG{\x}{\y}}
\draw[corAccent,dashed,thick](0.4,2.1)--(1.9,2.8)--(4.75,2.8);
\draw[corPrim,dashed,thick](0.4,2.1)--(0.25,2.3)--(3.45,2.3);
\node[font=\scriptsize]at(2.9,-0.55){planta elétrica coordenada};
\end{scope}
\end{tikzpicture}
\begin{caixaObs}
Portas definem comandos; mobiliário define tomadas; paredes e vigas limitam rotas; uso do ambiente define cargas e proteção.
\end{caixaObs}
\end{frame}

\begin{frame}{Siga um circuito da origem até cada ponto}
\centering
\begin{tikzpicture}[scale=0.82]
\draw[very thick](0,0)rectangle(11.8,5.4);
\draw[thick](4.2,0)--(4.2,5.4);\draw[thick](8.0,0)--(8.0,5.4);
\node[font=\scriptsize\bfseries]at(2.1,5.0){SALA};
\node[font=\scriptsize\bfseries]at(6.1,5.0){COZINHA};
\node[font=\scriptsize\bfseries]at(9.9,5.0){SERVIÇO};
\SimQD{0.55}{2.7}{QD1}
\SimLuz{2.1}{3.0}\SimLuz{6.1}{3.0}\SimLuz{9.9}{3.0}
\SimInterruptor{1.0}{0.45}{S1}\SimInterruptor{5.0}{0.45}{S2}\SimInterruptor{8.9}{0.45}{S3}
\draw[corAccent,line width=2.2pt,-{Stealth}](0.85,2.7)--(2.1,3.0)--(6.1,3.0)--(9.9,3.0);
\node[font=\scriptsize\bfseries,fill=white,corAccent]at(5.8,3.35){C1};
\draw[corNorma,dashed,thick](1.0,0.45)--(2.1,3.0);
\draw[corNorma,dashed,thick](5.0,0.45)--(6.1,3.0);
\draw[corNorma,dashed,thick](8.9,0.45)--(9.9,3.0);
\end{tikzpicture}
\vspace{1mm}
\begin{enumerate}
\item localizar C1 no quadro de cargas;
\item identificar origem, fase, neutro e PE no QD1;
\item seguir a rota e contar pontos/caixas;
\item conferir comandos S1--S3 e luminárias associadas;
\item cruzar seção, proteção, comprimento e queda.
\end{enumerate}
\end{frame}

\begin{frame}{A mesma instalação em quatro representações}
\scriptsize
\begin{columns}[T,onlytextwidth]
\column{0.24\textwidth}
\centering\textbf{Planta}\par\vspace{1mm}
\begin{tikzpicture}[scale=0.52]
\draw(0,0)rectangle(4,3);\SimQD{0.35}{1.5}{QD}\SimLuz{2.6}{1.7}\SimInterruptor{1.0}{0.35}{S}
\draw[corAccent,dashed,thick](0.65,1.5)--(2.6,1.7);\draw[corNorma,dashed](1.0,0.35)--(2.6,1.7);
\end{tikzpicture}
\column{0.24\textwidth}
\centering\textbf{Unifilar}\par\vspace{1mm}
\begin{circuitikz}[scale=0.58,transform shape]
\draw (0,2.6) node[left]{QD} to[short,o-] (0.8,2.6) to[ospst=$Q_1$] (2.0,2.6) to[lamp,l=$L_1$] (2.0,0) node[ground]{};
\end{circuitikz}
\column{0.24\textwidth}
\centering\textbf{Multifilar}\par\vspace{1mm}
\begin{circuitikz}[scale=0.50,transform shape]
\draw[corPrim] (0,3)--(3.4,3)node[right]{F};
\draw[corDef] (0,0)--(3.4,0)node[right]{N};
\draw[corNorma] (0,-0.55)--(3.4,-0.55)node[right]{PE};
\draw (0.65,3)to[ospst=$S_1$](2.25,3)to[lamp,l=$L_1$](2.25,0);
\end{circuitikz}
\column{0.24\textwidth}
\centering\textbf{Ligação física}\par\vspace{1mm}
\begin{tikzpicture}[scale=0.53]
\node[draw,minimum width=1.1cm,minimum height=1.0cm](s)at(0,0){S1};
\node[draw,circle,minimum size=0.9cm](l)at(2.8,2.3){L1};
\draw[corPrim,thick](s.west)--(-0.7,0)node[left]{F};
\draw[corAccent,thick](s.east)--(l.west);
\draw[corDef,thick](-0.7,2.7)node[left]{N}--(l.north);
\draw[corNorma,thick](-0.7,2.25)node[left]{PE}--(l.south);
\end{tikzpicture}
\end{columns}
\vspace{3mm}
\begin{caixaObs}
Planta localiza; unifilar resume; multifilar explicita condutores; ligação física mostra bornes e passagem real.
\end{caixaObs}
\end{frame}

\begin{frame}[shrink=4]{Como ler um diagrama unifilar}
\centering
\begin{tikzpicture}[
  b/.style={draw=corPrim,fill=corDef!7,rounded corners,minimum width=2.0cm,minimum height=0.7cm,align=center,font=\scriptsize}]
\node[b](orig)at(0,4.2){Rede 380/220 V};
\node[b](dg)at(0,3.1){DG 3P 63 A};
\node[b](dps)at(0,2.0){DPS + DR geral\\quando aplicável};
\node[b](qd)at(0,0.9){Barramentos\\F / N / PE};
\node[b](c1)at(-4.0,-0.8){C1 iluminação\\10 A / 1,5 mm$^2$};
\node[b](c2)at(-1.35,-0.8){C2 TUG\\16 A / 2,5 mm$^2$};
\node[b](c3)at(1.35,-0.8){C3 chuveiro\\32 A / 6 mm$^2$};
\node[b](c4)at(4.0,-0.8){C4 SAVE\\40 A / 10 mm$^2$};
\draw[-{Stealth},thick,corPrim](orig)--(dg)--(dps)--(qd);
\foreach \x in {c1,c2,c3,c4}{\draw[-{Stealth},thick](qd)--(\x);}
\end{tikzpicture}
\begin{caixaNorma}
Leia de cima para baixo: origem e tensão, seccionamento, proteção contra surtos/choque, barramentos, circuitos terminais e dados de cada saída.
\end{caixaNorma}
\end{frame}

\begin{frame}[shrink=4]{Como ler um diagrama multifilar}
\centering
\begin{circuitikz}[scale=0.92,transform shape]
\draw[corPrim,thick] (-5,2.4)--(5.2,2.4) node[right]{F};
\draw[corDef,thick] (-5,-0.9)--(5.2,-0.9) node[right]{N};
\draw[corNorma,thick] (-5,-1.55)--(5.2,-1.55) node[right]{PE};
\draw[corPrim] (-4.2,2.4) to[ospst,l=$Q_1$] (-2.3,2.4) to[ospst,l=$S_1$] (-0.3,2.4) to[lamp,l=$L_1$] (-0.3,-0.9);
\draw[corNorma,thick] (-0.3,-1.55)--(-0.3,-1.25);
\draw[corPrim] (1.0,2.4) to[ospst,l=$Q_2$] (2.9,2.4)--(2.9,0.35);
\draw[corDef] (2.9,0.35)--(2.9,-0.9);
\draw[corNorma] (3.45,0.35)--(3.45,-1.55);
\draw (2.55,0.35) rectangle (3.8,-0.35);
\node[font=\scriptsize] at (3.18,0){T1 2P+T};
\end{circuitikz}
\vspace{1mm}
\begin{columns}[T,onlytextwidth]
\column{0.49\textwidth}
\begin{itemize}
\item cada condutor é desenhado separadamente;
\item nós mostram conexões elétricas;
\item contatos aparecem na condição de repouso;
\end{itemize}
\column{0.49\textwidth}
\begin{itemize}
\item cores no desenho são didáticas, não substituem etiquetas;
\item acompanhar fase, neutro e PE até a carga;
\item conferir polos interrompidos e protegidos.
\end{itemize}
\end{columns}
\end{frame}

\begin{frame}{Funcional e ligação física respondem perguntas diferentes}
\small
\begin{columns}[T,onlytextwidth]
\column{0.49\textwidth}
\centering\textbf{Diagrama funcional}\par\vspace{1mm}
\begin{circuitikz}[scale=0.76,transform shape]
\draw (0,3)node[left]{F}to[ospst,l=$S_1$](2,3)to[lamp,l=$L_1$](2,0)node[left]{N};
\end{circuitikz}
\begin{itemize}
\item mostra a lógica de funcionamento;
\item simplifica caixas e percurso;
\item facilita análise de estados.
\end{itemize}
\column{0.49\textwidth}
\centering\textbf{Esquema de ligação}\par\vspace{1mm}
\begin{tikzpicture}[scale=0.7]
\node[draw,minimum width=1.3cm,minimum height=1.1cm](cx)at(0,1.4){4x2};
\node[draw,circle,minimum size=1.0cm](l)at(3.8,3.1){L1};
\draw[corPrim,thick](-1.3,1.4)node[left]{F}--(cx.west);
\draw[corAccent,thick](cx.east)--(l.west);
\draw[corDef,thick](-1.3,3.5)node[left]{N}--(l.north);
\draw[corNorma,thick](-1.3,2.9)node[left]{PE}--(l.south);
\end{tikzpicture}
\begin{itemize}
\item mostra caixas, bornes e condutores;
\item orienta passagem e terminação;
\item precisa corresponder ao funcional.
\end{itemize}
\end{columns}
\end{frame}

\begin{frame}{Layout do quadro não substitui o unifilar}
\small
\begin{columns}[T,onlytextwidth]
\column{0.49\textwidth}
\centering\textbf{Layout físico}\par\vspace{1mm}
\begin{tikzpicture}[scale=0.72]
\draw[very thick](0,0)rectangle(6.4,4.8);
\draw[thick](0.35,3.55)--(6.05,3.55);\draw[thick](0.35,1.75)--(6.05,1.75);
\foreach \x/\t in {0.75/DG,1.75/DPS,2.75/DR,3.75/C1,4.55/C2,5.35/C3}{
  \draw[fill=corDef!7,draw=corPrim](\x-0.35,2.2)rectangle(\x+0.35,3.3);
  \node[font=\tiny]at(\x,2.75){\t};}
\draw[corDef,thick](0.7,0.7)--(5.7,0.7)node[right,font=\tiny]{N};
\draw[corNorma,thick](0.7,0.35)--(5.7,0.35)node[right,font=\tiny]{PE};
\end{tikzpicture}
\column{0.49\textwidth}
\centering\textbf{Unifilar associado}\par\vspace{1mm}
\begin{tikzpicture}[
  b/.style={draw=corPrim,fill=corDef!7,minimum width=1.45cm,minimum height=0.55cm,font=\tiny,align=center}]
\node[b](dg)at(0,4.0){DG};\node[b](dps)at(0,3.1){DPS};\node[b](dr)at(0,2.2){DR};
\node[b](c1)at(-2.2,0.8){C1};\node[b](c2)at(0,0.8){C2};\node[b](c3)at(2.2,0.8){C3};
\draw[-{Stealth},thick,corPrim](dg)--(dps)--(dr);
\foreach \x in {c1,c2,c3}{\draw[-{Stealth},thick](dr)--(\x);}
\end{tikzpicture}
\end{columns}
\begin{caixaAt}
Layout indica posição e montagem. Unifilar indica relação elétrica. Ambos precisam usar os mesmos códigos e dispositivos.
\end{caixaAt}
\end{frame}

\begin{frame}[shrink=4]{Quadro de cargas conecta planta, diagrama e cálculo}
\scriptsize
\begin{tabularx}{\textwidth}{c X c c c c c c}
\toprule
C & Descrição & Pontos & Fase & $I_b$ & Cabo & DJ/DR & Prancha \\
\midrule
C1 & iluminação sala/cozinha & 7 & A & 4,8 A & 1,5 mm$^2$ & 10 A / RCBO & PE-02 \\
C2 & iluminação íntima & 6 & B & 6,4 A & 1,5 mm$^2$ & 10 A / RCBO & PE-02 \\
C3 & TUG sala/quartos & 17 & C & 7,7 A & 2,5 mm$^2$ & 16 A / DR & PE-03 \\
C5 & TUG cozinha & 6 & B & 9,5 A & 2,5 mm$^2$ & 16 A / DR & PE-03 \\
C9 & chuveiro social & 1 & A & 25,0 A & 6 mm$^2$ & 32 A / DR & PE-04 \\
\bottomrule
\end{tabularx}
\vspace{2mm}
\begin{caixaNorma}
Ao localizar C9 na planta, o quadro informa fase, corrente, seção e proteção; o unifilar mostra origem e dispositivo; o detalhe PE-04 mostra a ligação e a terminação.
\end{caixaNorma}
\end{frame}

\begin{frame}{Diagrama vertical: leia por pavimento e por sistema}
\centering
\begin{tikzpicture}[scale=0.8]
\draw[very thick](0,0)rectangle(10.5,5.2);
\foreach \y/\p in {0.8/Térreo,2.2/1º pav.,3.6/2º pav.,4.8/Cobertura}{
  \draw[black!35](0,\y)--(10.5,\y);\node[font=\tiny,anchor=east]at(-0.15,\y){\p};}
\draw[corPrim,line width=3pt](3.3,0.4)--(3.3,4.95);
\draw[corNorma,line width=3pt](6.7,0.4)--(6.7,4.95);
\node[font=\scriptsize\bfseries,corPrim]at(3.3,5.55){prumada de energia};
\node[font=\scriptsize\bfseries,corNorma]at(6.7,5.55){BEP / PE};
\foreach \y/\q [count=\i] in {0.8/QD-T,2.2/QD-1,3.6/QD-2}{
  \node[draw=corPrim,fill=corDef!7,minimum width=1.3cm,minimum height=0.48cm,font=\tiny](q\i)at(1.65,\y){\q};
  \draw[-{Stealth},thick,corPrim](3.3,\y)--(q\i);
  \draw[thick,corNorma](q\i)--(6.7,\y);}
\node[draw=corAccent,fill=corAccent!10,font=\tiny,align=center]at(8.7,4.35){SPDA\\cobertura};
\draw[thick,corNorma](8.7,4.0)--(6.7,3.6);
\end{tikzpicture}
\begin{caixaObs}
O vertical mostra origem, prumadas, quadros por pavimento, seções, proteção e equipotencialização. Ele não mostra a posição de cada tomada.
\end{caixaObs}
\end{frame}

\begin{frame}{Esquema de aterramento muda a leitura da proteção}
\scriptsize
\begin{columns}[T,onlytextwidth]
\column{0.32\textwidth}
\centering\textbf{TN-S}\par
\begin{tikzpicture}[scale=0.55]
\node[draw](f)at(0,2.8){Fonte};\node[draw](c)at(0,0){Carga};
\draw[corPrim,thick](f)--(-0.7,1.4)--(c);
\draw[corDef,thick](f)--(0,1.4)--(c);
\draw[corNorma,thick](f)--(0.7,1.4)--(c);
\node[font=\tiny]at(-0.9,1.4){F};\node[font=\tiny]at(-0.15,1.4){N};\node[font=\tiny]at(0.9,1.4){PE};
\end{tikzpicture}
\column{0.32\textwidth}
\centering\textbf{TT}\par
\begin{tikzpicture}[scale=0.55]
\node[draw](f)at(0,2.8){Fonte};\node[draw](c)at(0,0){Carga};
\draw[corPrim,thick](f)--(-0.55,1.4)--(c);
\draw[corDef,thick](f)--(0.15,1.4)--(c);
\draw[corNorma,thick](c.south)--(0,-0.8)node[ground]{};
\draw[corNorma,thick](f.south east)--(1.2,2.0)node[ground]{};
\node[draw=corAt,fill=corAt!7,font=\tiny]at(1.25,0.5){DR};
\end{tikzpicture}
\column{0.32\textwidth}
\centering\textbf{IT}\par
\begin{tikzpicture}[scale=0.55]
\node[draw](f)at(0,2.8){Fonte};\node[draw](c)at(0,0){Carga};
\draw[corPrim,thick](f)--(-0.45,1.4)--(c);
\draw[corDef,thick](f)--(0.25,1.4)--(c);
\draw[corNorma,thick](c.south)--(0,-0.8)node[ground]{};
\draw[dashed,corNorma](f.south east)--(1.2,2.0)node[ground]{};
\node[font=\tiny]at(1.0,1.25){isolado/impedância};
\end{tikzpicture}
\end{columns}
\begin{caixaAt}
Não interpretar proteção contra choque sem identificar o esquema. Caminho da falta, DR, equipotencialização e desligamento dependem dele.
\end{caixaAt}
\end{frame}

\begin{frame}{Referências cruzadas evitam repetir desenho}
\small
\begin{columns}[T,onlytextwidth]
\column{0.48\textwidth}
\begin{tikzpicture}[scale=0.76]
\draw[very thick](0,0)rectangle(6.5,4.5);
\node[draw=corAccent,circle,minimum size=0.75cm,font=\scriptsize\bfseries](d)at(3.2,2.4){4};
\draw[corAccent,thick](d)--(5.2,3.6);
\node[draw=corPrim,fill=corDef!7,font=\scriptsize]at(5.2,3.9){4/PE-08};
\node[font=\scriptsize]at(3.25,-0.55){planta PE-03};
\end{tikzpicture}
\column{0.50\textwidth}
\begin{tikzpicture}[scale=0.76]
\draw[very thick](0,0)rectangle(6.5,4.5);
\node[font=\scriptsize\bfseries]at(3.25,4.1){DETALHE 4 — PE-08};
\node[draw=corPrim,minimum width=2.1cm,minimum height=0.9cm]at(3.25,2.5){caixa / ligação};
\draw[corPrim,thick](3.25,2.05)--(3.25,0.7);
\node[font=\scriptsize]at(3.25,-0.55){prancha de detalhes};
\end{tikzpicture}
\end{columns}
\begin{caixaObs}
Leia o número do detalhe e a prancha de destino. Depois retorne à planta e confira se a condição detalhada é a mesma do ponto executado.
\end{caixaObs}
\end{frame}

\begin{frame}{Interprete por circuito, não apenas por prancha}
\small
\begin{enumerate}
\item escolha um circuito no quadro de cargas;
\item localize proteção e barramentos no unifilar;
\item encontre todos os pontos identificados na planta;
\item siga a rota e anote caixas, derivações e comprimentos;
\item abra os detalhes e esquemas de ligação associados;
\item confira seção, número de condutores e identificação;
\item verifique memorial, queda, DR, DPS e notas especiais;
\item confirme que a lista de materiais contempla o circuito.
\end{enumerate}
\begin{caixaProjeto}
Ao final, o aluno deve conseguir narrar o circuito: de onde vem, por onde passa, como é protegido, como funciona, onde termina e como será testado.
\end{caixaProjeto}
\end{frame}

\begin{frame}{Inconsistências entre planta, quadro e unifilar}
\scriptsize
\begin{tabularx}{\textwidth}{>{\bfseries}p{3.2cm}X}
\toprule
Conflito & Como perceber \\
\midrule
Circuito existe na planta, mas não no quadro & conferir todos os códigos Cx \\
Disjuntor e cabo divergem do unifilar & cruzar corrente, seção e proteção \\
Tomada sem tensão ou número de polos & verificar legenda e etiqueta \\
Interruptor não tem luminária associada & seguir linha de comando e identificação \\
\bottomrule
\end{tabularx}
\begin{caixaAt}
Não resolver divergência na obra por preferência pessoal. Emitir consulta técnica e registrar a resposta na revisão/as built.
\end{caixaAt}
\end{frame}

\begin{frame}{Inconsistências entre detalhes e execução}
\scriptsize
\begin{tabularx}{\textwidth}{>{\bfseries}p{3.2cm}X}
\toprule
Conflito & Como perceber \\
\midrule
PE ausente no esquema de ligação & conferir classe da carga e equipotencialização \\
Quadro sem espaço para dispositivos desenhados & comparar layout, módulos e reserva \\
Prumada não atende todos os pavimentos & conferir vertical e plantas por andar \\
Detalhe aponta para revisão diferente & validar carimbo e histórico de emissão \\
\bottomrule
\end{tabularx}
\begin{caixaAt}
Toda correção precisa alcançar o conjunto coordenado: prancha, quadro, diagrama, memorial, lista e as built.
\end{caixaAt}
\end{frame}

\begin{frame}{Competência de leitura esperada}
\small
Ao concluir esta parte, o aluno deve ser capaz de:
\begin{itemize}
\item reconhecer tipo, finalidade, revisão, escala, legenda e referências;
\item localizar pontos e rotas em plantas arquitetônicas e elétricas;
\item ler diagramas unifilares, multifilares, funcionais e verticais;
\item relacionar planta, quadro de cargas, memória e detalhes;
\item converter símbolo e etiqueta em condutores, bornes e dispositivos;
\item seguir um circuito completo da origem à carga;
\item detectar contradições antes de projetar ou executar;
\item explicar verbalmente o funcionamento e o critério de proteção.
\end{itemize}
\begin{caixaObs}
Ler um projeto não é decorar símbolos: é reconstruir a instalação e seu funcionamento a partir de documentos coordenados.
\end{caixaObs}
\end{frame}

\section{Fundamentos, potência, energia e demanda}

\begin{frame}{Grandezas em corrente alternada}
\small
\begin{columns}[T,onlytextwidth]
\column{0.54\textwidth}
\begin{align*}
v(t)&=V_m\sen(\omega t+\theta_v)\\
i(t)&=I_m\sen(\omega t+\theta_i)\\
V&=\frac{V_m}{\sqrt{2}},\qquad I=\frac{I_m}{\sqrt{2}}\\
\varphi&=\theta_v-\theta_i
\end{align*}
\begin{itemize}
\item RMS relaciona a grandeza CA ao efeito térmico equivalente em CC.
\item Em carga indutiva, corrente tende a atrasar a tensão.
\item Em carga capacitiva, corrente tende a adiantar a tensão.
\end{itemize}
\column{0.44\textwidth}
\begin{tikzpicture}
\begin{axis}[width=5.5cm,height=3.7cm,axis lines=middle,xlabel={$\omega t$},xtick={0,90,180,270,360},xticklabels={$0$,$90$,$180$,$270$,$360$},ytick=\empty,domain=0:360,samples=120,ymin=-1.3,ymax=1.3,grid=both]
\addplot[corPrim,thick] {sin(x)};
\addplot[corAccent,thick,dashed] {0.85*sin(x-35)};
\end{axis}
\node[corPrim,font=\scriptsize] at (1.25,3.25) {$v(t)$};
\node[corAccent,font=\scriptsize] at (4.25,2.45) {$i(t)$};
\end{tikzpicture}
\end{columns}
\end{frame}

\begin{frame}{Potência ativa, reativa e aparente}
\small
\begin{columns}[T,onlytextwidth]
\column{0.50\textwidth}
\begin{align*}
S &= VI \quad [\si{VA}]\\
P &= VI\cos\varphi \quad [\si{W}]\\
Q &= VI\sen\varphi \quad [\si{var}]\\
S^2 &= P^2+Q^2,\qquad fp=\frac{P}{S}
\end{align*}
\begin{itemize}
\item $P$: potência convertida em trabalho, calor, luz etc.
\item $Q$: intercâmbio de energia reativa.
\item $S$: capacidade elétrica exigida de cabos, quadros e transformadores.
\end{itemize}
\column{0.48\textwidth}
\centering
\begin{tikzpicture}[scale=0.9]
\draw[->,thick] (0,0)--(4.1,0) node[right]{$P$};
\draw[->,thick] (4.1,0)--(4.1,2.4) node[above]{$Q$};
\draw[->,thick,corPrim] (0,0)--(4.1,2.4) node[midway,above left]{$S$};
\draw (0.8,0) arc (0:30:0.8);
\node at (1.15,0.28) {$\varphi$};
\end{tikzpicture}
\end{columns}
\end{frame}

\begin{frame}{Sistemas trifásicos}
\small
\begin{columns}[T,onlytextwidth]
\column{0.52\textwidth}
\begin{align*}
P_{3\phi}&=\sqrt3\,V_{LL}I_L fp\\
Q_{3\phi}&=\sqrt3\,V_{LL}I_L\sen\varphi\\
S_{3\phi}&=\sqrt3\,V_{LL}I_L
\end{align*}
\begin{itemize}
\item Cargas equilibradas reduzem corrente de neutro fundamental.
\item Cargas monofásicas devem ser distribuídas entre as fases.
\item Harmônicas triplas podem somar no neutro mesmo com correntes eficazes semelhantes nas fases.
\end{itemize}
\column{0.46\textwidth}
\centering
\begin{tikzpicture}[scale=0.95]
\draw[->,thick,corPrim] (0,0)--(2,0) node[right]{$V_A$};
\draw[->,thick,corDef] (0,0)--({2*cos(120)},{2*sin(120)}) node[left]{$V_B$};
\draw[->,thick,corAccent] (0,0)--({2*cos(240)},{2*sin(240)}) node[left]{$V_C$};
\draw[thin,black!25] (0,0) circle (2);
\node at (0,-2.35) {defasagem de $120^\circ$};
\end{tikzpicture}
\end{columns}
\end{frame}

\begin{frame}{Corrente de projeto}
\small
\begin{columns}[T,onlytextwidth]
\column{0.53\textwidth}
\begin{caixaDef}
A corrente de projeto $I_b$ é a corrente prevista no circuito considerando potência, tensão, rendimento, fator de potência e regime de funcionamento.
\end{caixaDef}
\begin{align*}
I_{1\phi}&=\frac{P}{V\,fp\,\eta}\\
I_{3\phi}&=\frac{P}{\sqrt3\,V_{LL}\,fp\,\eta}
\end{align*}
\column{0.45\textwidth}
\begin{caixaEx}
Motor de $\SI{15}{kW}$, $\SI{380}{V}$, $fp=0{,}85$, $\eta=0{,}92$:
\[
I_b=\frac{15000}{\sqrt3\cdot380\cdot0{,}85\cdot0{,}92}\approx\SI{29.2}{A}.
\]
Esse valor ainda não resolve partida, cabo, disjuntor ou queda de tensão.
\end{caixaEx}
\end{columns}
\end{frame}

\begin{frame}{Carga instalada, demanda e diversidade}
\small
\begin{columns}[T,onlytextwidth]
\column{0.50\textwidth}
\begin{align*}
P_{inst}&=\sum P_k\\
D_{max}&=\max_t P(t)\\
F_d&=\frac{D_{max}}{P_{inst}}\\
F_{div}&=\frac{\sum D_{max,k}}{D_{max,grupo}}
\end{align*}
\begin{itemize}
\item A entrada não deve ser dimensionada supondo simultaneidade total sem justificativa.
\item Fator de demanda deve vir de regra aplicável, histórico, processo ou critério documentado.
\end{itemize}
\column{0.48\textwidth}
\begin{tikzpicture}
\begin{axis}[width=5.7cm,height=3.6cm,xlabel={hora},ylabel={$P$ (kW)},xmin=0,xmax=24,ymin=0,ymax=100,grid=both,xtick={0,6,12,18,24}]
\addplot[corPrim,thick,smooth] coordinates {(0,18)(3,15)(6,30)(9,54)(12,71)(15,63)(18,90)(21,50)(24,26)};
\addplot[corAccent,dashed,thick] coordinates {(0,90)(24,90)};
\end{axis}
\end{tikzpicture}
\end{columns}
\end{frame}

\begin{frame}{Levantamento de cargas moderno}
\scriptsize
\begin{tabularx}{\textwidth}{>{\bfseries}p{2.6cm}p{3.1cm}X}
\toprule
Carga & Dados mínimos & O que acrescentar ao projeto \\
\midrule
Iluminação & Potência, comando, ambiente & Iluminância, emergência, automação e eficiência. \\
TUG/TUE & Potência, tensão, uso & Circuitos, DR, simultaneidade, expansão. \\
Motores & kW, regime, partida & $I_N$, inrush, torque, proteção, VFD/soft-starter. \\
TI/eletrônica & W/VA, UPS & Harmônicas, DPS, continuidade, aterramento funcional. \\
Ar-condicionado & Potência e tecnologia & Inverter, corrente máxima, diversidade e partida. \\
VE & Potência de recarga, modo & Circuito dedicado, DR, gestão de carga, expansão. \\
FV/armazenamento & Potência, inversor & Proteções CA/CC, anti-ilhamento, fluxos bidirecionais. \\
Processo industrial & Ciclo e criticidade & Continuidade, redundância, seletividade, produção. \\
\bottomrule
\end{tabularx}
\end{frame}

%=====================================================================

\section{Projeto predial de baixa tensão}

\begin{frame}{Projeto predial: sequência correta}
\small
\centering
\begin{tikzpicture}[node distance=0.32cm,every node/.style={font=\scriptsize}]
\tikzset{et/.style={draw=corPrim,fill=corDef!7,rounded corners,minimum width=2.55cm,minimum height=0.72cm,align=center}}
\node[et] (a) {1. Planta e uso\\dos ambientes};
\node[et,right=of a] (b) {2. Pontos e\\cargas};
\node[et,right=of b] (c) {3. Circuitos\\terminais};
\node[et,below=of c] (d) {4. Condutores e\\eletrodutos};
\node[et,left=of d] (e) {5. Proteções\\e aterramento};
\node[et,left=of e] (f) {6. Quadro e\\balanceamento};
\node[et,below=of f] (g) {7. Entrada da\\concessionária};
\node[et,right=of g] (h) {8. Unifilar e\\memorial};
\node[et,right=of h] (i) {9. Verificação\\e as built};
\foreach \u/\v in {a/b,b/c,c/d,d/e,e/f,f/g,g/h,h/i}{\draw[-{Stealth},thick,corPrim] (\u)--(\v);}
\end{tikzpicture}
\vspace{2mm}
\begin{caixaAt}
Não se começa o projeto escolhendo disjuntor, cabo ou padrão de entrada. Esses elementos são consequência da arquitetura, das cargas, dos critérios normativos e das condições reais de instalação.
\end{caixaAt}
\end{frame}

\begin{frame}{Da arquitetura para os pontos elétricos}
\small
\begin{columns}[T,onlytextwidth]
\column{0.50\textwidth}
\begin{caixaDef}[title=Levantamento]
Antes de lançar símbolos elétricos, identificar:
\begin{itemize}
\item dimensões e perímetros dos ambientes;
\item portas, janelas e circulação;
\item bancadas, mobiliário e equipamentos fixos;
\item áreas molhadas, externas e sujeitas a lavagem;
\item posição tecnicamente adequada do QD.
\end{itemize}
\end{caixaDef}
\column{0.48\textwidth}
\begin{caixaProjeto}[title=Depois]
\begin{itemize}
\item calcular os mínimos normativos;
\item acrescentar pontos necessários ao uso real;
\item definir TUE pelas cargas previstas;
\item posicionar interruptores junto aos acessos;
\item lançar rotas compatíveis com a construção.
\end{itemize}
\end{caixaProjeto}
\end{columns}
\end{frame}

\begin{frame}{Iluminação e tomadas: mínimo não é projeto final}
\small
\begin{columns}[T,onlytextwidth]
\column{0.49\textwidth}
\begin{caixaNorma}[title=Previsão mínima de iluminação]
Em habitações, a previsão de carga parte de:
\begin{itemize}
\item 100 VA para os primeiros 6 m$^2$ ou fração;
\item 60 VA para cada acréscimo de 4 m$^2$ inteiros ou em fração.
\end{itemize}
A quantidade real de luminárias decorre do projeto luminotécnico, não dessa potência de previsão.
\end{caixaNorma}
\column{0.49\textwidth}
\begin{caixaNorma}[title=TUG]
A quantidade mínima depende do tipo e do perímetro do ambiente. Cozinha, copa, lavanderia e área de serviço possuem critérios mais exigentes que salas, dormitórios e outros cômodos.
\end{caixaNorma}
\begin{caixaAt}
Não desenhar tomadas apenas para atingir uma contagem: conferir portas, móveis, bancadas, eletrodomésticos e acessibilidade de uso.
\end{caixaAt}
\end{columns}
\end{frame}

\begin{frame}{Divisão de circuitos: critérios defensáveis}
\small
\begin{columns}[T,onlytextwidth]
\column{0.49\textwidth}
\begin{itemize}
\item separar iluminação e tomadas quando tecnicamente adequado;
\item separar cargas de potência elevada ou uso específico;
\item prever circuitos exclusivos para TUE quando requerido;
\item limitar a consequência de uma falha a uma área razoável;
\item facilitar identificação, manutenção e expansão.
\end{itemize}
\column{0.49\textwidth}
\begin{itemize}
\item respeitar ambientes molhados e externos;
\item evitar compartilhamento indevido de neutros;
\item considerar DR/DDR e sua seletividade funcional;
\item distribuir circuitos monofásicos entre as fases;
\item revisar simultaneidade e perfil das cargas.
\end{itemize}
\end{columns}
\begin{caixaObs}
A divisão de circuitos deve aparecer de forma idêntica na planta, no quadro de cargas, no diagrama unifilar e na identificação física do QD.
\end{caixaObs}
\end{frame}

\begin{frame}{Quadro de cargas: não usar números soltos}
\scriptsize
\begin{tabularx}{\textwidth}{c p{2.25cm} c c c c c c X}
\toprule
Circ. & Descrição & Fase & $P$ & $I_b$ & Cabo & DJ & DR & Verificações \\
 & & & W/VA & A & mm$^2$ & A & mA & \\
\midrule
C1 & Iluminação social & C & 720 & 3,27 & 1,5 & 10 & -- & $I_z$, queda, curto \\
C2 & Iluminação íntima & B & 860 & 3,91 & 1,5 & 10 & -- & $I_z$, queda, curto \\
C3 & TUG áreas secas & A & 1100 & 5,00 & 2,5 & 16 & conf. & uso e extensão \\
C4 & TUG cozinha & B & 1900 & 8,64 & 2,5 & 16 & 30 & ambiente molhado \\
C7 & Chuveiro & A & 5500 & 25,0 & 6 & 32 & 30 & dedicado \\
\bottomrule
\end{tabularx}
\vspace{2mm}
\begin{caixaAt}
Os valores acima pertencem ao projeto residencial desenvolvido adiante. Eles não são uma tabela universal para copiar em qualquer residência.
\end{caixaAt}
\end{frame}

\begin{frame}{Planta elétrica: o que precisa existir no desenho}
\small
\begin{columns}[T,onlytextwidth]
\column{0.47\textwidth}
\begin{itemize}
\item paredes, portas, janelas e nomes dos ambientes;
\item cotas ou referência dimensional confiável;
\item luminárias, interruptores e comandos;
\item TUG, TUE, dados e demais sistemas;
\item número do circuito junto aos pontos/rotas;
\item QD em local acessível e coerente.
\end{itemize}
\column{0.51\textwidth}
\begin{itemize}
\item eletrodutos ou rotas com leitura construtiva;
\item seção/quantidade de condutores quando pertinente;
\item alturas e observações de montagem;
\item legenda de símbolos;
\item compatibilização com hidráulica, estrutura, HVAC e dados;
\item detalhes dos locais especiais.
\end{itemize}
\begin{caixaProjeto}
Uma planta elétrica não é um diagrama abstrato colocado sobre retângulos. Ela deve permitir compreender \textbf{onde} instalar, \textbf{o que} instalar e \textbf{a qual circuito} cada ponto pertence.
\end{caixaProjeto}
\end{columns}
\end{frame}

\begin{frame}{Unifilar: fechar a lógica elétrica}
\small
\begin{columns}[T,onlytextwidth]
\column{0.48\textwidth}
\begin{itemize}
\item origem e tensão de fornecimento;
\item medição e proteção de entrada;
\item alimentador e número de condutores;
\item DG e barramentos do QD;
\item DPS e conexão ao sistema de equipotencialização;
\item DR/DDR/RCBO conforme circuitos protegidos.
\end{itemize}
\column{0.50\textwidth}
\begin{itemize}
\item disjuntores terminais;
\item fases atribuídas aos circuitos;
\item neutros corretamente segregados a jusante dos DR;
\item barramento PE contínuo e identificado;
\item seções e correntes nominais;
\item identificação coerente com planta e quadro de cargas.
\end{itemize}
\end{columns}
\begin{caixaAt}
O símbolo de terra isolado não substitui o desenho do PE, do BEP e da equipotencialização. Em projeto didático, esconder essa topologia ensina errado.
\end{caixaAt}
\end{frame}

%=====================================================================

%=====================================================================
% PROJETO PREDIAL - PREVISAO DE CARGAS E DIMENSIONAMENTO APLICADO
%=====================================================================

\section{Previsão de cargas e dimensionamento aplicado}

\begin{frame}{Cada ambiente precisa virar dados de projeto}
\centering
\begin{tikzpicture}[
  node distance=0.55cm and 0.35cm,
  b/.style={draw=corPrim,fill=corDef!7,rounded corners,minimum width=2.15cm,minimum height=0.78cm,align=center,font=\scriptsize},
  a/.style={-{Stealth},thick,corPrim}]
\node[b](arq){Arquitetura\\área e perímetro};
\node[b,right=of arq](uso){Uso real\\e equipamentos};
\node[b,right=of uso](prev){Previsão\\luz, TUG e TUE};
\node[b,below=of prev](circ){Circuitos\\e fases};
\node[b,left=of circ](dim){Condutor\\e proteção};
\node[b,left=of dim](doc){Planta, quadro\\e memorial};
\draw[a](arq)--(uso);\draw[a](uso)--(prev);\draw[a](prev)--(circ);
\draw[a](circ)--(dim);\draw[a](dim)--(doc);
\end{tikzpicture}
\vspace{4mm}
\begin{caixaAt}
Se um dado não aparece na planta, no quadro de cargas ou no memorial, ele não pode ser conferido na obra. Projeto completo transforma cada decisão em informação rastreável.
\end{caixaAt}
\end{frame}

\begin{frame}{Levantamento: dados que não podem ser presumidos}
\scriptsize
\begin{tabularx}{\textwidth}{>{\bfseries}p{2.6cm}X X}
\toprule
Grupo & Registrar & Impacto no projeto \\
\midrule
Arquitetura & área, perímetro, pé-direito, portas, janelas e mobiliário & pontos, comandos, rotas e comprimentos \\
Uso & permanência, umidade, acesso, público e manutenção & proteção, IP, DR, setorização e emergência \\
Equipamentos & potência, tensão, fases, $fp$, rendimento e regime & corrente, circuito, cabo e proteção \\
Ambiente & temperatura, insolação e agrupamento de linhas & capacidade de condução e vida útil \\
Conexão & tensão disponível, padrão, medição e limites da distribuidora & entrada, alimentador e documentação \\
Expansão & cargas futuras e espaços de reserva & quadro, eletroduto e demanda futura \\
\bottomrule
\end{tabularx}
\vspace{2mm}
\begin{caixaObs}
O levantamento termina com uma lista de premissas assinável: o que foi medido, informado pelo cliente, estimado ou deixado como reserva.
\end{caixaObs}
\end{frame}

\begin{frame}{Iluminação: previsão mínima em VA}
\small
Para habitações, a previsão mínima de carga de iluminação por cômodo é:
\[
S_{\mathrm{luz}}=
\begin{cases}
\SI{100}{VA}, & A\leq \SI{6}{m^2}\\[1mm]
\SI{100}{VA}+\SI{60}{VA}\left\lfloor\dfrac{A-6}{4}\right\rfloor,
& A>\SI{6}{m^2}
\end{cases}
\]
\begin{caixaEx}[title={Sala/jantar de 30,7 m²}]
\[
n=\left\lfloor\frac{30{,}7-6}{4}\right\rfloor=6
\qquad\Rightarrow\qquad
S_{\mathrm{luz}}=100+6\cdot60=\SI{460}{VA}
\]
\end{caixaEx}
\begin{caixaAt}
Os \SI{460}{VA} são previsão para dimensionamento. Não significam instalar luminárias que somem \SI{460}{W}.
\end{caixaAt}
\end{frame}

\begin{frame}{Luminotécnica: transformar iluminância em luminárias}
\small
\begin{columns}[T,onlytextwidth]
\column{0.48\textwidth}
No método dos lúmens:
\[
\Phi_{\mathrm{req}}=\frac{E_m A}{UF\,MF}
\]
\begin{itemize}
\item $E_m$: iluminância média de projeto;
\item $UF$: fator de utilização;
\item $MF$: fator de manutenção;
\item $\Phi$: fluxo luminoso total.
\end{itemize}
\column{0.50\textwidth}
\begin{caixaEx}[title=Critério declarado para a sala]
Adotando $E_m=\SI{150}{lx}$, $UF=0{,}75$ e $MF=0{,}80$:
\[
\Phi_{\mathrm{req}}=\frac{150\cdot30{,}7}{0{,}75\cdot0{,}80}
=\SI{7675}{lm}
\]
Seis luminárias de \SI{1300}{lm} fornecem \SI{7800}{lm}. A potência instalada vem do catálogo, não da previsão em VA.
\end{caixaEx}
\end{columns}
\begin{caixaObs}
Uniformidade, ofuscamento, reprodução de cor, temperatura de cor e cenas exigem verificação complementar ao método dos lúmens.
\end{caixaObs}
\end{frame}

\begin{frame}{TUG: quantidade e potência são cálculos diferentes}
\small
\begin{columns}[T,onlytextwidth]
\column{0.49\textwidth}
\begin{caixaNorma}[title=Cozinhas e locais análogos]
Nos critérios residenciais da NBR 5410, distribuir tomadas ao longo do perímetro e prever pelo menos uma a cada \SI{3.5}{m}, ou fração. A localização real deve atender bancadas e equipamentos.
\end{caixaNorma}
\[
N_{\mathrm{TUG}}=\left\lceil\frac{18{,}8}{3{,}5}\right\rceil=6
\]
\column{0.49\textwidth}
\begin{caixaEx}[title=Potência de previsão]
Para seis TUG de cozinha:
\[
S_{\mathrm{TUG}}=3(600)+3(100)=\SI{2100}{VA}
\]
Os três primeiros pontos recebem \SI{600}{VA} cada; os demais, \SI{100}{VA} cada.
\end{caixaEx}
\end{columns}
\begin{caixaAt}
O mínimo normativo não substitui a planta de mobiliário. Geladeira, micro-ondas, coifa e bancada precisam de pontos em posições utilizáveis.
\end{caixaAt}
\end{frame}

\begin{frame}{TUE e cargas específicas: usar dados de placa}
\scriptsize
\begin{tabularx}{\textwidth}{>{\bfseries}p{2.6cm}p{2.3cm}p{2.4cm}X}
\toprule
Carga & Dados mínimos & Circuito inicial & Verificações adicionais \\
\midrule
Chuveiro & $P$, $V$, comando & dedicado & DR, cabo, borne e queda \\
Ar-condicionado & $P_{in}$, $V$, fases, corrente & dedicado & partida/inverter e fabricante \\
Forno/cooktop & potência e simultaneidade & dedicado ou conforme fabricante & borne, ventilação e seccionamento \\
Bomba & $P_{out}$, $fp$, $\eta$, partida & dedicado & sobrecarga, falta de fase e comando \\
SAVE & potência, corrente e modo & dedicado & NBR 17019, DR compatível e gestão de carga \\
Fotovoltaico & potência e inversor & circuito de conexão & proteção CA/CC e norma da distribuidora \\
\bottomrule
\end{tabularx}
\vspace{2mm}
\begin{caixaObs}
Valores típicos servem apenas para anteprojeto. O executivo deve registrar fabricante, modelo e requisitos de instalação quando a carga já estiver definida.
\end{caixaObs}
\end{frame}

\begin{frame}[shrink=5]{Quadro de cargas começa no levantamento}
\scriptsize
\begin{tabularx}{\textwidth}{c X c c c c c}
\toprule
Circuito & Uso & Pontos & $S/P$ & $V$ & $I_b$ & Observação \\
\midrule
C1 & iluminação social/cozinha/serviço & 7 & 1,06 kVA & 220 & 4,8 A & RCBO \\
C2 & iluminação íntima/banhos/garagem & 6 & 1,40 kVA & 220 & 6,4 A & RCBO \\
C3 & TUG sala e dormitórios & 17 & 1,70 kVA & 220 & 7,7 A & área seca \\
C4 & TUG hall e garagem & 9 & 0,90 kVA & 220 & 4,1 A & uso geral \\
C5 & TUG cozinha & 6 & 2,10 kVA & 220 & 9,5 A & bancada \\
C6 & TUG serviço & 4 & 1,90 kVA & 220 & 8,6 A & área de serviço \\
C7 & TUG varanda gourmet & 4 & 1,90 kVA & 220 & 8,6 A & área externa coberta \\
C8 & TUG banheiros & 3 & 1,80 kVA & 220 & 8,2 A & DR obrigatório \\
C9--C16 & chuveiros, forno, lavadora e climatização & 8 & dados de placa & 220 & calcular & dedicados \\
\bottomrule
\end{tabularx}
\vspace{1mm}
\begin{caixaAt}
O quadro precisa informar método de instalação, seção, proteção, fase, comprimento, queda e critério de DR. Sem isso, ele é apenas inventário de cargas.
\end{caixaAt}
\end{frame}

\begin{frame}{Carga instalada, demanda e corrente não são sinônimos}
\small
\begin{columns}[T,onlytextwidth]
\column{0.52\textwidth}
\begin{align*}
P_{\mathrm{inst}}&=\sum_i P_i\\
P_{\mathrm{dem}}&=\sum_i f_{d,i}P_i\\
I_{b,1\phi}&=\frac{P}{Vfp\eta}\\
I_{b,3\phi}&=\frac{P}{\sqrt3V_{LL}fp\eta}
\end{align*}
\column{0.46\textwidth}
\begin{caixaNorma}[title=Dois usos distintos]
\begin{itemize}
\item a instalação interna usa correntes de projeto e simultaneidades justificadas;
\item a entrada segue o método e as faixas da distribuidora;
\item fatores não devem ser escolhidos apenas para reduzir cabo ou disjuntor.
\end{itemize}
\end{caixaNorma}
\end{columns}
\begin{caixaObs}
O projeto deve declarar qual potência alimenta cada equação: ativa de entrada, ativa de saída, aparente, instalada ou demandada.
\end{caixaObs}
\end{frame}

\begin{frame}{Divisão de circuitos: segurança, uso e manutenção}
\small
\begin{tabularx}{\textwidth}{>{\bfseries}p{3.0cm}X X}
\toprule
Decisão & Boa prática & Erro típico \\
\midrule
Iluminação e tomadas & separar funções e distribuir áreas & um único circuito para toda a residência \\
Áreas molhadas & agrupar apenas quando uso e proteção forem compatíveis & misturar cozinha, banheiro e área seca sem critério \\
Carga específica & circuito dedicado quando exigido ou tecnicamente necessário & ligar forno, chuveiro ou SAVE em TUG comum \\
Continuidade & evitar que uma falha apague toda a edificação & concentrar todos os pontos críticos no mesmo DR \\
Manutenção & identificar origem, destino e seccionamento & quadro sem legenda ou com circuitos genéricos \\
Expansão & reservar módulos e capacidade verificável & deixar ``reserva'' sem cabo, rota ou potência definida \\
\bottomrule
\end{tabularx}
\end{frame}

\begin{frame}{Balancear fases reduz corrente no neutro e assimetria}
\small
\begin{columns}[T,onlytextwidth]
\column{0.48\textwidth}
\begin{caixaEx}[title=Distribuição revisada]
\[
\begin{aligned}
S_A&=\SI{11.36}{kVA}\,,\\
S_B&=\SI{11.30}{kVA}\,,\\
S_C&=\SI{11.30}{kVA}\,.
\end{aligned}
\]
Em \SI{220}{V} fase-neutro:
\[
\begin{aligned}
I_A&=\SI{51.6}{A}\,,\\
I_B=I_C&=\SI{51.4}{A}\,.
\end{aligned}
\]
\end{caixaEx}
\column{0.50\textwidth}
\begin{tikzpicture}
\begin{axis}[
  ybar,bar width=13pt,width=6.0cm,height=4.2cm,
  ymin=0,ymax=13,xtick=data,xticklabels={A,B,C},
  ylabel={$S$ (kVA)},nodes near coords,
  every node near coord/.append style={font=\tiny},
  axis line style={corPrim},tick label style={font=\scriptsize}]
\addplot[fill=corDef!65,draw=corPrim] coordinates {(1,11.36) (2,11.30) (3,11.30)};
\end{axis}
\end{tikzpicture}
\end{columns}
\begin{caixaAt}
Balanceamento por potência é ponto de partida. Harmônicas triplas e cargas eletrônicas exigem análise adicional do neutro.
\end{caixaAt}
\end{frame}

\begin{frame}{Corrente de projeto: três exemplos}
\small
\begin{tabularx}{\textwidth}{>{\bfseries}p{2.5cm}X p{3.1cm}}
\toprule
Carga & Cálculo & Resultado \\
\midrule
Chuveiro resistivo & $I_b=5500/220$ & \SI{25.0}{A} \\
Forno com $fp=0{,}98$ & $I_b=3500/(220\cdot0{,}98)$ & \SI{16.2}{A} \\
Motor trifásico & $I_b=7500/(\sqrt3\cdot380\cdot0{,}86\cdot0{,}88)$ & \SI{15.1}{A} \\
\bottomrule
\end{tabularx}
\vspace{3mm}
\begin{caixaObs}
O disjuntor não é escolhido diretamente pela potência. Primeiro calcula-se $I_b$; depois verificam-se cabo, correções, partida, curva, capacidade de interrupção e fabricante.
\end{caixaObs}
\end{frame}

\begin{frame}{Capacidade de condução: corrigir antes de escolher a seção}
\small
\begin{columns}[T,onlytextwidth]
\column{0.48\textwidth}
Se a tabela fornece $I_{z,\mathrm{tab}}$, a capacidade corrigida é:
\[
I_z=I_{z,\mathrm{tab}}\,k_T\,k_G\,k_{outros}
\]
Ou, para consultar a tabela:
\[
I_{\mathrm{tab,req}}=\frac{I_n}{k_Tk_Gk_{outros}}
\]
\column{0.50\textwidth}
\begin{caixaEx}[title=Três circuitos agrupados]
Para $I_n=\SI{32}{A}$, $k_T=0{,}94$ e $k_G=0{,}70$:
\[
I_{\mathrm{tab,req}}=\frac{32}{0{,}94\cdot0{,}70}
=\SI{48.6}{A}
\]
Seleciona-se na tabela do método real uma seção cuja capacidade seja pelo menos \SI{48.6}{A}.
\end{caixaEx}
\end{columns}
\begin{caixaAt}
Não copiar ampacidade de catálogo sem identificar isolação, número de condutores carregados, método, temperatura e agrupamento.
\end{caixaAt}
\end{frame}

\begin{frame}[shrink=4]{Seções mínimas não encerram o dimensionamento}
\small
\begin{tabularx}{\textwidth}{>{\bfseries}p{3.6cm}c X}
\toprule
Uso do circuito & Cu mínimo & Ainda verificar \\
\midrule
Iluminação fixa & \SI{1.5}{mm^2} & ampacidade, queda, curto e proteção \\
Tomadas e força & \SI{2.5}{mm^2} & corrente, agrupamento e borne \\
Sinalização e controle & conforme aplicação & resistência mecânica e norma específica \\
Condutor de proteção & regra tabular ou adiabática & esquema, falta e desligamento \\
Alimentador & resultado dos critérios simultâneos & demanda, neutro, PE e expansão \\
\bottomrule
\end{tabularx}
\vspace{2mm}
\begin{caixaNorma}
Os valores acima são mínimos usuais da NBR 5410 para condutores de cobre em instalações fixas. A seção final é o maior resultado entre todos os critérios aplicáveis.
\end{caixaNorma}
\end{frame}

\begin{frame}{Queda de tensão: calcular o trecho real}
\small
Para circuito monofásico resistivo simplificado:
\[
\Delta V\%=100\frac{2\rho L I}{SV}
\]
\begin{caixaEx}[title=Chuveiro a \SI{28}{m} do quadro]
Com $\rho=\SI{0.0225}{\ohm.mm^2/m}$, $I=\SI{25}{A}$, $S=\SI{6}{mm^2}$ e $V=\SI{220}{V}$:
\[
\Delta V\%=100\frac{2(0{,}0225)(28)(25)}{6(220)}=2{,}39\%
\]
\[
\Delta V=0{,}0239\cdot220=\SI{5.26}{V}
\]
\end{caixaEx}
\begin{caixaAt}
Somar os trechos desde a origem considerada. Para cargas com baixo $fp$, usar resistência e reatância na temperatura de serviço.
\end{caixaAt}
\end{frame}

\begin{frame}{Curto-circuito: o cabo e o disjuntor precisam suportar a falta}
\small
\begin{columns}[T,onlytextwidth]
\column{0.49\textwidth}
Verificação térmica adiabática:
\[
S\geq\frac{I\sqrt{t}}{k}
\]
Corrente mínima de falta no fim do circuito:
\[
I_{k,\min}\approx\frac{U_0}{Z_{\text{fonte}}+Z_{\text{laço}}}
\]
\column{0.49\textwidth}
\begin{caixaEx}[title=Exemplo térmico]
Para $I=\SI{2}{kA}$, $t=\SI{0.10}{s}$ e $k=115$:
\[
S\geq\frac{2000\sqrt{0{,}10}}{115}=\SI{5.50}{mm^2}
\]
Adota-se \SI{6}{mm^2}, sujeito às demais verificações.
\end{caixaEx}
\end{columns}
\begin{caixaAt}
A capacidade de interrupção do disjuntor deve ser pelo menos igual à corrente de curto presumida no ponto de instalação.
\end{caixaAt}
\end{frame}

\begin{frame}{Coordenação completa: $I_b$, $I_n$, $I_z$ e atuação}
\small
\begin{columns}[T,onlytextwidth]
\column{0.48\textwidth}
\[
I_b\leq I_n\leq I_z
\qquad\text{e}\qquad
I_2\leq1{,}45I_z
\]
\begin{itemize}
\item $I_b$: corrente de projeto;
\item $I_n$: nominal/ajuste da proteção;
\item $I_z$: capacidade corrigida do cabo;
\item $I_2$: corrente convencional de atuação.
\end{itemize}
\column{0.50\textwidth}
\begin{caixaEx}[title=Circuito de chuveiro]
\[
I_b=25\,A,\quad I_n=32\,A,\quad I_z=41\,A
\]
\[
25\leq32\leq41
\]
Depois verificar curva/tempo, curto-circuito, queda, DR, borne e instruções do fabricante.
\end{caixaEx}
\end{columns}
\end{frame}

\begin{frame}{Neutro e PE também são dimensionados}
\small
\begin{columns}[T,onlytextwidth]
\column{0.49\textwidth}
\begin{caixaNorma}[title=Regra tabular usual do PE]
\begin{tabular}{cc}
\toprule
$S_f$ & $S_{PE}$ \\
\midrule
$S_f\leq16$ & $S_f$ \\
$16<S_f\leq35$ & \SI{16}{mm^2} \\
$S_f>35$ & $S_f/2$ \\
\bottomrule
\end{tabular}
\end{caixaNorma}
\column{0.49\textwidth}
\begin{caixaAt}[title=Neutro]
Não reduzir automaticamente quando houver:
\begin{itemize}
\item forte desequilíbrio monofásico;
\item harmônicas de ordem tripla;
\item cargas eletrônicas concentradas;
\item exigência do esquema ou da proteção.
\end{itemize}
\end{caixaAt}
\end{columns}
\begin{caixaObs}
O método adiabático do PE pode prevalecer. Verificar continuidade, impedância do laço e tempo de desligamento do esquema adotado.
\end{caixaObs}
\end{frame}

\begin{frame}{Eletroduto: a ocupação precisa ser calculada}
\small
Para três ou mais condutores, adota-se ocupação máxima de 40\% da área interna:
\[
A_{\mathrm{elet,min}}=\frac{\sum A_{\mathrm{cabos}}}{0{,}40}
\]
\begin{caixaEx}[title=Cinco cabos com diâmetro externo de \SI{4.8}{mm}]
\[
\sum A=5\frac{\pi(4{,}8)^2}{4}=\SI{90.5}{mm^2}
\]
\[
A_{\mathrm{elet,min}}=\frac{90{,}5}{0{,}40}=\SI{226.3}{mm^2}
\quad\Rightarrow\quad
D_{\mathrm{int,min}}=\sqrt{\frac{4A}{\pi}}=\SI{17.0}{mm}
\]
Escolher o diâmetro nominal usando o diâmetro interno real do fabricante e a facilidade de lançamento.
\end{caixaEx}
\end{frame}

\begin{frame}{Rotas, caixas e segregação também fazem parte do cálculo}
\small
\begin{columns}[T,onlytextwidth]
\column{0.50\textwidth}
\begin{tikzpicture}[scale=0.82]
\draw[corPrim,very thick] (0,0)--(2.2,0)--(2.2,2.0)--(5.2,2.0);
\draw[corAccent,very thick,dashed] (0,-0.45)--(3.0,-0.45)--(3.0,1.55)--(5.2,1.55);
\node[draw,fill=corDef!7,minimum size=0.65cm] at (0,0){QD};
\node[draw,fill=white,minimum size=0.55cm] at (2.2,0){C1};
\node[draw,fill=white,minimum size=0.55cm] at (2.2,2){C2};
\node[draw,fill=corAccent!10,minimum size=0.55cm] at (3,-0.45){D};
\node[anchor=west,font=\scriptsize] at (5.25,2){energia};
\node[anchor=west,font=\scriptsize] at (5.25,1.55){dados};
\node[font=\tiny,below] at (2.2,-0.4){caixa de passagem};
\end{tikzpicture}
\column{0.48\textwidth}
\begin{itemize}
\item limitar curvas e esforços de puxamento;
\item prever caixas acessíveis nos pontos de mudança;
\item separar energia, comunicação e sinais;
\item coordenar com estrutura, hidráulica, gás e climatização;
\item registrar comprimentos para queda e quantitativos.
\end{itemize}
\end{columns}
\begin{caixaAt}
Eletroduto é parte da linha elétrica. Uma rota impossível de executar invalida o dimensionamento que depende dela.
\end{caixaAt}
\end{frame}

\begin{frame}{Portão de qualidade antes de emitir o projeto}
\small
\begin{enumerate}
\item Todos os ambientes têm previsão de luz e TUG justificada?
\item Todas as cargas específicas têm dados elétricos e circuito definido?
\item O quadro fecha potência, corrente, fase, cabo, proteção e queda?
\item $I_b\leq I_n\leq I_z$ foi verificado com fatores de correção?
\item Curto-circuito, capacidade de interrupção e desligamento foram analisados?
\item Neutro, PE, DR, DPS, BEP e esquema de aterramento são coerentes?
\item Entrada, medição e documentação atendem à distribuidora vigente?
\item Planta, unifilar, memorial e lista de materiais dizem a mesma coisa?
\end{enumerate}
\begin{caixaProjeto}
Projeto pronto para obra é aquele que pode ser montado, inspecionado, ensaiado e mantido sem adivinhar decisões do projetista.
\end{caixaProjeto}
\end{frame}

\section{Condutores, infraestrutura e queda de tensão}

\begin{frame}{Dimensionamento: não existe um único critério}
\small
\begin{columns}[T,onlytextwidth]
\column{0.52\textwidth}
\begin{caixaNorma}[title=Fluxo de decisão]
A seção final deve atender simultaneamente a:
\begin{itemize}
\item seção mínima;
\item capacidade de condução de corrente;
\item queda de tensão;
\item curto-circuito/estabilidade térmica;
\item proteção contra choque/desligamento automático;
\item requisitos mecânicos e de instalação.
\end{itemize}
\end{caixaNorma}
\column{0.46\textwidth}
\centering
\begin{tikzpicture}[node distance=0.25cm,every node/.style={font=\scriptsize}]
\tikzset{b/.style={draw=corPrim,fill=corDef!7,rounded corners,align=center,minimum width=3.6cm,minimum height=0.58cm}}
\node[b] (a) {$I_b$ e regime};
\node[b,below=of a] (b) {Método de instalação};
\node[b,below=of b] (c) {Fatores de correção};
\node[b,below=of c] (d) {$I_b\le I_n\le I_z$};
\node[b,below=of d] (e) {Queda de tensão};
\node[b,below=of e] (f) {Curto + PE + desligamento};
\foreach \u/\v in {a/b,b/c,c/d,d/e,e/f}{\draw[-{Stealth},thick,corPrim] (\u)--(\v);}
\end{tikzpicture}
\end{columns}
\end{frame}

\begin{frame}{Coordenação carga--proteção--condutor}
\small
\begin{columns}[T,onlytextwidth]
\column{0.51\textwidth}
\begin{align*}
I_b &\le I_n \le I_z\\
I_2 &\le 1{,}45 I_z
\end{align*}
\begin{itemize}
\item $I_b$: corrente de projeto.
\item $I_n$: corrente nominal/ajuste do dispositivo.
\item $I_z$: capacidade do condutor após correções.
\item $I_2$: corrente convencional de atuação para sobrecarga.
\end{itemize}
\column{0.47\textwidth}
\centering
\begin{tikzpicture}[scale=0.9,every node/.style={font=\scriptsize}]
\draw[->] (0,0)--(5,0) node[right]{$I$};
\draw[thick,corDef] (1.0,0.1)--(1.0,1.1) node[above]{$I_b$};
\draw[thick,corAccent] (2.0,0.1)--(2.0,1.6) node[above]{$I_n$};
\draw[thick,corNorma] (3.3,0.1)--(3.3,1.2) node[above]{$I_z$};
\draw[decorate,decoration={brace,amplitude=5pt}] (1,-0.25)--(3.3,-0.25) node[midway,below=5pt]{coordenação};
\end{tikzpicture}
\end{columns}
\end{frame}

\begin{frame}{Fatores de correção térmica}
\small
\begin{tabularx}{\textwidth}{>{\bfseries}p{3.2cm}X}
\toprule
Fator & Consequência \\
\midrule
Temperatura ambiente & Aumenta temperatura de operação e reduz capacidade admissível. \\
Agrupamento & Circuitos próximos aquecem mutuamente. \\
Isolação térmica & Pode impedir dissipação e exigir forte redução de corrente. \\
Método de instalação & Eletroduto, leito, bandeja, enterrado e ao ar têm comportamentos diferentes. \\
Tipo de isolação & PVC, EPR/XLPE etc. possuem limites térmicos distintos. \\
Harmônicas & Aumentam perdas e podem elevar corrente no neutro. \\
\bottomrule
\end{tabularx}
\vspace{2mm}
\begin{caixaAt}
A corrente de catálogo do cabo não pode ser copiada sem identificar o método de referência e os fatores aplicáveis.
\end{caixaAt}
\end{frame}

\begin{frame}{Neutro em presença de harmônicas triplas}
\small
\begin{columns}[T,onlytextwidth]
\column{0.50\textwidth}
\begin{itemize}
\item Em sistema equilibrado senoidal, a fundamental do neutro tende a se cancelar.
\item Componentes de ordem 3, 9, 15... são de sequência zero e podem \textbf{somar} no neutro.
\item Fontes chaveadas monofásicas em grande quantidade exigem atenção especial.
\end{itemize}
\begin{caixaAt}
Não reduzir o neutro automaticamente em escritórios, data centers, laboratórios ou instalações com muita eletrônica.
\end{caixaAt}
\column{0.48\textwidth}
\begin{tikzpicture}
\begin{axis}[width=5.7cm,height=3.7cm,ybar,bar width=10pt,xlabel={ordem},ylabel={$I_h/I_1$},ymin=0,ymax=1.1,xtick={1,3,5,7,9},grid=major]
\addplot coordinates {(1,1)(3,0.55)(5,0.20)(7,0.14)(9,0.12)};
\end{axis}
\end{tikzpicture}
\end{columns}
\end{frame}

\begin{frame}[shrink=5]{Queda de tensão}
\small
\begin{columns}[T,onlytextwidth]
\column{0.53\textwidth}
\begin{align*}
\Delta V_{1\phi}&\approx2IL(R'\cos\varphi+X'\sen\varphi)\\
\Delta V_{3\phi}&\approx\sqrt3IL(R'\cos\varphi+X'\sen\varphi)\\
\Delta V(\%)&=100\frac{\Delta V}{V_n}
\end{align*}
\begin{itemize}
\item $L$ é o comprimento de ida; $R'$ e $X'$ são valores por unidade de comprimento, na condição de operação considerada. Usar unidades compatíveis, por exemplo km com $\Omega$/km.
\item Em $1\phi$, o fator 2 já representa ida e retorno; não duplicar novamente o comprimento.
\item Usar $V_n=V_{FN}$ no circuito monofásico e $V_n=V_{LL}$ no trifásico.
\item Circuitos longos podem ser dimensionados pela queda, não pela ampacidade.
\end{itemize}
\column{0.45\textwidth}
\begin{caixaEx}
Trifásico: $I=\SI{40}{A}$, $L=\SI{80}{m}$, $R'=\SI{3.08}{\ohm/km}$, $X'=\SI{0.08}{\ohm/km}$, $fp=0{,}85$.
\[
\Delta V\approx\SI{14.7}{V}\Rightarrow 3{,}9\%\;@\;\SI{380}{V}.
\]
A aceitação depende do limite aplicável ao trecho e ao sistema completo.
\end{caixaEx}
\end{columns}
\end{frame}

\begin{frame}{Efeito do comprimento: corrente de falta e queda}
\small
\centering
\begin{tikzpicture}
\begin{axis}[width=10.8cm,height=4.6cm,xlabel={distância ao quadro/fonte (m)},ylabel={valor normalizado},xmin=0,xmax=250,ymin=0,ymax=1.1,grid=both,legend pos=north east]
\addplot[corPrim,thick,smooth] coordinates {(0,1)(25,0.82)(50,0.67)(75,0.56)(100,0.47)(150,0.36)(200,0.29)(250,0.24)};
\addlegendentry{$I_{cc}$ relativo}
\addplot[corAccent,thick,dashed,smooth] coordinates {(0,0)(25,0.10)(50,0.20)(75,0.30)(100,0.40)(150,0.60)(200,0.80)(250,1.00)};
\addlegendentry{$\Delta V$ relativo}
\end{axis}
\end{tikzpicture}
\begin{caixaObs}
Quanto mais distante: maior queda de tensão e, em geral, menor corrente de curto. Isso pode alterar seção e tempo de desligamento.
\end{caixaObs}
\end{frame}

\begin{frame}{Curto-circuito: verificação térmica do condutor}
\small
\begin{columns}[T,onlytextwidth]
\column{0.51\textwidth}
\begin{align*}
S &\ge \frac{I_k\sqrt{t}}{k}
\end{align*}
\begin{itemize}
\item $S$: seção do condutor.
\item $I_k$: corrente de falta considerada.
\item $t$: duração até eliminação.
\item $k$: constante relacionada ao material e isolação.
\end{itemize}
\column{0.47\textwidth}
\begin{caixaEx}
Se um PE conduz $\SI{4}{kA}$ por $\SI{0.2}{s}$ e o valor adotado para $k$ for $115$:
\[
S\ge\frac{4000\sqrt{0{,}2}}{115}\approx\SI{15.6}{mm^2}.
\]
Escolhe-se seção comercial e verifica-se o conjunto de critérios.
\end{caixaEx}
\end{columns}
\end{frame}

\begin{frame}{Condutor de proteção: seção também é calculada}
\small
\begin{columns}[T,onlytextwidth]
\column{0.50\textwidth}
\begin{caixaNorma}[title=Método tabular simplificado]
Para condutores de fase e PE do mesmo material e nas condições previstas pelo método:
\[
S_{PE}=\begin{cases}
S_F, & S_F\le\SI{16}{mm^2}\\
\SI{16}{mm^2}, & \SI{16}{mm^2}<S_F\le\SI{35}{mm^2}\\
S_F/2, & S_F>\SI{35}{mm^2}
\end{cases}
\]
\end{caixaNorma}
\column{0.48\textwidth}
\begin{caixaAt}[title=O que ainda precisa fechar]
\begin{itemize}
\item seção mínima e proteção mecânica;
\item condição adiabática $S\ge I_k\sqrt{t}/k$;
\item impedância do caminho de falta e tempo de desligamento;
\item conexões, continuidade e identificação;
\item requisitos específicos de equipotencialização.
\end{itemize}
\end{caixaAt}
\end{columns}
\end{frame}

\begin{frame}{Eletrodutos, bandejas e segregação}
\small
\begin{columns}[T,onlytextwidth]
\column{0.50\textwidth}
\begin{itemize}
\item Ocupação deve permitir lançamento e substituição de cabos.
\item Respeitar raios de curvatura e esforços de tração.
\item Separar potência, controle, instrumentação e dados conforme compatibilidade e isolamento.
\item Considerar expansão e acessibilidade.
\item Bandejas metálicas precisam integrar o conceito de equipotencialização quando aplicável.
\end{itemize}
\column{0.48\textwidth}
\centering
\begin{tikzpicture}[scale=0.78]
\draw[thick,corPrim] (0,0) rectangle (5.8,2.1);
\foreach \x in {0.25,0.65,...,5.55}{\draw[black!25] (\x,0)--(\x,2.1);}
\foreach \x/\y/\c in {0.7/1.45/corPrim,1.5/0.7/corAccent,2.5/1.3/corAt,3.4/0.55/corNorma,4.3/1.5/corDef,5.1/0.85/corPrim}{\fill[\c!70] (\x,\y) circle (0.18);\draw (\x,\y) circle (0.18);}
\node at (2.9,-0.45) {leito/bandeja: rota também é parte do projeto};
\end{tikzpicture}
\end{columns}
\end{frame}

\begin{frame}{Identificação de condutores e circuitos}
\small
\begin{columns}[T,onlytextwidth]
\column{0.50\textwidth}
\begin{tabularx}{\textwidth}{>{\bfseries}p{2.1cm}X}
\toprule
Elemento & Princípio \\
\midrule
PE & Identificação exclusiva de proteção. \\
Neutro & Identificação própria e coerente. \\
Fases & Padrão consistente em toda a instalação. \\
Retornos & Anilhas/etiquetas; não confundir com PE/N. \\
Cabos & Origem, destino, circuito e seção. \\
\bottomrule
\end{tabularx}
\column{0.48\textwidth}
\begin{caixaProjeto}[title=As built útil]
Registrar:
\begin{itemize}
\item identificação real;
\item bornes e dispositivos;
\item ajustes de proteção;
\item cabos instalados;
\item alterações de campo;
\item reservas e circuitos desativados.
\end{itemize}
\end{caixaProjeto}
\end{columns}
\end{frame}

%=====================================================================

\section{Proteção, DR, DPS e aterramento}

\begin{frame}{Proteções: cada dispositivo resolve um problema}
\scriptsize
\begin{tabularx}{\textwidth}{>{\bfseries}p{3.0cm}X}
\toprule
Dispositivo & Função principal \\
\midrule
Disjuntor/fusível & Sobrecarga e curto-circuito; deve ter capacidade de interrupção adequada. \\
DR/IDR/RCBO & Corrente diferencial residual: choque/fuga conforme função, tipo e sensibilidade. \\
DPS & Limita sobretensões transitórias; atua em coordenação com aterramento, equipotencialização e SPDA. \\
Seccionadora & Isolação/manobra para operação e manutenção conforme aplicação. \\
Contator & Manobra frequente; não substitui proteção de curto. \\
Relé de motor & Sobrecarga, falta de fase, travamento e outras funções. \\
Proteção eletrônica & Ajustes de longo/curto/instantâneo, falta à terra, comunicação etc. \\
\bottomrule
\end{tabularx}
\end{frame}

\begin{frame}{Proteção contra choque: camadas}
\small
\begin{columns}[T,onlytextwidth]
\column{0.31\textwidth}
\begin{caixaDef}[title=Proteção básica]
Evitar contato com partes vivas em condição normal: isolação, barreiras, invólucros etc.
\end{caixaDef}
\column{0.31\textwidth}
\begin{caixaNorma}[title=Proteção supletiva]
Limitar risco se ocorrer falha: PE, equipotencialização, desligamento automático etc.
\end{caixaNorma}
\column{0.34\textwidth}
\begin{caixaAt}[title=Proteção adicional]
DR de alta sensibilidade em situações previstas. É complemento, não substituto das medidas básicas.
\end{caixaAt}
\end{columns}
\vspace{3mm}
\begin{caixaObs}
A proteção contra choque é um sistema de medidas coordenadas; não existe ``resolver tudo instalando um DR''.
\end{caixaObs}
\end{frame}

\begin{frame}{DR: princípio de funcionamento}
\small
\begin{columns}[T,onlytextwidth]
\column{0.50\textwidth}
\begin{caixaDef}
O DR compara a soma vetorial das correntes dos condutores ativos que atravessam seu toroide. Uma fuga que não retorna por esses condutores produz corrente residual.
\end{caixaDef}
\[
I_{\Delta}=\left|\sum I_k\right|.
\]
\begin{itemize}
\item Neutros de circuitos protegidos não podem ser misturados a jusante.
\item PE não atravessa o toroide como condutor ativo.
\end{itemize}
\column{0.48\textwidth}
\begin{circuitikz}[scale=0.76,every node/.style={font=\scriptsize}]
\draw (0,3) node[left]{F} -- (1,3) to[generic,l=DR] (2.4,3) -- (4.8,3) to[R,l=Carga] (4.8,1.2) -- (2.4,1.2) to[generic] (1,1.2) -- (0,1.2) node[left]{N};
\draw[->,thick,corPrim] (2.6,3.2)--(3.6,3.2) node[midway,above]{$I_F$};
\draw[->,thick,corPrim] (3.6,1.0)--(2.6,1.0) node[midway,below]{$I_N$};
\draw[->,thick,corAt] (4.8,2.1)--(5.6,1.2) node[right]{fuga};
\draw[corNorma,thick] (5.6,1.2)--(5.6,0.55)--(0.6,0.55) node[left]{PE};
\end{circuitikz}
\end{columns}
\end{frame}

\begin{frame}[shrink=8]{Tipos de DR e eletrônica de potência}
\scriptsize
\begin{tabularx}{\textwidth}{>{\bfseries}p{1.4cm}p{4.3cm}X}
\toprule
Tipo & Sensibilidade funcional & Aplicações/alertas \\
\midrule
AC & Corrente residual senoidal CA. & Cargas simples; pode ser inadequado para eletrônica moderna. \\
A & CA senoidal + componentes pulsantes CC previstas. & Eletrodomésticos e muitas cargas eletrônicas. \\
F & Inclui características para determinadas cargas monofásicas com conversores de frequência. & Alguns equipamentos inverter/velocidade variável. \\
B & Detecta também correntes residuais CC lisas dentro do escopo do dispositivo. & Pode ser requerido em VFD, FV, VE e outros conversores; não é escolha automática sem analisar o equipamento e a norma específica. \\
\bottomrule
\end{tabularx}
\vspace{2mm}
\begin{caixaAt}
A escolha não deve ser feita apenas por ``30 mA''. Verificar \textbf{tipo de forma de onda residual}, seletividade, imunidade, corrente de fuga normal, detecção de CC incorporada ao equipamento e norma específica. Em VE, por exemplo, uma solução com detecção de \SI{6}{mA} CC incorporada não é equivalente a instalar indiscriminadamente qualquer DR tipo A.
\end{caixaAt}
\end{frame}

\begin{frame}[shrink=4]{DPS: parâmetros e topologia que importam}
\small
\begin{columns}[T,onlytextwidth]
\column{0.47\textwidth}
\begin{itemize}
\item $U_c$: máxima tensão contínua aplicável.
\item $U_p$: nível de proteção de tensão.
\item $I_n$: corrente nominal de descarga em ensaio aplicável.
\item $I_{imp}$: corrente de impulso para DPS tipo 1.
\item Esquema de conexão compatível com o aterramento da instalação.
\item Proteção de retaguarda e capacidade de curto conforme fabricante.
\item Condutores de conexão tão curtos e retilíneos quanto praticável.
\end{itemize}
\column{0.51\textwidth}
\centering
\begin{tikzpicture}[scale=0.78,every node/.style={font=\scriptsize}]
\draw[thick] (0,3.0)--(5.3,3.0); \node[left] at (0,3.0){F};
\draw[thick] (0,2.1)--(5.3,2.1); \node[left] at (0,2.1){N};
\draw[thick,corNorma] (0,0.7)--(5.3,0.7); \node[left,corNorma] at (0,0.7){PE};
\node[draw=corAt,fill=corAt!8,rounded corners,minimum width=1.05cm,minimum height=0.55cm] (d1) at (2.0,1.55){DPS};
\node[draw=corAt,fill=corAt!8,rounded corners,minimum width=1.05cm,minimum height=0.55cm] (d2) at (3.7,1.55){DPS};
\draw[corAt,thick] (2.0,3.0)--(d1.north); \draw[corNorma,thick] (d1.south)--(2.0,0.7);
\draw[corAt,thick] (3.7,2.1)--(d2.north); \draw[corNorma,thick] (d2.south)--(3.7,0.7);
\node[font=\tiny,align=center] at (2.85,0.15){PE segue ao BEP / sistema\\de equipotencialização};
\end{tikzpicture}
\end{columns}
\begin{caixaObs}
O desenho é conceitual. A combinação F--PE, N--PE, F--N ou 3+1 depende do esquema de aterramento, do DPS selecionado e da coordenação requerida. Não se cria um ``terra local'' isolado para o DPS.
\end{caixaObs}
\end{frame}

\begin{frame}{Coordenação de DPS em cascata}
\small
\centering
\begin{tikzpicture}[scale=0.83,every node/.style={font=\scriptsize}]
\tikzset{b/.style={draw=corPrim,fill=corDef!7,rounded corners,minimum width=2.5cm,minimum height=0.7cm,align=center}}
\node[b] (a) at (0,2.2) {Entrada / QGBT\\DPS tipo 1 ou 1+2};
\node[b] (b) at (4.0,2.2) {QD\\DPS tipo 2};
\node[b] (c) at (8.0,2.2) {Equipamento\\proteção fina};
\draw[-{Stealth},thick,corPrim] (a)--(b);
\draw[-{Stealth},thick,corPrim] (b)--(c);
\node[font=\tiny\bfseries,corPrim,fill=white,inner sep=1pt] at (2.0,2.82){alimentador};
\node[font=\tiny\bfseries,corPrim,fill=white,inner sep=1pt] at (6.0,2.82){circuito};
% PE/equipotencialização contínuos, sem eletrodos independentes
\draw[corNorma,very thick] (-0.7,0.7)--(8.7,0.7);
\node[left,corNorma] at (-0.7,0.7){PE/BEP};
\draw[corNorma,thick] (0,1.85)--(0,0.7);
\draw[corNorma,thick] (4,1.85)--(4,0.7);
\draw[corNorma,thick] (8,1.85)--(8,0.7);
\draw[corNorma,thick] (1.0,0.7)--(1.0,-0.1) node[ground]{};
\node[font=\tiny,corNorma] at (2.35,0.25){sistema comum de equipotencialização/aterramento};
\end{tikzpicture}
\vspace{2mm}
\begin{caixaNorma}[title=Coordenação]
A cascata deve considerar categoria/local da proteção, tensão suportável do equipamento, distância, energia do surto, presença de SPDA, esquema de aterramento e instruções dos fabricantes.
\end{caixaNorma}
\end{frame}

\begin{frame}[shrink=2]{Funções do aterramento}
\small
\begin{columns}[T,onlytextwidth]
\column{0.52\textwidth}
\begin{itemize}
\item Referência e caminho de corrente de falta conforme esquema.
\item Limitação de tensões de toque e passo.
\item Viabilização do desligamento automático.
\item Equipotencialização de massas e elementos condutivos.
\item Dissipação/controle de surtos em conjunto com DPS/SPDA.
\item Compatibilidade eletromagnética e referência funcional quando aplicável.
\end{itemize}
\column{0.46\textwidth}
\centering
\begin{tikzpicture}[scale=0.82,every node/.style={font=\scriptsize}]
\draw[thick] (0,0) rectangle (3.8,2.0);
\node at (1.9,1.0) {Quadro / BEP};
\draw[thick,corNorma] (1.9,0) -- (1.9,-0.8) -- (0.6,-0.8) -- (0.6,-1.8);
\draw[thick,corNorma] (1.9,-0.8) -- (3.2,-0.8) -- (3.2,-1.8);
\foreach \x in {0.6,3.2}{\draw[thick,corNorma] (\x,-1.8)--(\x,-2.7);\draw[thick,corNorma] (\x-0.3,-2.7)--(\x+0.3,-2.7);}
\draw[corNorma,thick] (0.6,-2.25)--(3.2,-2.25);
\node[corNorma] at (1.9,-3.05) {eletrodos / malha};
\end{tikzpicture}
\end{columns}
\begin{caixaAt}
``Aterramento = uma haste'' é uma simplificação perigosa. O desempenho depende do esquema completo, impedâncias, equipotencialização e dispositivos de proteção.
\end{caixaAt}
\end{frame}

\begin{frame}{Esquemas TN, TT e IT}
\scriptsize
\begin{tabularx}{\textwidth}{>{\bfseries}p{1.5cm}p{3.5cm}X}
\toprule
Esquema & Característica & Pontos de projeto \\
\midrule
TN-S & N e PE separados. & Boa previsibilidade da corrente de falta; atenção a laços e EMC. \\
TN-C & PEN combinado em parte permitida do sistema. & Restrições importantes; não é solução genérica para circuitos terminais. \\
TN-C-S & PEN a montante e separação posterior de N/PE. & Ponto de separação deve ser controlado e não se recombina N com PE a jusante. \\
TT & Massas com eletrodo local, distinto funcionalmente do aterramento da fonte. & DR costuma ser essencial para desligamento automático. \\
IT & Fonte isolada da terra ou ligada por impedância. & Alta continuidade na primeira falta; exige monitoramento de isolamento e estratégia de localização. \\
\bottomrule
\end{tabularx}
\end{frame}

\begin{frame}{Equipotencialização principal e suplementar}
\small
\begin{columns}[T,onlytextwidth]
\column{0.49\textwidth}
\begin{caixaDef}
Equipotencializar é reduzir diferenças perigosas de potencial interligando, de forma planejada, massas e elementos condutivos relevantes ao sistema de equipotencialização.
\end{caixaDef}
\begin{itemize}
\item BEP na origem apropriada.
\item PE e eletrodo(s) de aterramento.
\item Partes metálicas relevantes e serviços que ingressem na edificação.
\item Ligações suplementares onde necessárias.
\end{itemize}
\column{0.49\textwidth}
\centering
\begin{tikzpicture}[scale=0.78,every node/.style={font=\scriptsize}]
\node[draw,fill=corDef!8] (bep) at (0,2) {BEP};
\node[draw] (qd) at (4,2) {QD};
\node[draw] (tub) at (0,0) {Tubulação};
\node[draw] (maq) at (4,0) {Massa};
\node[draw] (est) at (2,-1.2) {Estrutura};
\foreach \x in {qd,tub,maq,est}{\draw[thick,corNorma] (bep)--(\x);}
\draw[thick,corNorma] (bep)--(-1,1) node[ground]{};
\end{tikzpicture}
\end{columns}
\end{frame}

%=====================================================================

%=====================================================================
% PROJETO PREDIAL - ENTRADA, MEDICAO E INTERFACE COM A DISTRIBUIDORA
%=====================================================================

\section{Entrada, medição e padrão Equatorial Maranhão}

\begin{frame}{A fronteira da distribuidora precisa estar explícita}
\centering
\begin{tikzpicture}[
  b/.style={draw=corPrim,fill=corDef!7,rounded corners,minimum width=2.05cm,minimum height=0.72cm,align=center,font=\scriptsize},
  a/.style={-{Stealth},thick,corPrim}]
\node[b](rede)at(-5.0,0){Rede de\\distribuição};
\node[b](pc)at(-2.5,0){Ponto de\\conexão};
\node[b](ramal)at(0,0){Ramal de\\conexão};
\node[b](med)at(2.5,0){Medição e\\proteção geral};
\node[b](qg)at(5.0,0){Origem da instalação\\e QD/QGBT};
\draw[a](rede)--(pc);\draw[a](pc)--(ramal);\draw[a](ramal)--(med);\draw[a](med)--(qg);
\draw[corAccent,very thick,dashed] (1.25,-1.0)--(1.25,1.0);
\node[font=\tiny,corAccent,align=center] at (-1.0,-1.25){critérios da distribuidora};
\node[font=\tiny,corAccent,align=center] at (3.5,-1.25){instalação do consumidor};
\end{tikzpicture}
\vspace{5mm}
\begin{caixaObs}
Responsabilidade física, propriedade dos componentes e limite elétrico não devem ser presumidos. Confirmar a definição na norma vigente e no orçamento de conexão.
\end{caixaObs}
\end{frame}

\begin{frame}[shrink=4]{Documentos vigentes usados no Maranhão}
\scriptsize
\begin{tabularx}{\textwidth}{>{\bfseries}p{2.8cm}p{2.0cm}X}
\toprule
Documento & Revisão consultada & Aplicação \\
\midrule
NT.00001.EQTL & Rev. 09, 29/05/2025 & unidade individual em BT, carga instalada até \SI{75}{kW} \\
NT.00004.EQTL & Rev. 08, 31/12/2025 & empreendimentos de múltiplas unidades consumidoras \\
Anexo I da NT.00004 & Rev. 08, 21/05/2026 & cálculo de queda de tensão em EMUC \\
NT.00020.EQTL & versão vigente & micro e minigeração distribuída \\
NT.00030.EQTL & versão vigente & caixas de medição e proteção \\
NT.00042.EQTL & versão vigente & conexão de estação de recarga de veículo elétrico \\
\bottomrule
\end{tabularx}
\vspace{2mm}
\begin{caixaAt}
Revisões podem mudar durante o projeto. Registrar a data da consulta e conferir novamente antes de protocolar, comprar a caixa e executar o padrão.
\end{caixaAt}
\end{frame}

\begin{frame}[shrink=4]{Tabela 220/380 V: enquadramento inicial da unidade individual}
\scriptsize
\begin{tabularx}{\textwidth}{>{\bfseries}p{2.3cm}p{2.2cm}c c X}
\toprule
Carga instalada & Fornecimento & DG & Cu cliente & Observação \\
\midrule
até \SI{5}{kW} & monofásico 220 V & 25 A & \SI{4}{mm^2} & faixa inicial \\
5,1--\SI{8}{kW} & monofásico 220 V & 40 A & \SI{6}{mm^2} & -- \\
8,1--\SI{12}{kW} & monofásico 220 V & 63 A & \SI{10}{mm^2} & limite monofásico urbano no MA \\
até \SI{12}{kW} & trifásico 380/220 V & 20 A & \SI{6}{mm^2} & opção pode gerar custo de adaptação \\
12,1--\SI{24}{kW} & trifásico 380/220 V & 40 A & \SI{6}{mm^2} & -- \\
24,1--\SI{38}{kW} & trifásico 380/220 V & 63 A & \SI{10}{mm^2} & faixa do projeto residencial I \\
38,1--\SI{48}{kW} & trifásico 380/220 V & 80 A & \SI{16}{mm^2} & -- \\
48,1--\SI{60}{kW} & trifásico 380/220 V & 100 A & \SI{25}{mm^2} & -- \\
60,1--\SI{75}{kW} & trifásico 380/220 V & 125 A & \SI{35}{mm^2} & limite da tabela BT \\
\bottomrule
\end{tabularx}
\vspace{1mm}
\tiny Fonte técnica: Tabela 1 da NT.00001.EQTL Rev.09. Reproduzidos apenas os campos necessários ao enquadramento didático; consultar a tabela integral para ramal, eletrodutos, aterramento e notas.
\end{frame}

\begin{frame}{Exemplo: residência com \SI{33.96}{kW} instalados}
\small
\begin{columns}[T,onlytextwidth]
\column{0.48\textwidth}
\begin{align*}
P_{\mathrm{inst}}&=\SI{33.96}{kW}\\
24{,}1&<33{,}96\leq38
\end{align*}
\begin{caixaEx}[title=Enquadramento]
\begin{itemize}
\item fornecimento: 380/220 V, 3F+N;
\item disjuntor geral: 63 A;
\item condutor Cu do cliente: mínimo \SI{10}{mm^2};
\item verificar demais colunas e notas da tabela.
\end{itemize}
\end{caixaEx}
\column{0.50\textwidth}
\begin{caixaAt}[title=Não encerrar o projeto aqui]
A tabela da distribuidora dimensiona o padrão sob hipóteses próprias. O alimentador interno ainda deve ser recalculado por:
\begin{itemize}
\item capacidade de condução;
\item queda de tensão;
\item curto-circuito e proteção;
\item neutro, PE e método real.
\end{itemize}
\end{caixaAt}
\end{columns}
\end{frame}

\begin{frame}{Padrão aéreo: sequência física dos componentes}
\centering
\begin{tikzpicture}[scale=0.88,every node/.style={font=\scriptsize}]
\draw[very thick] (-5.2,-1.7)--(-5.2,3.3);
\draw[very thick] (-5.6,2.8)--(-4.8,2.8);
\draw[corPrim,thick] (-5.2,2.8).. controls (-2.5,2.5) and (-1.6,2.2)..(0,2.0);
\draw[very thick] (0,-1.7)--(0,3.0);
\node[draw=corPrim,fill=corDef!7,minimum width=1.8cm,minimum height=1.2cm,align=center](cx)at(0.95,0.75){caixa de\\medição};
\node[draw=corPrim,fill=corDef!7,minimum width=1.8cm,minimum height=0.8cm,align=center](dg)at(0.95,-0.45){proteção\\geral};
\draw[corPrim,thick] (0,2.0)--(cx.north);
\draw[corPrim,thick] (cx)--(dg);
\draw[corPrim,thick,-{Stealth}] (dg.east)--(3.15,-0.45)node[right]{alimentador};
\draw[corNorma,thick] (cx.south west)--(-0.65,-0.7)--(-0.65,-1.35)node[ground]{};
\node[align=center] at (-5.2,-2.1){poste da rede};
\node[align=center] at (0,-2.1){poste/muro do padrão};
\node[corAccent,align=center] at (-2.5,3.0){ramal de conexão};
\end{tikzpicture}
\begin{caixaObs}
Alturas, afastamentos, tipo de poste, eletroduto, caixa e ramal devem seguir os desenhos construtivos vigentes da Equatorial e a situação real do imóvel.
\end{caixaObs}
\end{frame}

\begin{frame}{Medição e proteção: acessibilidade sem acesso indevido}
\small
\begin{columns}[T,onlytextwidth]
\column{0.48\textwidth}
\begin{tikzpicture}[scale=0.83]
\draw[very thick,rounded corners] (0,0) rectangle (4.5,5.1);
\draw[thick] (0,2.0)--(4.5,2.0);
\node[draw=corPrim,circle,minimum size=1.35cm] at (2.25,3.55){kWh};
\node[draw=corPrim,fill=corDef!7,minimum width=1.7cm,minimum height=0.7cm] at (2.25,1.0){DG};
\node[font=\scriptsize] at (2.25,4.65){compartimento de medição};
\node[font=\scriptsize] at (2.25,1.65){proteção};
\draw[corAccent,thick,dashed] (0.25,2.25) rectangle (4.25,4.85);
\node[font=\tiny,corAccent] at (3.65,2.45){lacrável};
\end{tikzpicture}
\column{0.50\textwidth}
\begin{itemize}
\item caixa homologada para o tipo de fornecimento;
\item leitura e operação conforme acesso exigido;
\item compartimentos e pontos de lacre preservados;
\item condutores identificados e sem emendas indevidas;
\item raio de curvatura e espaço para terminais;
\item DG compatível com a tabela e capacidade de curto;
\item vedação, drenagem e resistência ao ambiente.
\end{itemize}
\end{columns}
\end{frame}

\begin{frame}{Aterramento do padrão não cria um PE isolado}
\centering
\begin{tikzpicture}[
  b/.style={draw=corPrim,fill=corDef!7,rounded corners,minimum width=2.25cm,minimum height=0.72cm,align=center,font=\scriptsize}]
\node[b](med)at(-4.2,1.2){Medição e DG};
\node[b](qd)at(-1.0,1.2){QD\\barras N e PE};
\node[b,draw=corNorma,fill=corNorma!7](bep)at(2.2,1.2){BEP\\equipotencialização};
\node[b,draw=corNorma,fill=corNorma!7](massas)at(5.2,1.2){Massas e partes\\condutivas};
\draw[-{Stealth},thick,corPrim](med)--(qd);
\draw[thick,corNorma](qd)--(bep)--(massas);
\draw[thick,corNorma](bep.south)--(2.2,-0.4)node[ground]{};
\node[font=\tiny,corNorma] at (3.0,-0.45){eletrodo(s)};
\end{tikzpicture}
\vspace{4mm}
\begin{caixaAt}
O condutor de proteção dos circuitos retorna ao barramento PE e ao BEP. Neutro e PE só se relacionam onde o esquema de aterramento determinar; não instalar pontes improvisadas nos quadros terminais.
\end{caixaAt}
\end{frame}

\begin{frame}{Ambiente costeiro altera materiais e manutenção}
\small
\begin{columns}[T,onlytextwidth]
\column{0.49\textwidth}
\begin{caixaNorma}[title=Até e após \SI{2}{km} da orla]
A NT.00001 diferencia materiais do ramal conforme a distância da orla marítima. O projeto deve registrar a localização e consultar a coluna correta da tabela.
\end{caixaNorma}
\column{0.49\textwidth}
\begin{itemize}
\item selecionar caixa, ferragens e fixações adequadas;
\item evitar pares galvânicos incompatíveis;
\item proteger entradas contra água e condensação;
\item aplicar torque e preparação conforme fabricante;
\item prever inspeção de corrosão e reaperto técnico.
\end{itemize}
\end{columns}
\begin{caixaAt}
Pintura ou fita não corrige material inadequado. A solução começa na especificação e segue na montagem e manutenção.
\end{caixaAt}
\end{frame}

\begin{frame}{EMUC: a unidade não define sozinha o empreendimento}
\centering
\begin{tikzpicture}[
  b/.style={draw=corPrim,fill=corDef!7,rounded corners,minimum width=2.0cm,minimum height=0.7cm,align=center,font=\scriptsize},
  a/.style={-{Stealth},thick,corPrim}]
\node[b](rede)at(0,3.4){Rede};
\node[b](med)at(0,2.1){Centro de medição};
\node[b](pr)at(0,0.8){Prumada / distribuição};
\node[b](u1)at(-3.6,-0.8){Unidade 1};
\node[b](u2)at(-1.2,-0.8){Unidade 2};
\node[b](u3)at(1.2,-0.8){Unidade 3};
\node[b](uc)at(3.6,-0.8){Área comum};
\draw[a](rede)--(med);\draw[a](med)--(pr);
\foreach \x in {u1,u2,u3,uc}{\draw[a](pr)--(\x);}
\end{tikzpicture}
\vspace{2mm}
\begin{caixaNorma}
Demanda do conjunto, centro de medição, prumadas, queda, proteção geral, espaço técnico e transferência de ativos seguem a NT.00004 e seus anexos vigentes.
\end{caixaNorma}
\end{frame}

\begin{frame}{Vistoria: conferir antes de solicitar a ligação}
\small
\begin{enumerate}
\item tipo de fornecimento, caixa e DG conferem com a carga declarada;
\item poste, ramal, eletrodutos, curvas e fixações seguem o desenho aplicável;
\item alturas, afastamentos e acesso ao medidor estão corretos;
\item condutores, terminais, identificação e raios de curvatura estão adequados;
\item aterramento, eletrodo e conexão estão instalados e acessíveis;
\item compartimentos lacráveis não foram invadidos;
\item padrão está íntegro, vedado, limpo e sem materiais improvisados;
\item dados do imóvel, carga e titular estão coerentes com a solicitação.
\end{enumerate}
\begin{caixaAt}
Não energizar a instalação interna como teste improvisado. Ensaios e energização seguem procedimento, instrumentos adequados e responsabilidades definidas.
\end{caixaAt}
\end{frame}

\begin{frame}{Erros que causam retrabalho no padrão}
\scriptsize
\begin{tabularx}{\textwidth}{>{\bfseries}p{3.4cm}X}
\toprule
Erro & Correção de projeto/obra \\
\midrule
Caixa comprada antes do enquadramento & fechar carga e tipo de fornecimento primeiro \\
DG fora da faixa normativa & usar tabela vigente e justificar caso especial \\
Poste ou eletroduto incompatível & seguir desenho construtivo aplicável \\
Medidor sem acesso regulamentar & reposicionar o padrão antes da execução final \\
Emenda em compartimento inadequado & prever comprimento e terminal correto \\
Neutro e PE unidos indevidamente & corrigir barras e esquema de aterramento \\
Material inadequado à orla & especificar conforme ambiente e tabela \\
Projeto interno divergente da carga declarada & revisar quadro, memorial e solicitação \\
\bottomrule
\end{tabularx}
\end{frame}

\section{SPDA e proteção contra descargas atmosféricas}

\begin{frame}{NBR 5419: atualização confirmada em 2026}
\small
\begin{columns}[T,onlytextwidth]
\column{0.52\textwidth}
\begin{caixaNorma}[title=ABNT NBR 5419-1:2026]
A edição 2026 da \textbf{Parte 1} está publicada e estabelece os requisitos para a determinação da necessidade de proteção contra descargas atmosféricas.
\end{caixaNorma}
\begin{itemize}
\item não assumir automaticamente que todas as partes da série mudaram de edição na mesma data;
\item no projeto, registrar a edição efetivamente adotada de cada parte aplicável;
\item verificar requisitos do órgão licenciador, Corpo de Bombeiros e contratante.
\end{itemize}
\column{0.46\textwidth}
\begin{caixaAt}[title=Regra para este material]
A referência ``NBR 5419:2015'' não é mais usada de forma genérica. Quando a Parte 1 for citada, o material indica explicitamente \textbf{ABNT NBR 5419-1:2026}; para as demais partes, exige-se conferência da edição vigente antes do projeto executivo.
\end{caixaAt}
\end{columns}
\end{frame}

\begin{frame}{PDA não é apenas o captor no telhado}
\small
\centering
\begin{tikzpicture}[scale=0.82,every node/.style={font=\scriptsize}]
\tikzset{b/.style={draw=corPrim,fill=corDef!7,rounded corners,minimum width=3.0cm,minimum height=0.7cm,align=center}}
\node[b] (r) at (0,2.6) {Avaliação da necessidade\\e análise de risco};
\node[b] (s) at (-3.4,1.1) {SPDA externo\\captação/descida/terra};
\node[b] (e) at (0,1.1) {Equipotencialização\\para descargas};
\node[b] (m) at (3.4,1.1) {MPS\\DPS, roteamento, blindagem};
\node[b] (i) at (0,-0.5) {Sistemas internos\\energia, dados e automação};
\draw[-{Stealth},thick,corPrim] (r)--(s);\draw[-{Stealth},thick,corPrim] (r)--(e);\draw[-{Stealth},thick,corPrim] (r)--(m);
\draw[-{Stealth},thick,corNorma] (s)--(i);\draw[-{Stealth},thick,corNorma] (e)--(i);\draw[-{Stealth},thick,corNorma] (m)--(i);
\end{tikzpicture}
\vspace{2mm}
\begin{caixaObs}
A interface entre SPDA, equipotencialização e DPS precisa ser projetada conjuntamente. Separar esses projetos em ``ilhas'' cria lacunas de proteção.
\end{caixaObs}
\end{frame}

\begin{frame}{Zonas de proteção e roteamento}
\small
\begin{columns}[T,onlytextwidth]
\column{0.49\textwidth}
\begin{itemize}
\item A energia do surto deve ser reduzida progressivamente entre zonas.
\item Entradas metálicas e cabos são pontos críticos de acoplamento.
\item Roteamento, blindagem, equipotencialização e DPS trabalham em conjunto.
\item Cabos sensíveis não devem percorrer rotas inadequadas junto a condutores de descida sem análise.
\end{itemize}
\column{0.49\textwidth}
\centering
\begin{tikzpicture}[scale=0.82,every node/.style={font=\scriptsize}]
\draw[thick,corAt] (0,0) rectangle (5.3,3.4);
\draw[thick,corAccent] (0.7,0.6) rectangle (4.6,2.8);
\draw[thick,corNorma] (1.5,1.2) rectangle (3.8,2.2);
\node[corAt] at (4.65,3.1) {ZPR externa};
\node[corAccent] at (3.9,2.55) {zona interna};
\node[corNorma] at (2.65,1.7) {equipamento sensível};
\draw[-{Stealth},thick,corPrim] (-0.6,1.7)--(1.4,1.7) node[midway,above]{linha};
\node at (0.7,1.35) {DPS};
\end{tikzpicture}
\end{columns}
\end{frame}

%=====================================================================
\section{Veículos elétricos e cargas modernas}

\begin{frame}{Recarga de VE: por que não é uma tomada comum}
\small
\begin{columns}[T,onlytextwidth]
\column{0.50\textwidth}
\begin{itemize}
\item Potência elevada durante longos períodos.
\item Carga eletrônica com requisitos próprios de proteção.
\item Possibilidade de corrente residual CC conforme arquitetura.
\item Expansão simultânea de vários carregadores.
\item Aumento de carga pode alterar o padrão de entrada.
\item Gestão de demanda pode se tornar parte do projeto.
\end{itemize}
\column{0.48\textwidth}
\begin{caixaNorma}[title=Referências]
\begin{itemize}
\item ABNT NBR 17019:2022 — instalação fixa para alimentação/recepção de energia de VE;
\item NT.00042.EQTL Rev.04 — conexão de estações de recarga na área da Equatorial;
\item NT.00001/NT.00004 — entrada individual ou empreendimento com múltiplas UCs, conforme o caso.
\end{itemize}
\end{caixaNorma}
\end{columns}
\end{frame}

\begin{frame}{Arquitetura correta de um ponto de recarga}
\small
\centering
\begin{tikzpicture}[scale=0.78,every node/.style={font=\scriptsize}]
\tikzset{b/.style={draw=corPrim,fill=corDef!7,rounded corners,minimum width=2.25cm,minimum height=0.68cm,align=center}}
\node[b] (q) at (0,2.1) {QD};
\node[b] (p) at (3.0,2.1) {Proteção dedicada\\sobrecorrente + DR};
\node[b] (e) at (6.0,2.1) {SAVE / EVSE};
\node[b] (v) at (9.0,2.1) {Veículo};
\node[b] (pe) at (4.5,0.35) {Barramento PE / BEP\\equipotencialização};
\draw[-{Stealth},thick,corPrim] (q)--(p)--(e)--(v);
\draw[corNorma,thick] (q.south)--(pe.west);
\draw[corNorma,thick] (p.south)--(pe);
\draw[corNorma,thick] (e.south)--(pe.east);
\draw[corNorma,thick] (pe.south)--(4.5,-0.45) node[ground]{};
\end{tikzpicture}
\vspace{2mm}
\begin{caixaAt}
O PE é contínuo e integrado ao sistema de equipotencialização. Não se desenha um eletrodo de terra independente para cada equipamento. O tipo de DR deve ser compatível com a corrente residual possível, a detecção CC incorporada ao SAVE e as instruções normativas/do fabricante.
\end{caixaAt}
\end{frame}

\begin{frame}{Exemplo: SAVE monofásico de 7,4 kW}
\small
\begin{columns}[T,onlytextwidth]
\column{0.49\textwidth}
\begin{align*}
P&=\SI{7.4}{kW},\quad V=\SI{220}{V}\\
I_b&\approx\frac{7400}{220}=\SI{33.6}{A}
\end{align*}
\begin{itemize}
\item circuito dedicado;
\item cabo por ampacidade + queda + método;
\item proteção contra sobrecorrente;
\item proteção diferencial compatível;
\item PE e DPS conforme o sistema.
\end{itemize}
\column{0.49\textwidth}
\begin{caixaProjeto}[title=Impacto no projeto residencial]
No exemplo-base deste material:
\[
33{,}96+7{,}40=\SI{41.36}{kW}.
\]
Logo, a inclusão do SAVE não é apenas ``mais um circuito'': no exemplo, a faixa da Tabela 1 da NT.00001 passa de 24,1--38 kW (63 A) para 38,1--48 kW (80 A). QD, alimentador, balanceamento e solicitação de aumento de carga devem ser revistos.
\end{caixaProjeto}
\end{columns}
\end{frame}

\begin{frame}{Gestão dinâmica de recarga}
\small
\centering
\begin{tikzpicture}[scale=0.80,every node/.style={font=\scriptsize}]
\tikzset{b/.style={draw=corPrim,fill=corDef!7,rounded corners,minimum width=2.25cm,minimum height=0.65cm,align=center}}
\node[b] (med) at (0,2.2) {Medição da demanda\\da instalação};
\node[b] (ctrl) at (3.2,2.2) {Controlador\\de carga};
\node[b] (ev1) at (6.4,3.2) {SAVE 1};
\node[b] (ev2) at (6.4,2.2) {SAVE 2};
\node[b] (ev3) at (6.4,1.2) {SAVE 3};
\node[b] (pred) at (0,0.5) {Demais cargas\\da edificação};
\draw[-{Stealth},thick,corPrim] (med)--(ctrl);
\foreach \x in {ev1,ev2,ev3}{\draw[-{Stealth},thick,corAccent] (ctrl)--(\x);}
\draw[-{Stealth},thick,corPrim] (med)--(pred);
\node[align=center] at (3.2,0.45) {limite disponível\\+ prioridades};
\end{tikzpicture}
\vspace{1mm}
\begin{caixaObs}
Gestão de carga reduz simultaneidade dos carregadores, mas não elimina as verificações de capacidade térmica, proteção, qualidade de energia e requisitos da distribuidora.
\end{caixaObs}
\end{frame}

\begin{frame}{Geração distribuída: interface com a instalação}
\small
\begin{columns}[T,onlytextwidth]
\column{0.49\textwidth}
\begin{itemize}
\item Corrente pode fluir em sentido diferente do projeto tradicional.
\item Proteções e seccionamento existem nos lados CC e CA.
\item DPS e equipotencialização precisam ser compatíveis com o SPDA.
\item Verificar capacidade de barramentos e condutores no ponto de conexão.
\item O estudo interno deve fechar com a solicitação de conexão à distribuidora.
\end{itemize}
\column{0.49\textwidth}
\begin{caixaNorma}[title=Equatorial]
A conexão de micro e minigeração distribuída é tratada pela \textbf{NT.00020.EQTL Rev.06}. Para unidades do Grupo B, usar também os formulários e anexos atuais disponibilizados pela distribuidora.
\end{caixaNorma}
\begin{caixaAt}
Geração distribuída não autoriza reduzir proteção, seção de condutor ou capacidade de barramento apenas porque parte da energia é produzida localmente.
\end{caixaAt}
\end{columns}
\end{frame}

\begin{frame}{UPS, geradores e transferência}
\small
\begin{columns}[T,onlytextwidth]
\column{0.49\textwidth}
\begin{itemize}
\item Definir cargas essenciais e não essenciais.
\item Analisar curto-circuito em modo rede e modo fonte alternativa.
\item A proteção que atua bem na rede pode não atuar igual com fonte limitada eletronicamente.
\item Verificar neutro, aterramento e chaveamento conforme a topologia.
\end{itemize}
\column{0.49\textwidth}
\centering
\begin{tikzpicture}[scale=0.78,every node/.style={font=\scriptsize}]
\tikzset{b/.style={draw=corPrim,fill=corDef!7,rounded corners,minimum width=2.1cm,minimum height=0.62cm,align=center}}
\node[b] (r) at (0,2.6) {Rede};
\node[b] (g) at (0,0.8) {Gerador};
\node[b] (ats) at (3.2,1.7) {ATS};
\node[b] (ups) at (6.0,1.7) {UPS};
\node[b] (crit) at (8.8,1.7) {Cargas críticas};
\draw[-{Stealth},thick,corPrim] (r)--(ats);\draw[-{Stealth},thick,corAccent] (g)--(ats);\draw[-{Stealth},thick,corPrim] (ats)--(ups)--(crit);
\end{tikzpicture}
\end{columns}
\end{frame}

%=====================================================================

\section{Qualidade de energia e eficiência}

\begin{frame}{Harmônicas: medida básica}
\small
\begin{columns}[T,onlytextwidth]
\column{0.50\textwidth}
\[
THD_I=100\frac{\sqrt{I_2^2+I_3^2+\cdots+I_n^2}}{I_1}.
\]
\begin{itemize}
\item THD não informa sozinho qual harmônica domina.
\item Analisar espectro e forma de onda.
\item Fonte, carga e impedância da rede determinam a distorção de tensão.
\end{itemize}
\column{0.48\textwidth}
\begin{tikzpicture}
\begin{axis}[width=5.8cm,height=3.8cm,ybar,bar width=10pt,xlabel={ordem},ylabel={$I_h/I_1$},ymin=0,ymax=1.1,xtick={1,3,5,7,9,11},grid=major]
\addplot coordinates {(1,1)(3,0.42)(5,0.22)(7,0.16)(9,0.09)(11,0.07)};
\end{axis}
\end{tikzpicture}
\end{columns}
\end{frame}

\begin{frame}{Efeitos de harmônicas}
\small
\begin{columns}[T,onlytextwidth]
\column{0.49\textwidth}
\begin{itemize}
\item Aquecimento extra de cabos e transformadores.
\item Sobrecarga do neutro por harmônicas triplas.
\item Perdas adicionais em motores.
\item Erros/interferências em equipamentos sensíveis.
\end{itemize}
\column{0.49\textwidth}
\begin{itemize}
\item Ressonância com bancos de capacitores.
\item Correntes elevadas em capacitores.
\item Atuação indevida de proteção.
\item Redução de vida útil de componentes.
\end{itemize}
\end{columns}
\begin{caixaProjeto}[title=Mitigação]
Reatores, filtros passivos/ativos, transformadores adequados, arquitetura de alimentação, seleção de conversores e projeto do banco de capacitores -- sempre após medição/estudo.
\end{caixaProjeto}
\end{frame}

\begin{frame}{Correção de fator de potência}
\small
\begin{columns}[T,onlytextwidth]
\column{0.51\textwidth}
\begin{align*}
Q_c=P\left(\tg\varphi_1-\tg\varphi_2\right)
\end{align*}
\begin{itemize}
\item Reduz corrente reativa circulante a montante do ponto de compensação.
\item Pode liberar capacidade e reduzir perdas.
\item Não corrige distorção harmônica.
\end{itemize}
\column{0.47\textwidth}
\begin{caixaEx}
$P=\SI{100}{kW}$, $fp_1=0{,}78$, $fp_2=0{,}95$:
\[
Q_c\approx\SI{47}{kvar}.
\]
Em ambiente com VFD, verificar ressonância e considerar banco dessintonizado/filtro conforme estudo.
\end{caixaEx}
\end{columns}
\end{frame}

\begin{frame}{Gestão de demanda}
\small
\centering
\begin{tikzpicture}
\begin{axis}[width=10.6cm,height=4.5cm,xlabel={hora},ylabel={kW},xmin=0,xmax=24,ymin=0,ymax=110,grid=both,legend pos=north west]
\addplot[corPrim,thick,smooth] coordinates {(0,25)(4,20)(8,65)(10,80)(12,72)(14,88)(17,100)(18,96)(20,62)(24,32)};\addlegendentry{sem gestão}
\addplot[corAccent,thick,dashed,smooth] coordinates {(0,25)(4,20)(8,62)(10,72)(12,69)(14,75)(17,78)(18,76)(20,60)(24,32)};\addlegendentry{com gestão}
\end{axis}
\end{tikzpicture}
\begin{caixaObs}
Deslocar cargas flexíveis, limitar recarga de VE e coordenar climatização pode reduzir pico sem reduzir energia útil do processo.
\end{caixaObs}
\end{frame}

\begin{frame}{Medição: o que registrar}
\scriptsize
\begin{tabularx}{\textwidth}{>{\bfseries}p{4.0cm}X}
\toprule
Grandeza & Uso \\
\midrule
$V$, $I$, $P$, $Q$, $S$, fp & Estado operacional e balanceamento. \\
Demanda e energia & Custos, contratos e eficiência. \\
THD/espectro & Diagnóstico de cargas não lineares. \\
Afundamentos e elevações & Falhas de processo e qualidade de suprimento. \\
Interrupções & Continuidade e confiabilidade. \\
Temperatura e termografia & Conexões, sobrecarga e ventilação. \\
Eventos de proteção & Causa-raiz e seletividade. \\
\bottomrule
\end{tabularx}
\end{frame}

%=====================================================================

\section{CAD, BIM, simulação e documentação}

\begin{frame}{Ferramentas: cada uma responde a uma pergunta}
\scriptsize
\begin{tabularx}{\textwidth}{>{\bfseries}p{2.6cm}X}
\toprule
Ferramenta & Uso \\
\midrule
CAD elétrico & Plantas, diagramas, detalhes, listas e documentação 2D. \\
BIM & Compatibilização espacial, parâmetros, quantitativos e coordenação multidisciplinar. \\
QElectroTech & Diagramas didáticos/funcionais de comando e potência. \\
ETAP / EasyPower / similares & Fluxo de carga, curto, seletividade, partida, arco e harmônicas conforme módulos. \\
Planilha / Python & Automação de cálculos, auditoria, geração de tabelas e checagens repetitivas. \\
Analisador de energia & Validar modelo com medições reais em retrofit/operação. \\
\bottomrule
\end{tabularx}
\end{frame}

\begin{frame}{Fluxo de simulação elétrica}
\small
\centering
\begin{tikzpicture}[node distance=0.30cm,every node/.style={font=\scriptsize}]
\tikzset{b/.style={draw=corPrim,fill=corDef!7,rounded corners,minimum width=2.35cm,minimum height=0.65cm,align=center}}
\node[b] (a) {Topologia e fontes};
\node[b,right=of a] (b) {Cabos / trafos};
\node[b,right=of b] (c) {Cargas / modos};
\node[b,below=of c] (d) {Fluxo de carga};
\node[b,left=of d] (e) {Curto-circuito};
\node[b,left=of e] (f) {Seletividade / arco};
\node[b,below=of f] (g) {Partida / harmônicas};
\node[b,right=of g] (h) {Relatórios};
\node[b,right=of h] (i) {Revisão do projeto};
\foreach \u/\v in {a/b,b/c,c/d,d/e,e/f,f/g,g/h,h/i}{\draw[-{Stealth},thick,corPrim] (\u)--(\v);}
\draw[-{Stealth},thick,corAccent] (i.east) to[out=0,in=-80] (a.south east);
\end{tikzpicture}
\vspace{2mm}
\begin{caixaAt}
Software calcula o modelo fornecido. Dado errado + modelo errado = resultado preciso e inútil.
\end{caixaAt}
\end{frame}

\begin{frame}{BIM: parâmetros elétricos úteis}
\small
\begin{columns}[T,onlytextwidth]
\column{0.49\textwidth}
\begin{itemize}
\item Circuito e quadro de origem.
\item Potência e tensão.
\item Número de fases.
\item Fator de potência e demanda.
\item Reserva e status do circuito.
\end{itemize}
\column{0.49\textwidth}
\begin{itemize}
\item Comprimento e rota.
\item Cabo/eletroduto/leito.
\item Dispositivo de proteção.
\item Classificação de carga crítica.
\item Código para operação/manutenção.
\end{itemize}
\end{columns}
\begin{caixaObs}
O valor do BIM aumenta quando os parâmetros alimentam cálculo, documentação e operação -- não apenas um modelo 3D bonito.
\end{caixaObs}
\end{frame}

\begin{frame}[fragile]{Automação de checagens com Python}
\scriptsize
\begin{lstlisting}[style=codPy]
from math import sqrt, acos, sin

def corrente_3f(P, Vll, fp=1.0, eta=1.0):
    return P/(sqrt(3)*Vll*fp*eta)

def queda_3f(I, L_m, R_ohm_km, X_ohm_km, fp, Vll):
    phi = acos(fp)
    dv = sqrt(3)*I*(L_m/1000)*(R_ohm_km*fp + X_ohm_km*sin(phi))
    return 100*dv/Vll

I = corrente_3f(42000, 380, 0.88)
dv = queda_3f(I, 60, 0.75, 0.08, 0.88, 380)
print(I, dv)
\end{lstlisting}
\begin{caixaObs}
O script deve ler dados rastreáveis e gerar alerta, não escolher cabo ``por mágica''. O engenheiro continua responsável pelas premissas e validação.
\end{caixaObs}
\end{frame}

\begin{frame}{Documentação para manutenção}
\small
\begin{columns}[T,onlytextwidth]
\column{0.49\textwidth}
\begin{itemize}
\item Unifilar atualizado.
\item Ajustes de relés/disjuntores.
\item Curto-circuito disponível por barra.
\item Identificação de fontes alternativas.
\item Procedimentos de manobra.
\end{itemize}
\column{0.49\textwidth}
\begin{itemize}
\item Etiquetas e riscos aplicáveis.
\item Histórico de intervenções.
\item Ensaios de DR, isolamento e aterramento.
\item Termografias e inspeções.
\item Lista de sobressalentes críticos.
\end{itemize}
\end{columns}
\end{frame}

%=====================================================================


% Bases gráficas para os projetos completos.
\input{projetos/base/plantas_arquitetonicas.tex}
% Portas e vaos complementares para as plantas realistas
% Implementacao simples e robusta, sem aritmetica adicional nas coordenadas TikZ.
\newcommand{\ArqVaoH}[3]{%
  \draw[white,line width=5.2pt] (#1,#3)--(#2,#3);
  \draw[black!20,thin] (#1,#3)--(#2,#3);}
\newcommand{\ArqVaoV}[3]{%
  \draw[white,line width=5.2pt] (#3,#1)--(#3,#2);
  \draw[black!20,thin] (#3,#1)--(#3,#2);}
\newcommand{\ArqPortaCorrerH}[3]{%
  \draw[white,line width=5.4pt] (#1,#3)--(#2,#3);
  \draw[black!55,thin,double=white,double distance=1.35pt] (#1,#3)--(#2,#3);}
\newcommand{\ArqPortaCorrerV}[3]{%
  \draw[white,line width=5.4pt] (#3,#1)--(#3,#2);
  \draw[black!55,thin,double=white,double distance=1.35pt] (#3,#1)--(#3,#2);}
\newcommand{\ArqPortaElevV}[3]{%
  \draw[white,line width=5.0pt] (#3,#1)--(#3,#2);
  \draw[black!65,thin,double=white,double distance=1.15pt] (#3,#1)--(#3,#2);}


%=====================================================================
% PROJETO RESIDENCIAL I — CASA TERREA COM DUAS SUITES
% Base eletrica: ABNT NBR 5410 + NT.00001.EQTL Rev.09
% Planta-base didatica com logica arquitetonica e mobiliario.
% Atualizacao: agosto de 2026
%=====================================================================

\ifdefined\ArqParedeExt\else\input{projetos/base/plantas_arquitetonicas.tex}\fi

\section{Projeto Residencial I — casa térrea com duas suítes}

% Símbolos elétricos locais
\providecommand{\PRLuz}[2]{\begin{scope}[shift={(#1,#2)}]\draw[corAccent,thick] (0,0) circle (0.14);\draw[corAccent,thick] (-0.10,-0.10)--(0.10,0.10);\draw[corAccent,thick] (-0.10,0.10)--(0.10,-0.10);\end{scope}}
\providecommand{\PRTUG}[2]{\begin{scope}[shift={(#1,#2)}]\draw[corPrim,thick] (0,0) circle (0.13);\draw[corPrim,thick] (-0.16,0)--(0.16,0);\end{scope}}
\providecommand{\PRTUE}[3]{\begin{scope}[shift={(#1,#2)}]\draw[corAt,thick] (0,0) circle (0.17);\node[font=\tiny\bfseries,corAt] at (0,0){#3};\end{scope}}
\providecommand{\PRQD}[3]{\begin{scope}[shift={(#1,#2)}]\draw[corPrim,fill=corDef!10,thick,rounded corners=1pt] (-0.38,-0.22) rectangle (0.38,0.22);\node[font=\tiny\bfseries] at (0,0){#3};\end{scope}}

%---------------------------------------------------------------------
% Planta 15,00 x 11,00 m. Frente em y=0.
% Ala íntima à esquerda; social ao centro; garagem/serviço à direita.
%---------------------------------------------------------------------
\newcommand{\CasaResidencialBase}{%
  % contorno
  \draw[arqExt] (0,0) rectangle (15,11);
  % ala íntima e hall
  \ArqParedeInt{5.2}{0}{5.2}{11}
  \ArqParedeInt{6.4}{0}{6.4}{11}
  \ArqParedeInt{0}{3.6}{5.2}{3.6}
  \ArqParedeInt{0}{7.2}{5.2}{7.2}
  % banheiros internos / social
  \ArqParedeInt{3.5}{0}{3.5}{2.0}
  \ArqParedeInt{3.5}{2.0}{5.2}{2.0}
  \ArqParedeInt{3.5}{4.8}{5.2}{4.8}
  \ArqParedeInt{3.5}{4.8}{3.5}{7.2}
  \ArqParedeInt{3.5}{8.4}{5.2}{8.4}
  \ArqParedeInt{3.5}{8.4}{3.5}{11}
  % social/cozinha
  \ArqParedeInt{6.4}{6.4}{11.2}{6.4}
  % garagem / serviço / varanda gourmet
  \ArqParedeInt{11.2}{0}{11.2}{11}
  \ArqParedeInt{11.2}{5.6}{15}{5.6}
  \ArqParedeInt{11.2}{8.2}{15}{8.2}

  % acessos: principal diretamente na sala; garagem para sala; serviço externo
  \ArqPortaHUp{7.25}{0}{1.10}
  \draw[white,line width=5pt] (11.2,2.2)--(11.2,3.2);\draw[black!60,thin] (11.2,2.2)--(10.2,2.2);
  \ArqPortaVLeft{15}{6.25}{0.90}
  \draw[white,line width=5pt] (11.2,9.1)--(11.2,10.7);\draw[arqGlass] (11.2,9.1)--(11.2,10.7);
  % hall interno aberto para a sala
  \draw[white,line width=5pt] (6.4,2.45)--(6.4,3.65);
  % portas dos quartos/banheiros para hall
  \ArqPortaVLeft{5.2}{2.45}{0.90}
  \ArqPortaVLeft{5.2}{3.85}{0.90}
  \ArqPortaVLeft{5.2}{7.45}{0.90}
  \ArqPortaVLeft{5.2}{0.55}{0.80}
  % banheiros das suítes: acesso exclusivo por dentro do quarto
  \ArqPortaVLeft{3.5}{5.05}{0.80}
  \ArqPortaVLeft{3.5}{8.65}{0.80}
  % cozinha -> serviço
  \ArqPortaVRight{11.2}{6.75}{0.85}
  % vão de garagem
  \draw[white,line width=5.5pt] (11.55,0)--(14.70,0);\draw[black!30,very thin] (11.55,0.10)--(14.70,0.10);

  % janelas
  \ArqJanV{0.45}{2.75}{0}
  \ArqJanV{4.15}{6.55}{0}
  \ArqJanV{7.75}{10.20}{0}
  \ArqJanH{8.0}{10.2}{0}
  \ArqJanH{7.3}{9.8}{11}
  \ArqJanH{0.7}{2.6}{11}
  \ArqJanH{12.0}{14.1}{11}

  % mobiliário - quarto 3
  \ArqCama{0.55}{0.45}{2.15}{2.65}
  \ArqArmario{2.85}{2.45}{0.55}{0.95}
  \node[arqLabel] at (1.75,3.25){QUARTO 3};
  % banho social
  \ArqChuveiro{3.65}{0.15}{1.25}{0.85}\ArqVaso{4.55}{1.35}\ArqPia{3.62}{1.20}{0.65}{0.45}
  \node[arqSub] at (4.30,1.82){BANHO SOCIAL};

  % suíte 2
  \ArqCama{0.55}{4.05}{2.20}{2.65}\ArqArmario{2.85}{3.95}{0.50}{1.05}
  \ArqChuveiro{3.65}{6.05}{1.25}{0.95}\ArqVaso{4.55}{5.35}\ArqPia{3.62}{5.05}{0.65}{0.45}
  \node[arqLabel] at (1.75,6.85){SUÍTE 2};\node[arqSub] at (4.35,6.90){B2};

  % suíte master
  \ArqCama{0.55}{7.75}{2.35}{2.70}\ArqArmario{2.95}{7.45}{0.45}{0.85}
  \ArqChuveiro{3.65}{10.00}{1.25}{0.85}\ArqVaso{4.55}{9.20}\ArqPia{3.62}{8.65}{0.65}{0.45}
  \node[arqLabel] at (1.85,10.65){SUÍTE MASTER};\node[arqSub] at (4.35,10.70){B1};
  \node[arqSub,rotate=90] at (5.80,5.50){HALL ÍNTIMO};

  % sala/jantar
  \ArqSofa{7.00}{0.85}{3.10}{1.05}\ArqBancada{10.50}{0.95}{0.35}{2.25}
  \ArqMesa{7.25}{4.20}{2.55}{1.10}
  \node[arqLabel] at (8.80,3.20){SALA / JANTAR};
  % cozinha
  \ArqBancada{6.75}{10.15}{3.50}{0.55}\ArqBancada{10.35}{7.05}{0.55}{3.65}
  \ArqFogao{9.25}{10.05}\ArqGeladeira{10.15}{6.65}\ArqBancada{7.40}{7.55}{2.30}{0.80}
  \node[arqLabel] at (8.60,9.45){COZINHA};
  % garagem
  \ArqCarro{12.00}{0.65}{2.15}{4.25}\node[arqLabel] at (13.10,5.20){GARAGEM};
  % serviço
  \ArqLavadora{12.00}{6.05}\ArqBancada{13.0}{6.0}{1.45}{0.55}\node[arqLabel] at (13.10,7.65){SERVIÇO};
  % gourmet / quintal coberto
  \ArqBancada{12.0}{9.85}{2.35}{0.55}\ArqMesa{12.20}{8.70}{1.55}{0.75}\node[arqLabel] at (13.10,10.75){VARANDA GOURMET};

  \CotaH{0}{15}{11.18}{15,00 m}\CotaV{0}{11}{15.18}{11,00 m}
}

\begin{frame}{Residencial I — planta arquitetônica funcional}
\centering
\begin{tikzpicture}[scale=0.49,every node/.style={font=\tiny}]
\CasaResidencialBase
\end{tikzpicture}
\vspace{0.5mm}
\scriptsize A entrada principal chega à sala. O hall é interno. A casa possui três banheiros completos: banho social, B2 e B1; estes dois últimos são acessados exclusivamente pelas suítes.
\end{frame}

\begin{frame}{Residencial I — iluminação e circuitos C1/C2}
\centering
\begin{tikzpicture}[scale=0.49,every node/.style={font=\tiny}]
\CasaResidencialBase
\PRQD{6.10}{3.10}{QD1}
\foreach \x/\y in {1.7/2.1,4.25/1.0,1.75/5.4,4.25/5.9,1.8/9.2,4.25/9.6,5.8/5.5,8.8/2.8,8.7/5.2,8.6/8.8,13.1/3.0,13.1/6.9,13.1/9.6}{\PRLuz{\x}{\y}}
% C1: sala, cozinha, servico e gourmet
\draw[corAccent!80!black,dashed,thick] (6.10,3.10)--(8.8,2.8)--(8.7,5.2)--(8.6,8.8)--(13.1,9.6);
\draw[corAccent!80!black,dashed,thick] (8.6,8.8)--(13.1,6.9);
% C2: ala intima, banheiros, hall e garagem
\draw[corAccent!55!black,dashed,thick] (6.10,3.10)--(5.8,5.5)--(1.75,5.4)--(4.25,5.9)--(1.8,9.2)--(4.25,9.6);
\draw[corAccent!55!black,dashed,thick] (6.10,3.10)--(1.7,2.1)--(4.25,1.0);
\draw[corAccent!55!black,dashed,thick] (6.10,3.10)--(13.1,3.0);
\node[font=\tiny\bfseries,fill=white,inner sep=1pt,corAccent!80!black] at (9.5,7.4){C1};
\node[font=\tiny\bfseries,fill=white,inner sep=1pt,corAccent!55!black] at (3.0,5.4){C2};
\end{tikzpicture}
\vspace{0.4mm}
\scriptsize Previsão calculada: C1 = 1,06 kVA e C2 = 1,40 kVA. Os 13 símbolos são pontos de luz; os 2,46 kVA resultam das áreas e dos acréscimos completos de 4 m$^2$, não da quantidade de símbolos.
\end{frame}

\begin{frame}{Residencial I — 43 TUG e oito cargas específicas}
\centering
\begin{tikzpicture}[scale=0.49,every node/.style={font=\tiny}]
\CasaResidencialBase
\PRQD{6.10}{3.10}{QD1}
% C3: sala e dormitorios — 17 pontos
\foreach \x/\y in {0.25/0.35,0.25/3.30,3.25/0.35,3.25/3.30,0.25/3.85,0.25/6.95,3.25/3.85,3.25/6.95,0.25/7.45,0.25/10.75,3.25/7.45,3.25/10.75,6.65/0.35,10.95/0.35,10.95/3.60,10.95/6.15,6.65/6.15}{\PRTUG{\x}{\y}}
% C4: hall e garagem — 9 pontos
\foreach \x/\y in {6.25/0.70,6.25/2.10,6.25/4.20,6.25/7.80,6.25/10.20,11.35/0.80,11.35/4.80,14.65/1.50,14.65/4.80}{\PRTUG{\x}{\y}}
% C5 cozinha; C6 servico; C7 gourmet; C8 banheiros
\foreach \x/\y in {6.70/10.75,8.00/10.75,9.30/10.75,10.95/7.00,10.95/8.50,10.95/10.00}{\PRTUG{\x}{\y}}
\foreach \x/\y in {11.35/5.85,14.65/5.85,11.35/7.95,14.65/7.95}{\PRTUG{\x}{\y}}
\foreach \x/\y in {11.35/8.45,14.65/8.45,11.35/10.75,14.65/10.75}{\PRTUG{\x}{\y}}
\foreach \x/\y in {5.00/1.80,4.95/5.00,4.95/8.65}{\PRTUG{\x}{\y}}
% TUE C9--C16
\PRTUE{4.35}{0.55}{C9}\PRTUE{4.35}{6.55}{C10}\PRTUE{4.35}{10.45}{C11}
\PRTUE{9.60}{10.45}{C12}\PRTUE{12.35}{6.40}{C13}
\PRTUE{0.45}{2.55}{C14}\PRTUE{0.45}{6.25}{C15}\PRTUE{0.45}{9.95}{C16}
\end{tikzpicture}
\vspace{0.3mm}
\tiny C3 sala/quartos (17) · C4 hall/garagem (9) · C5 cozinha (6) · C6 serviço (4) · C7 gourmet (4) · C8 banheiros (3) · C9--C11 chuveiros · C12 forno · C13 lavadora · C14--C16 ar-condicionado.
\end{frame}

\begin{frame}[shrink=7]{Residencial I — memória dos mínimos por ambiente}
\scriptsize
\begin{tabularx}{\textwidth}{>{\bfseries}p{2.55cm}c c c c >{\centering\arraybackslash}X}
\toprule
Ambiente & Área (m$^2$) & Perím. (m) & Luz (VA) & TUG & TUG (VA)\\
\midrule
Quarto 3 & 15,3 & 17,6 & 220 & 4 & 400\\
Suíte 2 & 14,6 & 17,6 & 220 & 4 & 400\\
Suíte master & 15,3 & 18,0 & 220 & 4 & 400\\
Banheiros (3) & 3,4/4,1/4,4 & critério próprio & 300 & 3 & 1.800\\
Hall íntimo & 13,2 & 24,4 & 160 & 5 & 500\\
Sala/jantar & 30,7 & 22,4 & 460 & 5 & 500\\
Cozinha & 22,1 & 18,8 & 340 & 6 & 2.100\\
Garagem & 21,3 & 18,8 & 280 & 4 & 400\\
Serviço & 9,9 & 12,8 & 100 & 4 & 1.900\\
Varanda gourmet$^*$ & 10,6 & 13,2 & 160 & 4 & 1.900\\
\midrule
\textbf{Total} & -- & -- & \textbf{2.460} & \textbf{43} & \textbf{10.300}\\
\bottomrule
\end{tabularx}
\begin{caixaObs}
Áreas e perímetros seguem a geometria interna desenhada; os dormitórios têm forma em L. $^*$A varanda gourmet foi tratada como ambiente análogo a cozinha por possuir bancada de preparo.
\end{caixaObs}
\end{frame}

\begin{frame}[shrink=6]{Residencial I — circuitos C1 a C8}
\scriptsize
\begin{tabularx}{\textwidth}{c p{3.25cm} c c c c c X}
\toprule
C & Uso & $S$ & $I_b$ & Cu & DJ & Fase & Solução adotada\\
\midrule
C1 & Luz social/cozinha/serviço/gourmet & 1,06 kVA & 4,8 A & 1,5 mm$^2$ & 10 A & A & RCBO 30 mA\\
C2 & Luz íntima/banhos/hall/garagem & 1,40 kVA & 6,4 A & 1,5 mm$^2$ & 10 A & B & RCBO 30 mA\\
C3 & TUG sala e dormitórios & 1,70 kVA & 7,7 A & 2,5 mm$^2$ & 16 A & A & RCBO 30 mA\\
C4 & TUG hall e garagem & 0,90 kVA & 4,1 A & 2,5 mm$^2$ & 16 A & B & RCBO 30 mA\\
C5 & TUG cozinha & 2,10 kVA & 9,5 A & 2,5 mm$^2$ & 16 A & C & RCBO 30 mA\\
C6 & TUG serviço & 1,90 kVA & 8,6 A & 2,5 mm$^2$ & 16 A & C & RCBO 30 mA\\
C7 & TUG gourmet & 1,90 kVA & 8,6 A & 2,5 mm$^2$ & 16 A & A & RCBO 30 mA\\
C8 & TUG banheiros & 1,80 kVA & 8,2 A & 2,5 mm$^2$ & 16 A & C & RCBO 30 mA\\
\bottomrule
\end{tabularx}
\begin{caixaObs}
RCBO individuais evitam neutros compartilhados a jusante e limitam a área afetada. Tipo A é a premissa inicial; confirmar forma de onda residual, corrente de fuga e fabricante.
\end{caixaObs}
\end{frame}

\begin{frame}[shrink=6]{Residencial I — circuitos C9 a C16}
\scriptsize
\begin{tabularx}{\textwidth}{c p{3.15cm} c c c c c X}
\toprule
C & Equipamento & $P$ & $I_b$ & Cu & DJ & Fase & Proteção/observação\\
\midrule
C9 & Chuveiro banho social & 4,50 kW & 20,5 A & 4 mm$^2$ & 25 A & A & dedicado; RCBO 30 mA\\
C10 & Chuveiro suíte 2 & 4,50 kW & 20,5 A & 4 mm$^2$ & 25 A & B & dedicado; RCBO 30 mA\\
C11 & Chuveiro suíte master & 4,50 kW & 20,5 A & 4 mm$^2$ & 25 A & C & dedicado; RCBO 30 mA\\
C12 & Forno elétrico & 3,50 kW & 15,9 A & 4 mm$^2$ & 20 A & B & dedicado; RCBO 30 mA\\
C13 & Lavadora & 1,20 kW & 5,5 A & 2,5 mm$^2$ & 16 A & A & dedicado; RCBO 30 mA\\
C14--C16 & Ar-condicionado quartos & 3$\times$1,00 kW & 4,5 A cada & 2,5 mm$^2$ & 16 A & C/B/A & dedicados; fabricante\\
\bottomrule
\end{tabularx}
\begin{caixaAt}
As seções e proteções são seleções iniciais. Antes do executivo: método de instalação, agrupamento, temperatura, $I_z$, queda, curto máximo/mínimo, capacidade de interrupção e manual de cada equipamento.
\end{caixaAt}
\end{frame}

\begin{frame}[shrink=4]{Residencial I — entrada e balanceamento fechados}
\small
\begin{columns}[T,onlytextwidth]
\column{0.49\textwidth}
\begin{caixaProjeto}[title=Balanceamento por valores instalados]
\begin{align*}
S_A&=\SI{11.36}{kVA}\quad (\SI{51.6}{A})\\
S_B&=\SI{11.30}{kVA}\quad (\SI{51.4}{A})\\
S_C&=\SI{11.30}{kVA}\quad (\SI{51.4}{A})
\end{align*}
\scriptsize A: C1+C3+C7+C9+C13+C16\\
B: C2+C4+C10+C12+C15\\
C: C5+C6+C8+C11+C14
\end{caixaProjeto}
\column{0.49\textwidth}
\begin{caixaNorma}[title=Equatorial Maranhão]
Hipótese declarada do estudo: \SI{33.96}{kW} equivalentes. Na Tabela 1 da NT.00001.EQTL Rev.09:
\begin{itemize}
\item faixa 24,1--38 kW;
\item atendimento 380/220 V, 3F+N;
\item disjuntor de entrada 63 A;
\item condutor Cu do cliente no padrão: mínimo 10 mm$^2$.
\end{itemize}
\end{caixaNorma}
\end{columns}
\begin{caixaAt}
Para ilustrar o enquadramento, os VA mínimos de luz/TUG foram tomados numericamente como W. Na ligação real, declarar as potências ativas nominais efetivamente instaladas conforme a definição da NT.00001; demanda não substitui carga instalada nessa tabela. O alimentador interno é recalculado pela NBR 5410.
\end{caixaAt}
\end{frame}

\begin{frame}{Residencial I — unifilar coerente com 16 circuitos}
\centering
\begin{tikzpicture}[scale=0.76,every node/.style={font=\tiny}]
\tikzset{b/.style={draw=corPrim,fill=corDef!7,rounded corners,minimum width=2.05cm,minimum height=0.58cm,align=center}}
\node[b](r)at(0,4.6){Rede EQTL\\380/220 V};
\node[b](m)at(0,3.70){Medição\\padrão de entrada};
\node[b](dg)at(0,2.80){DG 3P 63 A\\+ DPS coordenado};
\node[b](qd)at(0,1.85){QD1\\barras N e PE separadas};
\node[b](g1)at(-4.2,0.45){C1--C8\\luz/TUG + RCBO};
\node[b](g2)at(-1.4,0.45){C9--C11\\chuveiros + RCBO};
\node[b](g3)at(1.4,0.45){C12--C13\\forno/lavadora};
\node[b](g4)at(4.2,0.45){C14--C16\\ar-condicionado};
\node[b,draw=corNorma,fill=corNorma!7](bep)at(5.0,2.65){BEP\\PE + equipot.};
\draw[-{Stealth},thick,corPrim](r)--(m)--(dg)--(qd);
\foreach \x in {g1,g2,g3,g4}{\draw[-{Stealth},thick](qd)--(\x);}
\draw[corNorma,thick](qd)--(bep);
\draw[corNorma,thick](bep.south)--(5.0,1.35)node[ground]{};
\node[font=\tiny,corNorma,anchor=west] at (5.15,1.35){eletrodo(s)};
\end{tikzpicture}
\vspace{1mm}
\scriptsize O PE não é um ``terra isolado'' por circuito; todos os circuitos retornam ao barramento PE, ligado ao BEP e ao sistema comum de aterramento/equipotencialização.
\end{frame}

% Residencial II - apartamento com varanda, duas suites e acessos coerentes
\section{Projeto Residencial II — apartamento com varanda}

\newcommand{\AptoFinalBase}{%
\draw[arqExt] (0,0) rectangle (12,9);\draw[arqExt] (12,4) rectangle (14,9);
\ArqParedeInt{4.2}{0}{4.2}{9}\ArqParedeInt{5.2}{2}{5.2}{6.2}\ArqParedeInt{0}{4.4}{4.2}{4.4}
\ArqParedeInt{2.6}{0}{2.6}{2.1}\ArqParedeInt{2.6}{2.1}{4.2}{2.1}
\ArqParedeInt{2.6}{6.8}{4.2}{6.8}\ArqParedeInt{2.6}{6.8}{2.6}{9}
\ArqParedeInt{5.2}{3.8}{12}{3.8}\ArqParedeInt{7.2}{0}{7.2}{3.8}\ArqParedeInt{10}{0}{10}{3.8}
\ArqParedeInt{5.2}{2.2}{6.6}{2.2}\ArqParedeInt{6.6}{2.2}{6.6}{3.8}
% entrada e acessos
\ArqPortaHUp{5.7}{0}{1.0}\node[arqSub] at (6.2,0.45){ENTRADA / HALL COMUM};
\ArqVaoH{5.65}{7.05}{3.8} % foyer -> sala
\ArqVaoH{7.65}{9.45}{3.8} % cozinha -> sala
\ArqPortaVLeft{5.2}{2.45}{0.85} % foyer -> hall intimo
\ArqPortaVLeft{4.2}{2.55}{0.85}\ArqPortaVLeft{4.2}{4.90}{0.85} % suites
% banhos compactos: portas de correr evitam conflito com loucas e box
\ArqPortaCorrerH{2.95}{3.75}{2.1}\ArqPortaCorrerH{2.95}{3.75}{6.8}
\ArqPortaHUp{5.45}{2.2}{0.75} % lavabo abre para o foyer
\ArqPortaVRight{10}{1.10}{0.80} % cozinha -> servico
\ArqPortaCorrerV{5.0}{8.2}{12} % sala -> varanda
% janelas
\ArqJanV{0.55}{3.4}{0}\ArqJanV{5.25}{8.35}{0}\ArqJanH{7.65}{9.4}{0}\ArqJanH{10.35}{11.6}{0}\ArqJanH{0.6}{2.3}{9}
% suite 2
\ArqCama{0.45}{0.50}{1.95}{2.55}\ArqArmario{0.4}{3.45}{2.8}{0.5}\ArqChuveiro{2.78}{0.15}{1.15}{0.78}\ArqVaso{3.55}{1.35}\ArqPia{2.72}{1.15}{0.55}{0.4}\node[arqLabel] at (1.55,3.9){SUÍTE 2};\node[arqSub] at (3.45,1.92){B2};
% suite master
\ArqCama{0.45}{4.95}{2.05}{2.70}\ArqArmario{0.45}{8.05}{1.85}{0.5}\ArqChuveiro{2.78}{8.05}{1.15}{0.78}\ArqVaso{3.55}{7.35}\ArqPia{2.72}{7}{0.55}{0.4}\node[arqLabel] at (1.5,8.55){SUÍTE MASTER};\node[arqSub] at (3.45,8.72){B1};
\node[arqSub,rotate=90] at (4.70,4.25){HALL ÍNTIMO};
% foyer e lavabo
\node[arqSub] at (6.2,1.45){FOYER};\ArqVaso{5.75}{2.85}\ArqPia{5.35}{3.25}{0.55}{0.38}\node[arqSub] at (5.9,3.6){LAVABO};
% cozinha e servico
\ArqBancada{7.45}{0.25}{2.15}{0.55}\ArqBancada{9.35}{0.8}{0.55}{2.5}\ArqFogao{8.55}{0.18}\ArqGeladeira{7.35}{2.7}\node[arqLabel] at (8.55,3.35){COZINHA};
\ArqLavadora{10.35}{0.35}\ArqBancada{11.1}{0.3}{0.55}{2.65}\node[arqLabel] at (11.05,3.35){SERVIÇO};
% estar e varanda
\ArqSofa{5.8}{4.55}{3.2}{1.05}\ArqBancada{5.55}{7.75}{0.38}{1.0}\ArqMesa{8.65}{6.35}{2.15}{1.0}\node[arqLabel] at (8.65,5.95){ESTAR / JANTAR};
\ArqSofa{12.45}{5.05}{1.1}{2.2}\ArqMesa{12.55}{7.65}{0.85}{0.7}\node[arqLabel,rotate=90] at (13.55,6.5){VARANDA};
\CotaH{0}{12}{9.18}{12,00 m}\CotaV{0}{9}{14.18}{9,00 m}
}

\begin{frame}{Residencial II — acessos refinados}
\centering\begin{tikzpicture}[scale=0.53,every node/.style={font=\tiny}]\AptoFinalBase\end{tikzpicture}
\vspace{0.4mm}\scriptsize Entrada pelo hall comum; lavabo voltado ao foyer; cozinha integrada à sala; serviço acessado pela cozinha; B1/B2 por dentro das suítes, com portas de correr; varanda pela sala.
\end{frame}

\begin{frame}{Residencial II — iluminação A1/A2}
\centering\begin{tikzpicture}[scale=0.53,every node/.style={font=\tiny}]
\AptoFinalBase\SimQD{6.85}{3.30}{QD-A}
\foreach \x/\y in {1.45/2.8,3.35/1,1.5/6,3.35/8,4.7/4.3,6.2/1.2,5.9/2.9,8.5/1.9,11/1.9,8/5.1,10.1/6.5,13/6.6}{\SimLuz{\x}{\y}}
\foreach \x/\y in {3.92/2.82,3.12/2.32,3.92/5.12,3.12/6.58,5.42/0.42,5.42/2.55,6.82/3.52,9.82/1.28,5.42/4.18,11.72/5.18}{\SimInterruptor{\x}{\y}{S}}
\draw[corAccent!80!black,dashed,thick] (6.85,3.30)--(8.0,5.1)--(10.1,6.5)--(13,6.6);
\draw[corAccent!80!black,dashed,thick] (6.85,3.30)--(8.5,1.9)--(11,1.9);
\draw[corAccent!55!black,dashed,thick] (6.85,3.30)--(6.2,1.2)--(5.9,2.9)--(4.7,4.3)--(1.5,6)--(3.35,8);
\draw[corAccent!55!black,dashed,thick] (4.7,4.3)--(1.45,2.8)--(3.35,1);
\node[font=\tiny\bfseries,fill=white,inner sep=1pt,corAccent!80!black] at (9.5,4.8){A1};
\node[font=\tiny\bfseries,fill=white,inner sep=1pt,corAccent!55!black] at (3.2,4.3){A2};
\end{tikzpicture}
\vspace{0.3mm}\scriptsize Previsão: A1 = 0,94 kVA e A2 = 0,94 kVA. Os comandos estão junto aos acessos; a potência mínima total é 1,88 kVA, calculada pelas áreas e pelos acréscimos completos de 4 m$^2$.
\end{frame}

\begin{frame}{Residencial II — 27 TUG e sete cargas específicas}
\centering\begin{tikzpicture}[scale=0.53,every node/.style={font=\tiny}]
\AptoFinalBase\SimQD{6.85}{3.30}{QD-A}
% A3: suites, hall, foyer, sala e varanda — 16 pontos
\foreach \x/\y in {0.25/0.55,0.25/4.15,2.35/0.25,3.95/4.15,0.25/4.85,0.25/8.75,2.35/8.75,3.95/4.65,4.35/4.30,5.35/0.55,5.45/4.05,7.40/4.05,9.50/4.05,11.70/4.10,11.70/8.70,13.70/8.70}{\SimTUG{\x}{\y}}
% A4 cozinha, A5 servico, A6 banheiros
\foreach \x/\y in {7.45/0.25,8.55/0.25,9.75/1.50,7.35/3.55}{\SimTUGM{\x}{\y}}
\foreach \x/\y in {10.15/0.25,11.75/0.25,10.15/3.55,11.75/3.55}{\SimTUGM{\x}{\y}}
\foreach \x/\y in {3.95/1.65,3.95/7.15,5.45/3.50}{\SimTUG{\x}{\y}}
% A7--A13
\SimTUE{3.4}{0.4}{A7}\SimTUE{3.4}{8.55}{A8}\SimTUE{8.9}{0.35}{A9}\SimTUE{10.7}{0.45}{A10}
\SimTUE{0.3}{3.1}{A11}\SimTUE{0.3}{7.95}{A12}\SimTUE{11.55}{7.8}{A13}
\end{tikzpicture}
\vspace{0.3mm}\tiny A3 áreas secas/varanda (16) · A4 cozinha (4) · A5 serviço (4) · A6 banheiros/lavabo (3) · A7--A8 chuveiros · A9 forno · A10 lavadora · A11--A13 ar-condicionado.
\end{frame}

\begin{frame}[shrink=8]{Residencial II — memória dos mínimos por ambiente}
\scriptsize
\begin{tabularx}{\textwidth}{>{\bfseries}p{2.35cm}c c c c >{\centering\arraybackslash}X}
\toprule
Ambiente & Área (m$^2$) & Perím. (m) & Luz (VA) & TUG & TUG (VA)\\
\midrule
Suíte 2 / B2 & 15,1 / 3,4 & 17,2 / próprio & 220 / 100 & 4 / 1 & 400 / 600\\
Suíte master / B1 & 15,8 / 3,5 & 17,6 / próprio & 220 / 100 & 4 / 1 & 400 / 600\\
Hall íntimo & 4,2 & critério por área & 100 & 1 & 100\\
Foyer & 5,4 & critério por área & 100 & 1 & 100\\
Lavabo & 2,2 & critério próprio & 100 & 1 & 600\\
Sala/jantar & 35,4 & 24,0 & 520 & 5 & 500\\
Cozinha & 10,6 & 13,2 & 160 & 4 & 1.900\\
Serviço & 7,6 & 11,6 & 100 & 4 & 1.900\\
Varanda & 10,0 & ao menos uma & 160 & 1 & 100\\
\midrule
\textbf{Total} & -- & -- & \textbf{1.880} & \textbf{27} & \textbf{7.200}\\
\bottomrule
\end{tabularx}
\end{frame}

\begin{frame}[shrink=6]{Residencial II — circuitos A1 a A6}
\scriptsize
\begin{tabularx}{\textwidth}{c p{3.5cm} c c c c c X}
\toprule
C & Uso & $S$ & $I_b$ & Cu & DJ & Fase & Solução adotada\\
\midrule
A1 & Luz social/cozinha/serviço/varanda & 0,94 kVA & 4,3 A & 1,5 mm$^2$ & 10 A & A & RCBO 30 mA\\
A2 & Luz íntima/foyer/lavabo & 0,94 kVA & 4,3 A & 1,5 mm$^2$ & 10 A & B & RCBO 30 mA\\
A3 & TUG áreas secas e varanda & 1,60 kVA & 7,3 A & 2,5 mm$^2$ & 16 A & C & RCBO 30 mA\\
A4 & TUG cozinha & 1,90 kVA & 8,6 A & 2,5 mm$^2$ & 16 A & B & RCBO 30 mA\\
A5 & TUG serviço & 1,90 kVA & 8,6 A & 2,5 mm$^2$ & 16 A & C & RCBO 30 mA\\
A6 & TUG banheiros/lavabo & 1,80 kVA & 8,2 A & 2,5 mm$^2$ & 16 A & C & RCBO 30 mA\\
\bottomrule
\end{tabularx}
\end{frame}

\begin{frame}[shrink=6]{Residencial II — circuitos A7 a A13}
\scriptsize
\begin{tabularx}{\textwidth}{c p{3.1cm} c c c c c X}
\toprule
C & Equipamento & $P$ & $I_b$ & Cu & DJ & Fase & Proteção/observação\\
\midrule
A7 & Chuveiro B2 & 4,50 kW & 20,5 A & 4 mm$^2$ & 25 A & A & dedicado; RCBO 30 mA\\
A8 & Chuveiro B1 & 4,50 kW & 20,5 A & 4 mm$^2$ & 25 A & B & dedicado; RCBO 30 mA\\
A9 & Forno elétrico & 3,50 kW & 15,9 A & 4 mm$^2$ & 20 A & C & dedicado; RCBO 30 mA\\
A10 & Lavadora & 1,20 kW & 5,5 A & 2,5 mm$^2$ & 16 A & B & dedicado; RCBO 30 mA\\
A11--A13 & Ar-condicionado & 3$\times$1,00 kW & 4,5 A cada & 2,5 mm$^2$ & 16 A & A/A/A & dedicados; fabricante\\
\bottomrule
\end{tabularx}
\begin{caixaAt}
Seleções iniciais: recalcular $I_z$, agrupamento, queda, curto, capacidade de interrupção e requisitos do fabricante. A proteção diferencial adotada deve ser compatível com cargas eletrônicas/inverter.
\end{caixaAt}
\end{frame}

\begin{frame}{Residencial II — balanceamento e interface EMUC}
\small\begin{columns}[T,onlytextwidth]
\column{0.49\textwidth}
\begin{caixaProjeto}[title=Unidade revisada]
\begin{align*}
S_A&=\SI{8.44}{kVA}\;(\SI{38.4}{A})\\
S_B&=\SI{8.54}{kVA}\;(\SI{38.8}{A})\\
S_C&=\SI{8.80}{kVA}\;(\SI{40.0}{A})
\end{align*}
Total equivalente: \SI{25.78}{kW}. Proteção de 63 A é a premissa inicial da unidade. Para esta ilustração, os VA de previsão foram tomados numericamente como W; a declaração real usa potências ativas nominais. Verificar também cenários de demanda, pois o balanceamento instalado concentra os três aparelhos de ar-condicionado na fase A.
\end{caixaProjeto}
\column{0.49\textwidth}
\begin{caixaNorma}[title=Edifício multifamiliar]
Centro de medição, demanda do conjunto, prumadas, queda, proteção geral e solução de conexão são definidos pela \textbf{NT.00004.EQTL Rev.08} e seu Anexo I vigente. Não multiplicar a carga desta unidade pelo número de apartamentos sem aplicar o método do empreendimento.
\end{caixaNorma}
\end{columns}
\end{frame}

\begin{frame}{Residencial II — unifilar da unidade}
\centering\begin{tikzpicture}[scale=0.74,every node/.style={font=\tiny}]
\tikzset{b/.style={draw=corPrim,fill=corDef!8,rounded corners,minimum width=2cm,minimum height=0.56cm,align=center}}
\node[b](cm)at(0,4.4){Centro de medição\\EMUC};\node[b](pr)at(0,3.5){Prumada 3F+N+PE};\node[b](qd)at(0,2.55){QD-A\\DG 3P 63 A + DPS};
\node[b](l)at(-3.6,0.9){A1--A6\\luz/TUG + RCBO};\node[b](ch)at(-1.2,0.9){A7--A8\\chuveiros};\node[b](co)at(1.2,0.9){A9--A10\\forno/lavadora};\node[b](ac)at(3.6,0.9){A11--A13\\climatização};
\node[b,draw=corNorma,fill=corNorma!7](bep)at(4.8,2.9){BEP do edifício\\PE/equipot.};
\draw[-{Stealth},thick,corPrim](cm)--(pr)--(qd);\foreach \x in {l,ch,co,ac}{\draw[-{Stealth},thick](qd)--(\x);}\draw[corNorma,thick](pr)--(bep);\draw[corNorma,thick](qd)--(bep);\draw[corNorma,thick](bep.south)--(4.8,1.55)node[ground]{};
\end{tikzpicture}\end{frame}

%=====================================================================
% PROJETO COMERCIAL — LOJA DE SHOPPING COM FOCO EM LUMINOTECNICA
% Caso didatico completo: leitura, calculo, documentacao e execucao.
%=====================================================================

\section{Projeto Comercial — loja de shopping e luminotécnica}

\newcommand{\LojaEnvelope}{%
  \draw[arqExt] (0,0) rectangle (18,10);
  \ArqParedeInt{0}{7.3}{18}{7.3}
  \ArqParedeInt{7.0}{7.3}{7.0}{10}
  \ArqParedeInt{10.0}{7.3}{10.0}{10}
  \ArqParedeInt{15.0}{7.3}{15.0}{10}
  \ArqParedeInt{11.7}{7.3}{11.7}{10}
  \ArqParedeInt{13.35}{7.3}{13.35}{10}
  \ArqJanH{0.25}{7.15}{0}
  \ArqVaoH{7.15}{10.85}{0}
  \ArqJanH{10.85}{17.75}{0}
  \ArqPortaHDown{2.7}{7.3}{0.90}
  \ArqPortaHDown{8.0}{7.3}{0.90}
  \ArqPortaHDown{10.35}{7.3}{0.75}
  \ArqPortaHDown{12.05}{7.3}{0.75}
  \ArqPortaHDown{13.7}{7.3}{0.75}
  \ArqPortaHDown{16.0}{7.3}{0.80}
  \node[arqLabel] at (3.5,9.55){ESTOQUE};
  \node[arqLabel] at (8.5,9.55){APOIO};
  \node[arqSub] at (10.85,9.55){PROV. 1};
  \node[arqSub] at (12.52,9.55){PROV. 2};
  \node[arqSub] at (14.18,9.55){PROV. PCD};
  \node[arqLabel] at (16.5,9.55){ELÉTRICA / TI};
  \node[arqSub] at (9,0.45){ACESSO PELO MALL};
  \CotaH{0}{18}{10.18}{18,00 m}
  \CotaV{0}{10}{18.18}{10,00 m}}

\newcommand{\LojaBase}{%
  \LojaEnvelope
  \draw[arqMob,rounded corners=2pt] (0.45,0.45) rectangle (3.1,1.15);
  \draw[arqMob,rounded corners=2pt] (14.9,0.45) rectangle (17.55,1.15);
  \node[arqSub] at (1.78,0.80){VITRINE A};
  \node[arqSub] at (16.22,0.80){VITRINE B};
  \foreach \x/\y in {3.0/2.4,7.4/2.4,11.8/2.4,3.0/4.6,7.4/4.6,11.8/4.6}{%
    \draw[arqMob,rounded corners=1pt] (\x,\y) rectangle ++(2.4,0.8);}
  \draw[arqMob,rounded corners=1pt] (14.8,5.35) rectangle (17.35,6.35);
  \node[arqLabel] at (16.08,5.85){CAIXA / POS};
  \foreach \y in {7.75,8.45,9.15}{\draw[arqEquip] (0.35,\y) rectangle (6.55,\y+0.25);}
  \ArqDesk{7.45}{8.25}{2.1}{0.65}
  \foreach \x in {10.3,11.95,13.6}{\draw[arqMob] (\x,8.0) rectangle ++(1.0,0.45);}
  \node[arqLabel] at (8.8,6.8){ÁREA DE VENDAS};}

\begin{frame}{Loja de shopping — o que o aluno precisa entregar}
\small
\begin{columns}[T,onlytextwidth]
\column{0.49\textwidth}
\begin{caixaProjeto}[title=Projeto]
\begin{itemize}
\item briefing, arquitetura e tarefas visuais;
\item critérios luminotécnicos por zona;
\item cálculo, seleção e arranjo das luminárias;
\item teto refletido, comandos, cargas e circuitos;
\item quadro, alimentador, unifilar e detalhes.
\end{itemize}
\end{caixaProjeto}
\column{0.49\textwidth}
\begin{caixaNorma}[title=Execução e entrega]
\begin{itemize}
\item compatibilização com o manual do shopping;
\item montagem de trilhos, drivers, quadro e comandos;
\item focalização das luminárias de destaque;
\item ensaios elétricos e medição de iluminância;
\item as built, cenas, regulagens e comissionamento.
\end{itemize}
\end{caixaNorma}
\end{columns}
\end{frame}

\begin{frame}{Briefing e arquitetura da loja didática}
\centering
\begin{tikzpicture}[scale=0.44,every node/.style={font=\tiny}]
\LojaBase
\end{tikzpicture}
\vspace{1mm}
\begin{caixaObs}
Unidade de \SI{180}{m^2}: frente envidraçada, área de vendas, duas vitrines, caixa, estoque, apoio, provadores e sala elétrica/TI. A alimentação nasce no quadro do lojista; o ponto de entrega interno é definido pelo shopping.
\end{caixaObs}
\end{frame}

\begin{frame}{Documentos de entrada: não projetar a loja isoladamente}
\scriptsize
\begin{tabularx}{\textwidth}{>{\bfseries}p{3.0cm}X X}
\toprule
Documento & Informação que deve ser lida & Decisão afetada \\
\midrule
Manual técnico do shopping & tensão, demanda concedida, curto, proteção, medição, horários e materiais & alimentador, QDL, seletividade e obra \\
Projeto arquitetônico & layout, pé-direito, forro, vitrines, provadores e acessos & zonas, focos, comandos e cotas \\
Projeto de incêndio & rota, iluminação/sinalização de emergência e sprinklers & teto refletido e circuitos de segurança \\
Ar-condicionado & fan-coils, condensadoras, drenos e grelhas & carga, quadros e interferências \\
Automação, segurança e TI & POS, rede, CFTV, alarme e controle & tomadas dedicadas, segregação e reserva \\
Catálogos fotométricos & fluxo, distribuição, potência, IRC, CCT, UGR e vida & cálculo e especificação \\
\bottomrule
\end{tabularx}
\begin{caixaAt}
A NT.00001 da distribuidora não dimensiona diretamente o circuito interno da loja. Primeiro vale o ponto de entrega e o manual técnico do empreendimento; a NBR 5410 continua aplicável à instalação de baixa tensão.
\end{caixaAt}
\end{frame}

\begin{frame}{Normas e critérios: declarar a hierarquia}
\small
\begin{columns}[T,onlytextwidth]
\column{0.52\textwidth}
\begin{itemize}
\item ABNT NBR 5410 — instalação elétrica em baixa tensão;
\item ABNT NBR ISO/CIE 8995-1 — iluminação de ambientes de trabalho interiores;
\item ABNT NBR 10898 — sistema de iluminação de emergência;
\item normas do CBMMA, especialmente NT 18 e projeto aprovado;
\item regras de acessibilidade e do código local aplicáveis;
\item manual do shopping e fichas dos fabricantes.
\end{itemize}
\column{0.46\textwidth}
\begin{caixaNorma}[title=Regra de uso]
Confirmar a edição vigente e adquirir os documentos licenciados no início do executivo. Valores deste caso são \textbf{premissas didáticas}, não transcrição de tabelas normativas.
\end{caixaNorma}
\begin{caixaAt}
Conflito entre critérios deve ser resolvido e registrado antes da obra; não escolher silenciosamente o valor mais conveniente.
\end{caixaAt}
\end{columns}
\end{frame}

\begin{frame}{A loja usa camadas de luz, não uma única malha}
\centering
\begin{tikzpicture}[
  b/.style={draw=corPrim,fill=corDef!7,rounded corners,minimum width=2.35cm,minimum height=0.78cm,align=center,font=\scriptsize},
  a/.style={-{Stealth},thick,corPrim}]
\node[b](amb)at(-4.6,1.2){Ambiente\\circulação e leitura geral};
\node[b](ac)at(-1.55,1.2){Destaque\\produto e manequim};
\node[b](vit)at(1.55,1.2){Vitrine\\atração e modelagem};
\node[b](tar)at(4.6,1.2){Tarefa\\caixa, estoque, provador};
\node[b,draw=corNorma,fill=corNorma!7](seg)at(0,-1.0){Emergência\\saída e abandono seguro};
\draw[a](amb)--(ac)--(vit)--(tar);
\draw[a](seg)--(0,0.45);
\end{tikzpicture}
\vspace{4mm}
\begin{caixaObs}
O ambiente fornece base uniforme; o destaque cria hierarquia; a vitrine comunica; a luz de tarefa permite conferir cor, preço e pagamento; a emergência permanece funcional quando a iluminação normal falha.
\end{caixaObs}
\end{frame}

\begin{frame}{Metas do estudo luminotécnico}
\scriptsize
\begin{tabularx}{\textwidth}{>{\bfseries}p{2.45cm}c c c X}
\toprule
Zona & $\overline{E}_m$ & Uniformidade & IRC & Critério complementar \\
\midrule
Área de vendas & 300 lx horizontal & $U_0\geq0{,}40$ & $R_a\geq90$ & contraste por destaque, sem ofuscamento incapacitante \\
Exposição vertical & 500 lx vertical & verificar malha & $R_a\geq90$ & cor e modelagem coerentes com o produto \\
Vitrine & 750 lx de referência & por cena & $R_a\geq90$ & controlar reflexo no vidro e luz do mall \\
Caixa/POS & 500 lx & $U_0\geq0{,}60$ & $R_a\geq80$ & leitura de tela, dinheiro e etiquetas \\
Provadores & 300 lx vertical & simetria facial & $R_a\geq90$ & evitar sombras duras e distorção de cor \\
Estoque & 200 lx & $U_0\geq0{,}40$ & $R_a\geq80$ & leitura de etiquetas e prateleiras \\
\bottomrule
\end{tabularx}
\begin{caixaAt}
Metas são adotadas para o exercício. No projeto real, fechar requisitos com a edição vigente das normas, o cliente, o tipo de mercadoria e o manual do shopping.
\end{caixaAt}
\end{frame}

\begin{frame}{Método dos lúmens — iluminação ambiente da área de vendas}
\small
\[
N=\frac{E_m\,A}{\Phi_L\,UF\,MF}
\]
\begin{caixaEx}[title=Premissas do estudo]
\begin{align*}
A&=18\times7{,}3=131{,}4\ \mathrm{m^2}, & E_m&=300\ \mathrm{lx},\\
\Phi_L&=3\,200\ \mathrm{lm}, & UF&=0{,}72,\quad MF=0{,}80.
\end{align*}
\[
N=\frac{300\times131{,}4}{3\,200\times0{,}72\times0{,}80}=21{,}4
\quad\Longrightarrow\quad \boxed{24\ \text{luminárias}}
\]
\end{caixaEx}
\begin{caixaObs}
Arredondar e distribuir não encerra o projeto: é obrigatório simular/verificar a malha, uniformidade, sombras, iluminância vertical, ofuscamento e interferências no teto.
\end{caixaObs}
\end{frame}

\begin{frame}{UF e MF não podem ser escolhidos por palpite}
\small
\begin{columns}[T,onlytextwidth]
\column{0.48\textwidth}
\begin{caixaNorma}[title=Fator de utilização — UF]
Depende de:
\begin{itemize}
\item distribuição fotométrica da luminária;
\item índice do recinto;
\item refletâncias de teto, paredes e piso;
\item plano de trabalho e geometria real.
\end{itemize}
Obter no arquivo/catálogo fotométrico ou no software validado.
\end{caixaNorma}
\column{0.50\textwidth}
\begin{caixaNorma}[title=Fator de manutenção — MF]
Depende de:
\begin{itemize}
\item depreciação do fluxo do LED;
\item sujeira na luminária e no recinto;
\item falha/substituição e regime de operação;
\item intervalo e método de limpeza.
\end{itemize}
O memorial deve registrar o plano de manutenção associado.
\end{caixaNorma}
\end{columns}
\end{frame}

\begin{frame}{Teto refletido — arranjo e identificação}
\centering
\begin{tikzpicture}[scale=0.44,every node/.style={font=\tiny}]
\LojaEnvelope
\foreach \x in {1.5,4.5,7.5,10.5,13.5,16.5}{
  \foreach \y in {1.35,3.05,4.75,6.45}{
    \draw[corAccent,fill=corAccent!12,rounded corners=1pt] (\x-0.28,\y-0.10) rectangle (\x+0.28,\y+0.10);}}
\foreach \x/\y in {2.2/2.15,4.0/2.15,6.6/2.15,8.4/2.15,11.0/2.15,12.8/2.15,2.2/5.35,4.0/5.35,6.6/5.35,8.4/5.35,11.0/5.35,12.8/5.35,1.4/0.8,3.0/0.8,15.0/0.8,16.6/0.8,15.5/5.6,17.0/5.6}{
  \draw[corFix,fill=corFix!10] (\x,\y) circle (0.13);}
\foreach \x in {1.0,3.0,5.0}{\draw[corAccent,fill=corAccent!12] (\x,8.4) rectangle ++(0.8,0.16);}
\foreach \x in {10.45,12.1,13.75}{\draw[corAccent,fill=corAccent!12] (\x,8.6) circle (0.13);}
\SimQD{16.5}{8.55}{QDL}
\draw[corAt,thick,-{Stealth}] (8.95,0.55)--(8.95,0.05);
\node[corAt,font=\tiny\bfseries] at (9.5,0.35){SAÍDA};
\end{tikzpicture}
\vspace{0.5mm}
\scriptsize Retângulos âmbar: ambiente L-A; círculos magenta: destaque L-D; fundo: estoque/provadores; QDL e saída permanecem acessíveis. A prancha executiva deve acrescentar cotas, códigos de circuito, altura, modelo e orientação.
\end{frame}

\begin{frame}{Luz de destaque — mirar o plano vertical correto}
\small
\begin{columns}[T,onlytextwidth]
\column{0.50\textwidth}
\begin{tikzpicture}[scale=0.82]
\draw[very thick] (-0.2,0)--(-0.2,5.0);
\draw[very thick] (-0.2,5.0)--(5.4,5.0);
\draw[arqMob] (3.35,0) rectangle (4.75,3.0);
\node[font=\scriptsize] at (4.05,1.5){expositor};
\draw[corFix,fill=corFix!12] (2.0,4.85) circle (0.18);
\draw[corFix!75,fill=corFix!8] (2.0,4.65)--(3.45,0.2)--(4.65,0.2)--cycle;
\draw[dashed,corAccent] (2.0,4.65)--(4.05,1.5);
\draw[corPrim,<->] (2.0,4.25) arc (270:305:0.65);
\node[font=\tiny,corPrim] at (2.75,4.05){ângulo de foco};
\end{tikzpicture}
\column{0.48\textwidth}
\begin{itemize}
\item usar fotometria e ângulo de abertura reais;
\item evitar foco atrás da pessoa, que cria sombra sobre o produto;
\item reduzir reflexo especular em vidro, tela e embalagem;
\item limitar luminância direta no campo visual;
\item prever ajuste e travamento após o merchandising;
\item registrar posição e direção no as built.
\end{itemize}
\end{columns}
\end{frame}

\begin{frame}{Qualidade da luz — especificação mínima rastreável}
\scriptsize
\begin{tabularx}{\textwidth}{>{\bfseries}p{2.55cm}X X}
\toprule
Parâmetro & O que significa & Como fechar no projeto \\
\midrule
Fluxo e fotometria & quantidade e distribuição de luz & arquivo IES/LDT da referência ofertada \\
Potência e corrente & carga real do circuito/driver & ficha do conjunto, inclusive perdas \\
IRC e fidelidade & reprodução das cores do produto & $R_a$ e métricas adicionais exigidas pela marca \\
CCT e consistência & aparência de cor e variação entre peças & temperatura escolhida e tolerância/binning \\
Ofuscamento & desconforto por luminância elevada & posição, óptica, acabamento e avaliação do ambiente \\
Flicker e estroboscopia & modulação perceptível ou por câmera & compatibilidade driver, dimerização e filmagem \\
Vida e manutenção & depreciação e falhas do conjunto & LxBy, garantia, acesso e reposição compatível \\
\bottomrule
\end{tabularx}
\begin{caixaAt}
Não aceitar substituição apenas por “mesma potência”. Duas luminárias de 20 W podem produzir distribuições, ofuscamento, cor e durabilidade totalmente diferentes.
\end{caixaAt}
\end{frame}

\begin{frame}{Cenas e comandos — o funcional antes da automação}
\centering
\begin{tikzpicture}[scale=0.84,transform shape,
  b/.style={draw=corPrim,fill=corDef!7,rounded corners,minimum width=2.15cm,minimum height=0.66cm,align=center,font=\scriptsize},
  a/.style={-{Stealth},thick,corPrim}]
\node[b](q)at(0,3.5){QDL — circuitos de força};
\node[b](ctl)at(0,2.25){Controlador / contatores\\relógio + interface};
\node[b](amb)at(-4.5,0.6){C1 ambiente};
\node[b](des)at(-2.25,-0.5){C2 destaque};
\node[b](vit)at(0,0.6){C3 vitrine/fachada};
\node[b](tar)at(2.25,-0.5){C4 caixa/provador};
\node[b](est)at(4.5,0.6){C5 estoque/sensor};
\draw[a](q)--(ctl);
\foreach \x in {amb,des,vit,tar,est}{\draw[a](ctl)--(\x);}
\end{tikzpicture}
\vspace{2mm}
\begin{caixaProjeto}[title=Cenas previstas]
Limpeza (ambiente 100\%), abertura (ambiente + destaque), operação (todas as zonas úteis), vitrine após fechamento e desligamento geral com exceções declaradas. Emergência não depende do controlador de cenas.
\end{caixaProjeto}
\end{frame}

\begin{frame}{Iluminação de emergência — sistema independente e coordenado}
\centering
\begin{tikzpicture}[scale=0.44,every node/.style={font=\tiny}]
\LojaEnvelope
\draw[corAt,line width=2.1pt,-{Stealth}] (3.0,8.5)--(3.0,6.9)--(8.9,6.9)--(8.9,0.15);
\foreach \x/\y in {3.0/8.5,3.0/6.9,8.9/6.9,8.9/3.8,8.9/0.7}{\SimEmerg{\x}{\y}}
\node[draw=corAt,fill=corAt!10,minimum width=1.2cm,minimum height=0.45cm] at (8.9,0.32){SAÍDA};
\node[font=\tiny,corAt] at (5.9,7.15){rota de abandono do estudo};
\end{tikzpicture}
\vspace{1mm}
\begin{caixaAt}
Tipo de sistema, autonomia, níveis, espaçamento, alimentação e sinalização devem seguir o projeto de incêndio aprovado, a NBR 10898 e a NT 18 do CBMMA. O desenho da loja deve coincidir com o desenho de segurança do shopping.
\end{caixaAt}
\end{frame}

\begin{frame}{Potência de iluminação — fechar pelo equipamento especificado}
\scriptsize
\begin{tabularx}{\textwidth}{>{\bfseries}p{2.6cm}r r r X}
\toprule
Sistema & Quant. & Pot. unit. & Potência & Circuito \\
\midrule
Ambiente linear & 24 & 28 W & 672 W & C1 \\
Destaque em trilho & 18 & 18 W & 324 W & C2 \\
Vitrine & 8 & 24 W & 192 W & C3 \\
Caixa/provadores & 10 & 12 W & 120 W & C4 \\
Estoque/apoio & 6 & 18 W & 108 W & C5 \\
Fachada/letreiro & 1 conjunto & 150 W & 150 W & C6 \\
Emergência & 5 & 3 W & 15 W & CE \\
\midrule
\textbf{Total} & -- & -- & \textbf{1 581 W} & -- \\
\bottomrule
\end{tabularx}
\[
\mathrm{DPI}=\frac{1\,581}{180}=8{,}8\ \mathrm{W/m^2}
\]
\begin{caixaObs}
Potência específica permite comparar alternativas, mas não comprova iluminância ou qualidade. Eficiência energética e desempenho visual precisam ser atendidos juntos.
\end{caixaObs}
\end{frame}

\begin{frame}{Quadro de cargas da loja — hipótese didática}
\scriptsize
\begin{tabularx}{\textwidth}{c l c c c c X}
\toprule
Circ. & Uso & Fase & $P/S$ & $I_b$ & DJ & Condutor inicial \\
\midrule
C1 & ambiente & L1-N & 0,67 kW & 3,2 A & 10 A & 1,5 mm$^2$ \\
C2 & destaque & L2-N & 0,32 kW & 1,6 A & 10 A & 1,5 mm$^2$ \\
C3 & vitrine & L3-N & 0,19 kW & 0,9 A & 10 A & 1,5 mm$^2$ \\
C4 & caixa/provadores & L1-N & 0,12 kW & 0,6 A & 10 A & 1,5 mm$^2$ \\
C5 & estoque/apoio & L2-N & 0,11 kW & 0,5 A & 10 A & 1,5 mm$^2$ \\
C6 & letreiro & L3-N & 0,15 kW & 0,7 A & 10 A & 1,5 mm$^2$ \\
C7--C9 & TUG/POS/TI & L1/L2/L3-N & 4,0 kVA & -- & 16/20 A & 2,5 mm$^2$ \\
C10 & climatização & 3F & 12,0 kW & 19,8 A & 25 A 3P & 6 mm$^2$ \\
CE & emergência & circuito próprio & 0,02 kW & -- & projeto & projeto de incêndio \\
\bottomrule
\end{tabularx}
\begin{caixaAt}
Seções são ponto de partida. Recalcular por método de instalação, temperatura, agrupamento, harmônicas, queda, curto, proteção e terminal do equipamento.
\end{caixaAt}
\end{frame}

\begin{frame}{Alimentador do lojista — corrente e critérios simultâneos}
\small
\begin{columns}[T,onlytextwidth]
\column{0.48\textwidth}
Hipóteses: $P_{dem}=17{,}2\ \mathrm{kW}$, $V=380\ \mathrm{V}$, $fp=0{,}92$, trifásico.
\[
I_b=\frac{17\,200}{\sqrt3\,380\,0{,}92}=28{,}4\ \mathrm{A}
\]
\begin{caixaEx}[title=Seleção preliminar]
DJ 3P 32 A e cobre 10 mm$^2$, desde que:
\[
I_b\leq I_n\leq I_z.
\]
\end{caixaEx}
\column{0.50\textwidth}
Para $I_{z,tab}=57$ A, $k_T=0{,}87$ e $k_G=0{,}80$:
\[
I_z=57\times0{,}87\times0{,}80=39{,}7\ \mathrm{A}.
\]
Com $L=65$ m e cobre a quente, estimativa:
\[
\Delta V\%=\frac{100\sqrt3\rho L I\cos\varphi}{SV}=1{,}74\%.
\]
\begin{caixaAt}
Confirmar curto, seletividade e orçamento de queda permitido pelo shopping.
\end{caixaAt}
\end{columns}
\end{frame}

\begin{frame}{Balanceamento: ler circuito, fase e potência juntos}
\small
\begin{columns}[T,onlytextwidth]
\column{0.58\textwidth}
\begin{tabular}{l r r r}
\toprule
Grupo monofásico & L1 & L2 & L3 \\
\midrule
Iluminação & 0,79 & 0,43 & 0,34 \\
TUG de vendas & 0,70 & 0,70 & 0,60 \\
POS/TI/segurança & 1,20 & 0,40 & 0,40 \\
Reserva operacional & 0,50 & 0,50 & 0,50 \\
\midrule
\textbf{Total (kVA)} & \textbf{3,19} & \textbf{2,03} & \textbf{1,84} \\
\bottomrule
\end{tabular}
\column{0.40\textwidth}
\begin{caixaAt}[title=Leitura crítica]
O quadro ainda está desequilibrado. Transferir parte de POS/TI ou TUG de L1 para L2/L3, respeitando continuidade, circuitos dedicados e operação.
\end{caixaAt}
\begin{caixaObs}
Carga trifásica de climatização não corrige o desequilíbrio entre cargas monofásicas.
\end{caixaObs}
\end{columns}
\end{frame}

\begin{frame}{Unifilar do QDL — da alimentação aos usos}
\centering
\begin{tikzpicture}[scale=0.84,transform shape,
  b/.style={draw=corPrim,fill=corDef!7,rounded corners,minimum width=1.85cm,minimum height=0.62cm,align=center,font=\tiny},
  a/.style={-{Stealth},thick,corPrim}]
\node[b](orig)at(0,4.5){QD do shopping\\380/220 V};
\node[b](alim)at(0,3.5){DJ origem + alimentador\\5G10 mm$^2$};
\node[b](dg)at(0,2.5){QDL: DG 3P 32 A\\DPS + seccionamento};
\node[b](dr1)at(-3.9,1.2){DR/RCBO grupo A\\conforme análise};
\node[b](dr2)at(0,1.2){Barramento protegido\\climatização 3F};
\node[b](dr3)at(3.9,1.2){DR/RCBO grupo B\\conforme análise};
\node[b](c1)at(-5.0,-0.3){C1/C4\\L1-N};
\node[b](c2)at(-3.0,-0.3){C2/C5\\L2-N};
\node[b](c10)at(0,-0.3){C10 AC\\3F + PE};
\node[b](c3)at(3.0,-0.3){C3/C6\\L3-N};
\node[b](t)at(5.0,-0.3){C7--C9\\TUG/POS/TI};
\node[b,draw=corNorma,fill=corNorma!7](pe)at(5.0,3.35){Barra PE\\BEP do shopping};
\draw[a](orig)--(alim)--(dg);
\draw[a](dg)--(dr1);\draw[a](dg)--(dr2);\draw[a](dg)--(dr3);
\draw[a](dr1)--(c1);\draw[a](dr1)--(c2);\draw[a](dr2)--(c10);\draw[a](dr3)--(c3);\draw[a](dr3)--(t);
\draw[corNorma,thick](pe)--(dg);
\end{tikzpicture}
\vspace{2mm}
\begin{caixaObs}
O executivo informa polos, curva, capacidade de interrupção, esquema de aterramento, ligação do DPS, barramentos de neutro por DR e terminais. “DPS + DR” sem esses dados não é diagrama executivo.
\end{caixaObs}
\end{frame}

\begin{frame}{Detalhe de montagem — trilho, driver e manutenção}
\small
\begin{columns}[T,onlytextwidth]
\column{0.52\textwidth}
\begin{tikzpicture}[scale=0.80]
\draw[very thick] (-0.4,4.8)--(6.0,4.8);
\draw[fill=black!10,draw=black!50] (0.4,4.45) rectangle (5.2,4.75);
\draw[corPrim,thick] (0.8,4.60)--(4.8,4.60);
\draw[fill=corDef!10,draw=corPrim,rounded corners] (1.0,3.45) rectangle (3.0,4.05);
\node[font=\tiny] at (2.0,3.75){driver acessível};
\draw[corPrim,thick] (2.0,4.45)--(2.0,4.05);
\draw[fill=corFix!12,draw=corFix,thick] (4.0,3.7) circle (0.35);
\draw[corFix,thick] (4.0,4.45)--(4.0,4.05);
\draw[corFix!55,fill=corFix!7] (4.0,3.35)--(3.15,0.4)--(4.85,0.4)--cycle;
\draw[dashed,corAt] (0.2,3.15) rectangle (3.35,4.2);
\node[font=\tiny,corAt] at (1.8,2.9){alçapão / acesso};
\node[font=\tiny] at (4.0,0.1){produto};
\end{tikzpicture}
\column{0.46\textwidth}
\begin{enumerate}
\item confirmar capacidade e continuidade do trilho;
\item fixar na estrutura prevista, não apenas na placa do forro;
\item alimentar com caixa/terminal acessível;
\item garantir ventilação e acesso ao driver;
\item identificar circuito e fase;
\item orientar, focalizar e travar o projetor;
\item registrar modelo, óptica e posição final.
\end{enumerate}
\end{columns}
\end{frame}

\begin{frame}{Compatibilização do teto — checagem antes de furar}
\scriptsize
\begin{tabularx}{\textwidth}{>{\bfseries}p{3.0cm}X X}
\toprule
Interferência & Risco & Solução de projeto \\
\midrule
Sprinkler e detector & obstrução, sombra ou descumprimento do incêndio & respeitar cobertura e afastamentos do sistema aprovado \\
Difusor de ar & oscilação, sujeira, condensação ou desconforto & coordenar eixo e distância com HVAC \\
Alçapão & equipamento inacessível & reservar faixa de manutenção e raio de abertura \\
Estrutura/suporte & fixação sem capacidade & indicar substrato, suporte e carga \\
Comunicação visual & sombra, reflexo ou conflito de altura & verificar corte, elevação e cena noturna \\
Eletrocalha/dados & cruzamentos e segregação inadequada & compatibilizar rota e caixas de derivação \\
\bottomrule
\end{tabularx}
\begin{caixaAt}
Uma planta de teto refletido sem as demais disciplinas é apenas um desenho de intenção. A liberação de obra exige sobreposição e análise de interferências.
\end{caixaAt}
\end{frame}

\begin{frame}{Comissionamento luminotécnico — medir uma malha, não um ponto}
\small
\begin{columns}[T,onlytextwidth]
\column{0.48\textwidth}
\begin{tikzpicture}[scale=0.72]
\draw[very thick] (0,0) rectangle (7.2,4.8);
\foreach \x in {0.9,2.7,4.5,6.3}{\foreach \y in {0.8,2.4,4.0}{\fill[corAccent] (\x,\y) circle (2.4pt);}}
\draw[dashed,corPrim] (0.9,0.8) grid[xstep=1.8,ystep=1.6] (6.3,4.0);
\node[font=\scriptsize] at (3.6,-0.45){malha no plano de medição};
\end{tikzpicture}
\column{0.50\textwidth}
\begin{enumerate}
\item estabilizar a instalação e registrar a cena;
\item controlar ou registrar contribuição da luz externa;
\item usar luxímetro adequado, orientado e sem sombra do operador;
\item medir todos os pontos e a vertical onde aplicável;
\item calcular média, mínimo e uniformidade;
\item ajustar focos, cenas e dimerização;
\item repetir e anexar resultados ao comissionamento.
\end{enumerate}
\[
U_0=\frac{E_{min}}{\overline{E}}.
\]
\end{columns}
\end{frame}

\begin{frame}{Exemplo de aceitação — interpretar os números}
\small
\begin{columns}[T,onlytextwidth]
\column{0.55\textwidth}
\begin{tabular}{c r@{\quad}c r@{\quad}c r}
\toprule
Ponto & lx & Ponto & lx & Ponto & lx \\
\midrule
1 & 322 & 5 & 301 & 9 & 315 \\
2 & 336 & 6 & 289 & 10 & 298 \\
3 & 310 & 7 & 274 & 11 & 307 \\
4 & 295 & 8 & 281 & 12 & 324 \\
\bottomrule
\end{tabular}
\[
\overline{E}=304\ \mathrm{lx},\qquad E_{min}=274\ \mathrm{lx},\qquad U_0=0{,}90.
\]
\column{0.43\textwidth}
\begin{caixaEx}[title=Conclusão do exemplo]
A meta horizontal de 300 lx e a uniformidade adotada foram atendidas pela média/mínimo medidos.
\end{caixaEx}
\begin{caixaAt}
Isso ainda não aceita a loja inteira: conferir vertical de produtos, vitrine, caixa, provadores, emergência, ofuscamento, cor, cenas e potência.
\end{caixaAt}
\end{columns}
\end{frame}

\begin{frame}{Entrega executiva da loja — conjunto mínimo}
\small
\begin{enumerate}
\item planta de pontos e teto refletido, cotados e compatibilizados;
\item legenda, códigos, detalhes de fixação e cortes relevantes;
\item memória luminotécnica, arquivos fotométricos e critérios;
\item quadro de cargas, balanceamento, alimentador e queda;
\item unifilar, diagrama de comandos/cenas e lista de cabos;
\item especificação de luminárias, drivers, trilhos e controles;
\item lista de materiais e plano de montagem por etapa;
\item relatórios de ensaios elétricos e medições luminotécnicas;
\item as built, regulagens finais, cenas e plano de manutenção.
\end{enumerate}
\begin{caixaProjeto}[title=Competência avaliada]
O aluno deve conseguir localizar um ponto na planta, achar seu circuito no quadro, interpretar sua ligação no diagrama, executar a montagem, medir o resultado e justificar qualquer divergência.
\end{caixaProjeto}
\end{frame}

\begin{frame}{Exercício de leitura integrada — loja de shopping}
\small
Sem alterar as premissas, produza uma revisão R01 que:
\begin{enumerate}
\item transfira cargas monofásicas para reduzir o desequilíbrio;
\item identifique na planta todas as luminárias de C2 e sua orientação;
\item acrescente as referências cruzadas entre teto, unifilar e detalhe de trilho;
\item indique quais circuitos permanecem após o fechamento da loja;
\item estabeleça uma malha de medição para vitrine e provadores;
\item atualize quadro, lista de cabos, lista de materiais e memorial;
\item gere uma ordem de serviço de montagem e uma ficha de comissionamento.
\end{enumerate}
\begin{caixaAt}
Critério de correção: a revisão deve permanecer coerente em todos os documentos. Alterar apenas a planta não vale como projeto revisado.
\end{caixaAt}
\end{frame}


\section{Exemplos resolvidos}

\begin{frame}{Exemplo 1 — circuito de TUG em área seca}
\small
\begin{columns}[T,onlytextwidth]
\column{0.49\textwidth}
\begin{caixaEx}
Circuito em 220 V, $P=\SI{1100}{VA}$, comprimento de estudo $L=\SI{30}{m}$:
\[
I_b=\frac{1100}{220}=\SI{5.00}{A}.
\]
Adota-se inicialmente Cu 2,5 mm$^2$ e disjuntor 16 A.
\end{caixaEx}
\column{0.49\textwidth}
\textbf{A seleção só fecha após verificar:}
\begin{enumerate}
\item método de instalação;
\item temperatura e agrupamento;
\item $I_b\le I_n\le I_z$;
\item queda de tensão;
\item curto e desligamento;
\item necessidade de proteção diferencial.
\end{enumerate}
\end{columns}
\end{frame}

\begin{frame}{Exemplo 2 — circuito de chuveiro}
\small
\begin{columns}[T,onlytextwidth]
\column{0.49\textwidth}
\begin{caixaEx}
$P=\SI{5.5}{kW}$, $V=\SI{220}{V}$:
\[
I_b=\frac{5500}{220}=\SI{25}{A}.
\]
Para este exemplo de 5,5 kW, adota-se inicialmente circuito dedicado, Cu 6 mm$^2$ e DJ 32 A, após as verificações de capacidade de condução e queda. O projeto residencial adiante usa chuveiros de 4,5 kW e outra seleção.
\end{caixaEx}
\column{0.49\textwidth}
\begin{caixaNorma}[title=Proteção adicional]
O circuito que serve ponto de utilização em local contendo chuveiro deve receber proteção diferencial-residual de alta sensibilidade conforme os requisitos aplicáveis.
\end{caixaNorma}
\begin{caixaAt}
Não escolher seção apenas por uma tabela de potência de chuveiro. Comprimento, método, temperatura e agrupamento mudam o resultado.
\end{caixaAt}
\end{columns}
\end{frame}

\begin{frame}{Exemplo 3 — balanceamento antes do padrão de entrada}
\small
\begin{columns}[T,onlytextwidth]
\column{0.52\textwidth}
No projeto residencial desenvolvido neste material:
\[
S_A=\SI{11.36}{kVA},\qquad S_B=S_C=\SI{11.30}{kVA}.
\]
Em 220 V fase-neutro:
\[
I_A\approx\SI{51.6}{A},\qquad I_B=I_C\approx\SI{51.4}{A}.
\]
\column{0.46\textwidth}
\begin{caixaProjeto}[title=Distribuição]
\begin{itemize}
\item A: C1+C3+C7+C9+C13+C16;
\item B: C2+C4+C10+C12+C15;
\item C: C5+C6+C8+C11+C14.
\end{itemize}
A soma fecha com os 16 circuitos e mantém um chuveiro em cada fase. Este é o balanceamento dos valores instalados; cenários de demanda também devem ser verificados.
\end{caixaProjeto}
\end{columns}
\end{frame}

\begin{frame}{Exemplo 4 — enquadramento Equatorial Maranhão}
\small
\begin{columns}[T,onlytextwidth]
\column{0.50\textwidth}
Carga instalada do projeto:
\[
P_{inst,eq}=\SI{33.96}{kW}.
\]
A NT.00001.EQTL Rev.09 enquadra a faixa de 24,1 a 38 kW, na Tabela 1 de 220/380 V, em ligação trifásica com disjuntor termomagnético de entrada de 63 A.
\column{0.48\textwidth}
\begin{caixaNorma}[title=No Maranhão]
A ligação trifásica urbana é 380/220 V, 3F+N. A mesma tabela indica 10 mm$^2$ de cobre como mínimo do condutor do cliente no padrão correspondente.
\end{caixaNorma}
\begin{caixaAt}
O alimentador interno é recalculado pela NBR 5410; não se copia automaticamente o mínimo do padrão da concessionária.
\end{caixaAt}
\end{columns}
\end{frame}

\begin{frame}{Exemplo 5 — recarga de VE e aumento de carga}
\small
\begin{columns}[T,onlytextwidth]
\column{0.49\textwidth}
Se ao projeto de $\SI{33.96}{kW}$ for acrescentado um SAVE de $\SI{7.4}{kW}$:
\[
P_{inst,novo}=33{,}96+7{,}40=\SI{41.36}{kW}.
\]
O exemplo passa à faixa de 38,1 a 48 kW, associada a 80 A na Tabela 1; o enquadramento da entrada precisa ser revisto.
\column{0.49\textwidth}
\begin{caixaProjeto}[title=Consequência]
Além do novo circuito e da capacidade do QD/alimentador, a conexão da estação de recarga deve observar a NT.00042.EQTL Rev.04 e eventual necessidade de aumento de carga perante a distribuidora.
\end{caixaProjeto}
\end{columns}
\end{frame}

%=====================================================================

\section{Exercícios N1--N4}

\begin{frame}{N1 — leitura da planta e requisitos mínimos}
\small
\begin{enumerate}
\item Para uma sala de 18 m$^2$ e perímetro de 17 m, determine a previsão mínima de carga de iluminação e o número mínimo de TUG.
\item Explique por que a potência de previsão de iluminação não determina diretamente a quantidade de luminárias.
\item Em uma planta residencial, indique os locais adequados para QD, interruptores e tomadas considerando portas, circulação e mobiliário.
\item Diferencie TUG, TUE, circuito dedicado, PE, DR/DDR e DPS.
\end{enumerate}
\end{frame}

\begin{frame}{N2 — dimensionamento de circuitos}
\small
\begin{enumerate}
\item Dimensione um circuito de TUG em 220 V, 1.900 VA e 24 m, verificando corrente, proteção, seção inicial e queda de tensão.
\item Analise um chuveiro de 5,5 kW em 220 V a 28 m do QD; verifique a seção por ampacidade e queda de tensão.
\item Proponha a divisão de circuitos para uma residência com sala, cozinha, área de serviço, dois quartos, banheiro e escritório.
\item Explique quais desses circuitos exigem proteção diferencial-residual de alta sensibilidade e por quê.
\end{enumerate}
\end{frame}

\begin{frame}{N3 — quadro, balanceamento e entrada Equatorial}
\small
\begin{enumerate}
\item Distribua 11 circuitos monofásicos de potências distintas entre A, B e C em uma rede 380/220 V, minimizando o desbalanceamento da carga instalada.
\item Para uma unidade residencial no Maranhão com 21,5 kW instalados, determine o enquadramento de fornecimento usando a NT.00001.EQTL Rev.09 e identifique quais dados ainda precisam ser verificados no padrão de entrada.
\item Mostre por que o condutor mínimo indicado pela concessionária para o padrão não pode ser copiado automaticamente para o alimentador interno.
\item Desenhe um unifilar com DG, DPS, barramentos N/PE, circuitos com DDR/RCBO e identificação coerente com o quadro de cargas.
\end{enumerate}
\end{frame}

\begin{frame}{N4 — desafio de projeto residencial completo}
\small
Receba uma planta arquitetônica residencial cotada e, sem alterar sua geometria, desenvolva o projeto elétrico completo:
\begin{enumerate}
\item calcule os pontos mínimos de iluminação e TUG e justifique pontos adicionais;
\item posicione TUE e cargas específicas conforme uso real;
\item divida e identifique os circuitos na planta;
\item dimensione condutores, eletrodutos e proteções;
\item verifique queda de tensão, proteção diferencial, DPS, PE e equipotencialização;
\item balanceie as fases e enquadre a entrada conforme Equatorial Maranhão;
\item gere quadro de cargas, unifilar, memorial de cálculo e lista de materiais;
\item confira se planta, quadro e unifilar apresentam exatamente os mesmos circuitos.
\end{enumerate}
\end{frame}

%=====================================================================


\section{Referências e rastreabilidade}

\begin{frame}[shrink=4]{Fontes oficiais verificadas nesta revisão}
\small
\begin{itemize}
\item \textbf{ABNT Catálogo}: confirmação da ABNT NBR 5419-1:2026 e consulta das edições indicadas.\\
\href{https://www.abntcatalogo.com.br/}{\texttt{abntcatalogo.com.br}}
\item \textbf{Equatorial Energia Maranhão}: lista vigente das normas de conexão e respectivos anexos, incluindo NT.00001 Rev.09, NT.00002 Rev.10, NT.00004 Rev.08, NT.00020 Rev.06, NT.00030 Rev.03 e NT.00042 Rev.04.\\
\href{https://ma.equatorialenergia.com.br/institucional/leis-e-normas-tecnicas/normas-tecnicas/instalacao-de-padrao/}{\texttt{Equatorial MA — Normas Técnicas de Conexão}}
\item \textbf{Ministério do Trabalho e Emprego}: NR-10 vigente até 31/05/2027 e nova redação com vigência em 01/06/2027.\\
\href{https://www.gov.br/trabalho-e-emprego/pt-br/acesso-a-informacao/participacao-social/conselhos-e-orgaos-colegiados/comissao-tripartite-partitaria-permanente/normas-regulamentadora/normas-regulamentadoras-vigentes/norma-regulamentadora-no-10-nr-10}{\texttt{MTE — página oficial da NR-10}}
\end{itemize}
\begin{caixaAt}
Consulta realizada em 13/08/2026. A lista da distribuidora pode mudar; no início de cada projeto, conferir novamente revisão, anexos e período de transição. Este material resume critérios e não reproduz nem substitui o texto integral das normas ABNT.
\end{caixaAt}
\end{frame}

\begin{frame}{Referências para a loja de shopping}
\small
\begin{itemize}
\item \textbf{ABNT NBR ISO/CIE 8995-1}: iluminação de ambientes de trabalho interiores; confirmar edição vigente no \href{https://www.abntcatalogo.com.br/}{Catálogo ABNT}.
\item \textbf{ABNT NBR 10898}: sistema de iluminação de emergência; usar em conjunto com o projeto de segurança contra incêndio aprovado.
\item \textbf{CBMMA}: \href{https://cbm.ssp.ma.gov.br/unidades-bm/unidades-administrativas/dat/legislacao-tecnica-dat/}{legislação técnica oficial}, incluindo NT 18 para iluminação de emergência e demais normas aplicáveis ao shopping.
\item \textbf{Manual técnico do empreendimento}: ponto de entrega interno, demanda concedida, proteção, medição, materiais, procedimentos de obra e comissionamento.
\item \textbf{Dados fotométricos e fichas dos fabricantes}: arquivos IES/LDT, fluxo, distribuição, potência, IRC, CCT, compatibilidade de controles, torque e manutenção.
\end{itemize}
\begin{caixaAt}
As metas numéricas do estudo da loja são premissas didáticas declaradas. O executivo deve confrontá-las com as edições normativas vigentes, a atividade visual real, a marca e as exigências do shopping.
\end{caixaAt}
\end{frame}

\section{Fechamento}
\begin{frame}{Síntese}
\small
\begin{caixaProjeto}[title=Competência de projeto]
O dimensionamento elétrico não é uma sequência isolada de fórmulas. Um projeto consistente coordena \textbf{carga, demanda, condutor, proteção, queda de tensão, aterramento, segurança, documentação e operação} como partes de um único sistema.
\end{caixaProjeto}
\vspace{2mm}
\begin{caixaObs}
Os exemplos são didáticos. Projetos executivos devem utilizar dados reais da instalação, critérios normativos vigentes, informações da concessionária e responsabilidade técnica profissional.
\end{caixaObs}
\end{frame}



\end{document}
